% Polygones réguliers — Mathématiques 3ème — Cours signé OnBuch+
% Programme officiel MINESEC. Compiler avec : tectonic N13-polygones-reguliers.tex
\newcommand{\DOCMATIERE}{Mathématiques}
\newcommand{\DOCNIVEAU}{3ème}
\newcommand{\DOCMODULE}{Mathématiques – classe de 3ème}
\newcommand{\DOCLECON}{N13}
\newcommand{\DOCTITRE}{Polygones réguliers}
% OnBuch+ — préambule « cours signés » ENRICHI (compilé avec tectonic / XeLaTeX)
% Système visuel « Pop » d'OnBuch+ : fond crème (jamais blanc), bordures encre
% épaisses, ombres « dures » décalées, titres Archivo Black en capitales.
% Version MATHÉMATIQUES : graphes (pgfplots/tikz), boîtes pédagogiques
% (définition, propriété, méthode, attention, le savais-tu, exercice résolu,
% exercice, TP, bilan), page de garde et signature OnBuch+ ancrée en bas.
%
% Macros attendues dans main.tex AVANT \input{preamble} :
%   \DOCMATIERE  \DOCNIVEAU  \DOCMODULE  \DOCLECON (numéro)  \DOCTITRE
\documentclass[11pt]{article}

\usepackage{fontspec}
\usepackage[french]{babel}
\usepackage[a4paper,margin=2cm,top=3.4cm,bottom=2.8cm,headheight=1.4cm,footskip=1.3cm]{geometry}
\usepackage{amsmath,amssymb,amsfonts}
\usepackage{mathtools} % \xmapsto, \coloneqq... souvent utilisés par les modèles
\usepackage{cancel} % simplifications barrées
% \abs{z} (module, valeur absolue) et \norm{u} (norme), utilisés par les modèles
\providecommand{\abs}{}\let\abs\relax\DeclarePairedDelimiter{\abs}{\lvert}{\rvert}
\providecommand{\norm}{}\let\norm\relax\DeclarePairedDelimiter{\norm}{\lVert}{\rVert}
\usepackage[table]{xcolor} % [table] : fournit aussi \rowcolors (alternance de lignes)
\usepackage{pagecolor}
\usepackage{tikz}
\usetikzlibrary{arrows.meta,positioning,calc,shapes.geometric,shapes.misc,shapes.arrows,decorations.pathmorphing,decorations.pathreplacing,decorations.markings,patterns,shadows,mindmap,trees,fit,backgrounds,matrix,intersections,angles,quotes}
\usepackage{pgfplots}
\usepackage{pgf-pie} % diagrammes circulaires \pie (statistiques)
\pgfplotsset{compat=1.18}
\usepgfplotslibrary{fillbetween,groupplots} % aires entre courbes (calcul intégral)
\usepackage{enumitem}
\usepackage{fancyhdr}
% [most] charge toutes les bibliothèques tcolorbox (listings, minted, raster,
% documentation...) et gonfle inutilement la mémoire TeX sur un document long
% (~300 boîtes) : on ne charge que ce qui est réellement utilisé.
\usepackage[skins,breakable]{tcolorbox}
\usepackage{graphicx}
\usepackage{colortbl}
\usepackage{booktabs}
\usepackage{tabularx}
\usepackage{diagbox} % case d'en-tête diagonale des tableaux à double entrée
% Colonnes extensibles centrée / gauche / droite (C, L, R), souvent utilisées par les modèles
\newcolumntype{C}{>{\centering\arraybackslash}X}
\newcolumntype{L}{>{\raggedright\arraybackslash}X}
\newcolumntype{R}{>{\raggedleft\arraybackslash}X}
\usepackage{array}
\usepackage{multicol}
\usepackage{siunitx}
% parse-numbers=false : accepte aussi \SI{2,5 \times 10^{-3}}{...}
\sisetup{locale=FR,output-decimal-marker={,},group-minimum-digits=4,parse-numbers=false}
\DeclareSIUnit\ohm{\ensuremath{\symup{\Omega}}} % U+2126 absent de la police
\usepackage{qrcode}
% \degree : commande fréquemment utilisée par les modèles pour un angle ou une
% température en dehors de \SI{}{} (non fournie par les paquets chargés).
\providecommand{\degree}{^{\circ}}
% \qed (fin de démonstration), utilisé par les modèles sans amsthm
\providecommand{\qed}{\hfill$\square$}
% \barre : double barre des valeurs interdites dans un tableau de variations (tkz-tab)
\providecommand{\barre}{\ensuremath{\|}}
% environnement proof (amsthm non chargé) utilisé par les modèles
\ifdefined\proof\else\newenvironment{proof}[1][Démonstration]{\par\noindent\textit{#1.}\ }{\qed\par}\fi
\usepackage{lastpage}
\usepackage{hyperref}
\hypersetup{hidelinks}

% ============================================================
% Palette « Pop » OnBuch+ — mêmes hex que la classe P (plus_ui.dart)
% ============================================================
\definecolor{poporange}{HTML}{F59321}
\definecolor{popdark}{HTML}{E07A0C}
\definecolor{popcream}{HTML}{FFF6E8}
\definecolor{popink}{HTML}{1C1714}
\definecolor{poppurple}{HTML}{7A5AE0}
\definecolor{popgreen}{HTML}{2E9E6A}
\definecolor{poppink}{HTML}{E0457B}
\definecolor{popblue}{HTML}{2F7DE1}
\definecolor{popgold}{HTML}{C98A12}
\definecolor{popmuted}{HTML}{8A7F76}
\definecolor{popline}{HTML}{EADFD2}
\definecolor{popgray}{HTML}{8A7F76}
\colorlet{popgrey}{popgray}
% Alias tolérants (le modèle nomme parfois les couleurs différemment)
\colorlet{popwhite}{white}
\colorlet{popblack}{popink}
\colorlet{popred}{poppink}
\colorlet{popyellow}{popgold}
% Teintes claires pour fonds de tableaux / zones
\colorlet{popblueL}{popblue!10!white}
\colorlet{poporangeL}{poporange!14!white}
\colorlet{popgreenL}{popgreen!12!white}
\colorlet{poppurpleL}{poppurple!12!white}
\colorlet{poppinkL}{poppink!12!white}
\colorlet{popgoldL}{popgold!14!white}

\pagecolor{popcream}

% ============================================================
% Typographie
% ============================================================
\defaultfontfeatures{Ligatures=TeX}
\setmainfont{PlusJakartaSans-Medium.ttf}[Path=../../../fonts/,BoldFont=PlusJakartaSans-Bold.ttf]
% Maths en sans-serif (Fira Math) avec chiffres et lettres droites de Plus
% Jakarta, pour que les formules se fondent dans le texte.
\usepackage[math-style=ISO,bold-style=ISO]{unicode-math}
\setmathfont{FiraMath-Regular.otf}[Path=../../../fonts/]
\setmathfont{PlusJakartaSans-Medium.ttf}[Path=../../../fonts/,range={up/num,up/{Latin,latin},\mathup}]
\usepackage{icomma} % virgule décimale sans espace en mode math
% Caractères mathématiques Unicode absents des polices de texte (Plus Jakarta,
% Archivo Black) : rendus par leur équivalent mathématique, en texte comme en
% formule — sinon ils s'affichent en carré vide (ex. « Arithmétique dans ℤ »).
\usepackage{newunicodechar}
\newunicodechar{ℝ}{\ensuremath{\mathbb{R}}}
\newunicodechar{ℤ}{\ensuremath{\mathbb{Z}}}
\newunicodechar{ℂ}{\ensuremath{\mathbb{C}}}
\newunicodechar{ℕ}{\ensuremath{\mathbb{N}}}
\newunicodechar{ℚ}{\ensuremath{\mathbb{Q}}}
\newunicodechar{𝔻}{\ensuremath{\mathbb{D}}}
\newunicodechar{α}{\ensuremath{\alpha}}
\newunicodechar{β}{\ensuremath{\beta}}
\newunicodechar{γ}{\ensuremath{\gamma}}
\newunicodechar{δ}{\ensuremath{\delta}}
\newunicodechar{ε}{\ensuremath{\varepsilon}}
\newunicodechar{θ}{\ensuremath{\theta}}
\newunicodechar{λ}{\ensuremath{\lambda}}
\newunicodechar{μ}{\ensuremath{\mu}}
\newunicodechar{σ}{\ensuremath{\sigma}}
\newunicodechar{τ}{\ensuremath{\tau}}
\newunicodechar{φ}{\ensuremath{\varphi}}
\newunicodechar{ψ}{\ensuremath{\psi}}
\newunicodechar{ω}{\ensuremath{\omega}}
\newunicodechar{Σ}{\ensuremath{\Sigma}}
% majuscules produites par \MakeUppercase dans les titres (α-aminés -> Α)
\newunicodechar{Α}{\ensuremath{\alpha}}
\newunicodechar{Β}{\ensuremath{\beta}}
\newunicodechar{Γ}{\ensuremath{\Gamma}}
\newunicodechar{Δ}{\ensuremath{\Delta}}
\newunicodechar{Ε}{\ensuremath{\varepsilon}}
\newunicodechar{Θ}{\ensuremath{\Theta}}
\newunicodechar{Λ}{\ensuremath{\Lambda}}
\newunicodechar{Μ}{\ensuremath{\mu}}
\newunicodechar{Τ}{\ensuremath{\tau}}
\newunicodechar{Φ}{\ensuremath{\Phi}}
\newunicodechar{Ψ}{\ensuremath{\Psi}}
\newunicodechar{Ω}{\ensuremath{\Omega}}
\newunicodechar{↦}{\ensuremath{\mapsto}}
\newunicodechar{∈}{\ensuremath{\in}}
\newunicodechar{∉}{\ensuremath{\notin}}
\newunicodechar{∧}{\ensuremath{\wedge}}
\newunicodechar{⇔}{\ensuremath{\Leftrightarrow}}
\newunicodechar{⟺}{\ensuremath{\Longleftrightarrow}}
\newunicodechar{⇒}{\ensuremath{\Rightarrow}}
\newunicodechar{⟹}{\ensuremath{\Longrightarrow}}
\newunicodechar{∘}{\ensuremath{\circ}}
\newunicodechar{∪}{\ensuremath{\cup}}
\newunicodechar{∩}{\ensuremath{\cap}}
\newunicodechar{≡}{\ensuremath{\equiv}}
\newunicodechar{⊂}{\ensuremath{\subset}}
\newunicodechar{⊥}{\ensuremath{\perp}}
\newunicodechar{∃}{\ensuremath{\exists}}
\newunicodechar{∀}{\ensuremath{\forall}}
\newunicodechar{∛}{\ensuremath{\sqrt[3]{\,}}}
\newunicodechar{∜}{\ensuremath{\sqrt[4]{\,}}}
\newunicodechar{⃗}{\ensuremath{{}^{\to}}}
\newunicodechar{ⁿ}{\ifmmode^{n}\else\textsuperscript{\ensuremath{n}}\fi}
\newunicodechar{ˣ}{\ifmmode^{x}\else\textsuperscript{\ensuremath{x}}\fi}
\newunicodechar{ᵇ}{\ifmmode^{b}\else\textsuperscript{\ensuremath{b}}\fi}
\newunicodechar{ʳ}{\ifmmode^{r}\else\textsuperscript{\ensuremath{r}}\fi}
\newunicodechar{ᵘ}{\ifmmode^{u}\else\textsuperscript{\ensuremath{u}}\fi}
\newunicodechar{ʸ}{\ifmmode^{y}\else\textsuperscript{\ensuremath{y}}\fi}
\newunicodechar{ᵃ}{\ifmmode^{a}\else\textsuperscript{\ensuremath{a}}\fi}
\newunicodechar{ᵉ}{\ifmmode^{e}\else\textsuperscript{\ensuremath{e}}\fi}
\newunicodechar{ᵏ}{\ifmmode^{k}\else\textsuperscript{\ensuremath{k}}\fi}
\newunicodechar{ᵖ}{\ifmmode^{p}\else\textsuperscript{\ensuremath{p}}\fi}
\newunicodechar{ᴺ}{\ifmmode^{N}\else\textsuperscript{\ensuremath{N}}\fi}
\newunicodechar{ᶜ}{\ifmmode^{c}\else\textsuperscript{\ensuremath{c}}\fi}
\newunicodechar{ʰ}{\ifmmode^{h}\else\textsuperscript{\ensuremath{h}}\fi}
\newunicodechar{ᵠ}{\ifmmode^{q}\else\textsuperscript{\ensuremath{q}}\fi}
\newunicodechar{ₙ}{\ifmmode_{n}\else\textsubscript{\ensuremath{n}}\fi}
\newunicodechar{ₐ}{\ifmmode_{a}\else\textsubscript{\ensuremath{a}}\fi}
\newunicodechar{ᵢ}{\ifmmode_{i}\else\textsubscript{\ensuremath{i}}\fi}
\newunicodechar{ᵣ}{\ifmmode_{r}\else\textsubscript{\ensuremath{r}}\fi}
\newunicodechar{ₖ}{\ifmmode_{k}\else\textsubscript{\ensuremath{k}}\fi}
\newunicodechar{ᵦ}{\ifmmode_{\beta}\else\textsubscript{\ensuremath{\beta}}\fi}
\newunicodechar{ₓ}{\ifmmode_{x}\else\textsubscript{\ensuremath{x}}\fi}
\newfontfamily\popdisplay{ArchivoBlack-Regular.ttf}[Path=../../../fonts/]
\newfontfamily\poplogo{SpaceGrotesk-Bold.ttf}[Path=../../../fonts/]
\newfontfamily\popsemi{PlusJakartaSans-SemiBold.ttf}[Path=../../../fonts/]
\newfontfamily\popxbold{PlusJakartaSans-ExtraBold.ttf}[Path=../../../fonts/]
\newfontfamily\popxboldls{PlusJakartaSans-ExtraBold.ttf}[Path=../../../fonts/,LetterSpace=3.0]

\setlist[enumerate]{leftmargin=1.7em,itemsep=5pt,topsep=4pt}
\setlist[itemize]{leftmargin=1.7em,itemsep=4pt,topsep=4pt,label={\color{poporange}\textbullet}}
\setlength{\parindent}{0pt}
\setlength{\parskip}{5pt}
\renewcommand{\arraystretch}{1.25}

% pgfplots : style Pop par défaut
\pgfplotsset{
  every axis/.append style={axis line style={popink,line width=1pt},tick style={popink},
    grid style={popline,line width=0.5pt},label style={font=\small},tick label style={font=\footnotesize}},
  popcurve/.style={poporange,line width=1.8pt,no markers,smooth},
}
% Même style disponible dans un \draw TikZ simple (hors pgfplots)
\tikzset{popcurve/.style={poporange,line width=1.8pt}}

% ============================================================
% Pill / squiggle
% ============================================================
\newcommand{\poppill}[3]{%
  \tikz[baseline=-0.55ex]\node[draw=popink,line width=1.2pt,rounded corners=2.6pt,%
    fill=#2,inner xsep=6.5pt,inner ysep=3pt]{%
    \popxboldls\fontsize{7}{7}\selectfont\color{#3}\MakeUppercase{#1}};%
}
\newcommand{\popsquiggle}[1]{%
  \tikz[baseline=-0.4ex]{
    \draw[#1,line width=2.6pt,line cap=round]
      plot [smooth] coordinates {(0,0.09) (0.34,0.01) (0.68,0.21) (1.02,0.01) (1.36,0.10)};
  }%
}
% Mot-clé surligné (marqueur orange sous le texte)
\newcommand{\cle}[1]{{\popxbold\color{popdark}#1}}

% ============================================================
% En-tête / pied
% ============================================================
\pagestyle{fancy}
\fancyhf{}
\renewcommand{\headrulewidth}{0pt}
\renewcommand{\footrulewidth}{0pt}
\lhead{\raisebox{-0.15ex}{\poplogo\fontsize{13}{13}\selectfont\color{popink}OnBuch\color{poporange}+}}
\rhead{\poppill{\DOCMATIERE\ \textbullet\ \DOCNIVEAU\ \textbullet\ Leçon \DOCLECON}{white}{popink}}
\lfoot{\popxbold\fontsize{7.6}{7.6}\selectfont\color{popmuted}\MakeUppercase{Cours signé OnBuch+}\ \textbullet\ {\popsemi\DOCTITRE}}
\rfoot{\tikz[baseline=-0.6ex]\node[rounded corners=8.5pt,fill=popink,minimum height=17pt,minimum width=17pt,inner xsep=5pt,inner sep=0]{\popxbold\fontsize{8}{8}\selectfont\color{popcream}\thepage\,/\,\pageref*{LastPage}};}

% ============================================================
% Titres
% ============================================================
\newcommand{\chaptitre}[2]{%
  \begin{tcolorbox}[enhanced,colback=poporange,colframe=popink,boxrule=2.6pt,arc=11pt,
      drop shadow=popink,left=15pt,right=15pt,top=12pt,bottom=11pt,before skip=3pt,after skip=18pt]
    \poppill{#2}{popink}{popcream}\par\vspace{9pt}
    {\raggedright\hyphenpenalty=10000\exhyphenpenalty=10000\popdisplay\fontsize{22}{24}\selectfont\color{popcream}\MakeUppercase{#1}\par}
  \end{tcolorbox}
}
% Section numérotée : pastille orange avec le numéro + titre Archivo
\newcounter{popsec}
\newcommand{\coursec}[1]{%
  \stepcounter{popsec}\setcounter{popsub}{0}%
  \par\vspace{16pt}\noindent
  \begin{minipage}{\linewidth}
  \tikz[baseline=-0.7ex]\node[draw=popink,line width=1.6pt,fill=poporange,rounded corners=4pt,minimum size=22pt,inner sep=0]{\popdisplay\fontsize{12}{12}\selectfont\color{popcream}\thepopsec};%
  \hspace{8pt}{\popdisplay\fontsize{15}{17}\selectfont\color{popink}\MakeUppercase{#1}}\par
  \vspace{4pt}\noindent{\color{poporange}\rule{36pt}{3.6pt}}
  \end{minipage}\par\nopagebreak\vspace{8pt}
}
\newcounter{popsub}
\newcommand{\courssub}[1]{%
  \stepcounter{popsub}%
  \par\vspace{9pt}\noindent{\popxbold\fontsize{11.5}{13}\selectfont\color{popink}{\color{poporange}\thepopsec.\thepopsub}\hspace{6pt}#1}\par\nopagebreak\vspace{4pt}
}

% ============================================================
% Boîtes pédagogiques — même recette : fond blanc, bordure épaisse colorée,
% ombre dure encre, badge plein. Titre optionnel [#1].
% ============================================================
\tcbset{popbase/.style={breakable,enhanced,colback=white,boxrule=2.3pt,arc=9pt,
  shadow={3pt}{-3pt}{0pt}{popink},left=14pt,right=14pt,top=11pt,bottom=11pt,before skip=11pt,after skip=15pt,
  fontupper=\popsemi\fontsize{10.6}{14.5}\selectfont\color{popink},
  fontlower=\popsemi\fontsize{10.6}{14.5}\selectfont\color{popink}}}
% Coche dessinée (le glyphe ✓ n'existe pas dans les polices du document)
\newcommand{\popcheck}[1]{\tikz[baseline=-0.5ex]\draw[#1,line width=1.6pt,line cap=round,line join=round](-0.13,0.0)--(-0.04,-0.09)--(0.13,0.1);}
\newcommand{\popboxhead}[3]{\poppill{#1}{#2}{white}\ifx\relax#3\relax\else\hspace{6pt}{\popxbold\fontsize{10.6}{12}\selectfont\color{popink}#3}\fi\par\vspace{7pt}}

\newtcolorbox{aretenir}[1][]{popbase,colframe=popgreen,before upper={\popboxhead{À retenir}{popgreen}{#1}}}
\newtcolorbox{exemplebox}[1][]{popbase,colframe=poporange,before upper={\popboxhead{Exemple}{poporange}{#1}}}
\newtcolorbox{definition}[1][]{popbase,colframe=popblue,before upper={\popboxhead{Définition}{popblue}{#1}}}
\newtcolorbox{propriete}[1][]{popbase,colframe=poppurple,before upper={\popboxhead{Propriété}{poppurple}{#1}}}
\newtcolorbox{methode}[1][]{popbase,colframe=popgold,before upper={\popboxhead{Méthode}{popgold}{#1}}}
\newtcolorbox{attention}[1][]{popbase,colframe=poppink,before upper={\popboxhead{Attention}{poppink}{#1}}}
\newtcolorbox{experience}[1][]{popbase,colframe=popblue,colback=popblueL,before upper={\popboxhead{Expérience}{popblue}{#1}}}
% « Le savais-tu ? » : carte violette pleine (contexte, culture, Cameroun)
\newtcolorbox{savaistu}[1][]{popbase,colback=poppurple,colframe=popink,
  fontupper=\popsemi\fontsize{10.4}{14.2}\selectfont\color{white},
  before upper={\poppill{Le savais-tu\,?}{white}{poppurple}\ifx\relax#1\relax\else\hspace{6pt}{\popxbold\color{white}#1}\fi\par\vspace{7pt}}}
% Exercice résolu : énoncé en haut, \tcblower puis la solution sur fond crème
\newcounter{popexo}\newcounter{popexr}
\newtcolorbox{exoresolu}[1][]{popbase,colframe=poporange,colbacklower=popcream,
  segmentation style={popink,line width=1.2pt,dashed},
  before upper={\stepcounter{popexr}\popboxhead{Exercice résolu \thepopexr}{poporange}{#1}},
  before lower={\poppill{Solution}{popink}{popcream}\par\vspace{6pt}}}
% Exercice (non résolu en place) — la correction est dans la partie Corrigés
\newtcolorbox{exercice}[1][]{popbase,colframe=popink,
  before upper={\stepcounter{popexo}\popboxhead{Exercice \thepopexo}{popink}{#1}}}
% Corrigé
\newtcolorbox{corrige}[1][]{popbase,colframe=popgreen,colback=popgreenL,
  before upper={\popboxhead{Corrigé}{popgreen}{#1}}}
% Objectifs / prérequis / bilan
\newtcolorbox{objectifs}{popbase,colframe=popink,colback=poporangeL,
  before upper={\popboxhead{Objectifs de la leçon}{popink}{}}}
\newtcolorbox{prerequis}{popbase,colframe=popink,colback=popblueL,
  before upper={\popboxhead{Prérequis}{popblue}{}}}
\newtcolorbox{bilan}{popbase,colframe=popink,colback=white,boxrule=2.8pt,
  before upper={\popboxhead{Fiche bilan}{poporange}{L'essentiel à connaître pour l'examen}}}
% Figure encadrée avec légende
\newtcolorbox{popfigure}[1][]{enhanced,breakable=false,colback=white,colframe=popline,boxrule=1.4pt,arc=8pt,
  left=10pt,right=10pt,top=10pt,bottom=8pt,before skip=10pt,after skip=12pt,halign=center,
  lower separated=false,halign lower=center,
  fontlower=\popsemi\fontsize{9}{11.5}\selectfont\color{popmuted},
  #1}
\newcommand{\legende}[1]{\par\vspace{6pt}{\popsemi\fontsize{9}{11.5}\selectfont\color{popmuted}#1}}

% ============================================================
% Signature OnBuch+
% ============================================================
\newcommand{\poplogoinline}[1]{{\poplogo\fontsize{#1}{#1}\selectfont\color{popink}OnBuch\color{poporange}+}}
\newcommand{\signatureonbuch}{%
  \begin{tcolorbox}[enhanced,colback=popink,colframe=popink,boxrule=2.2pt,arc=10pt,
      drop shadow=poporange,left=14pt,right=14pt,top=10pt,bottom=10pt,before skip=0pt,after skip=0pt]
    \tikz[baseline=-0.7ex]\node[circle,fill=popgreen,draw=popcream,line width=1.3pt,minimum size=22pt,inner sep=0]{\popcheck{white}};%
    \hspace{9pt}\begin{minipage}[c]{0.62\linewidth}
      {\popxbold\fontsize{12}{13}\selectfont\color{popcream}Signé\ \textbullet\ {\poplogo OnBuch\color{poporange}+}}\par\vspace{2pt}
      {\raggedright\popsemi\fontsize{8}{10.5}\selectfont\color{popline}\MakeUppercase{Équipe pédagogique OnBuch+}\\\MakeUppercase{Conforme au programme officiel MINESEC}\par}
    \end{minipage}\hfill
    {\poplogo\fontsize{22}{22}\selectfont\color{popcream}OnBuch\color{poporange}+}
  \end{tcolorbox}%
}

% ============================================================
% Page de garde
% #1 titre · #2 matière · #3 niveau · #4 module · #5 n° leçon · #6 durée ·
% #7 description / objectifs (texte ou itemize)
% ============================================================
\newcommand{\pagegarde}[7]{%
  \thispagestyle{empty}
  \newgeometry{a4paper,margin=1.7cm,top=1.6cm,bottom=1.4cm}%
  \begin{tikzpicture}[remember picture,overlay]
    \fill[poporange!18] ([xshift=-3.2cm,yshift=-3.6cm]current page.north east) circle (4.6cm);
    \fill[poppurple!14] ([xshift=2.4cm,yshift=7.6cm]current page.south west) circle (3.8cm);
  \end{tikzpicture}%
  {\poplogo\fontsize{30}{30}\selectfont\color{popink}OnBuch\color{poporange}+}\hfill
  \raisebox{0.6ex}{\poppill{Cours signé \textbullet\ Premium}{popink}{popcream}}\par
  \vspace{1pt}\popsquiggle{poppurple}\par
  \vspace{8pt}
  \begin{tcolorbox}[enhanced,colback=poporange,colframe=popink,boxrule=2.8pt,arc=13pt,
      drop shadow=popink,left=18pt,right=18pt,top=12pt,bottom=13pt,after skip=10pt]
    \poppill{#2 \textbullet\ #3}{popink}{popcream}\hspace{5pt}\poppill{#4}{white}{popink}\par\vspace{8pt}
    {\popdisplay\fontsize{14}{14}\selectfont\color{popink}LEÇON #5}\par\vspace{4pt}
    {\raggedright\hyphenpenalty=10000\exhyphenpenalty=10000\popdisplay\fontsize{25}{27}\selectfont\color{popcream}\MakeUppercase{#1}\par}
    \vspace{8pt}
    {\popxbold\fontsize{9.5}{9.5}\selectfont\color{popink}\MakeUppercase{Durée conseillée : #6}}
  \end{tcolorbox}
  \begin{tcolorbox}[enhanced,colback=white,colframe=popink,boxrule=2.2pt,arc=10pt,
      drop shadow=popink,left=16pt,right=16pt,top=10pt,bottom=10pt,after skip=8pt]
    \poppill{Dans cette leçon}{poppurple}{white}\par\vspace{5pt}
    \popsemi\fontsize{9.3}{12.6}\selectfont\color{popink}#7
  \end{tcolorbox}
  \noindent\begin{minipage}[t]{0.487\linewidth}
    \begin{tcolorbox}[enhanced,colback=popgreen,colframe=popink,boxrule=2.2pt,arc=10pt,
        drop shadow=popink,left=11pt,right=11pt,top=10pt,bottom=10pt,height=3.1cm,valign=center]
      \tcbox[colback=white,colframe=popink,boxrule=1.3pt,arc=5pt,boxsep=0pt,nobeforeafter,
        left=3pt,right=3pt,top=3pt,bottom=3pt,valign=center]{\qrcode[height=1.9cm]{https://play.google.com/store/apps/details?id=com.luvvix.onbuch&hl=fr}}%
      \hspace{8pt}\begin{minipage}[c]{0.5\linewidth}
        {\popdisplay\fontsize{11}{12.5}\selectfont\color{white}\MakeUppercase{Télécharge OnBuch}}\par\vspace{4pt}
        {\raggedright\popsemi\fontsize{8.4}{11}\selectfont\color{white}Gratuit \textbullet\ Google Play\\Tuteur IA, annales, concours, résultats\par}
      \end{minipage}
    \end{tcolorbox}
  \end{minipage}\hfill
  \begin{minipage}[t]{0.487\linewidth}
    \begin{tcolorbox}[enhanced,colback=poppurple,colframe=popink,boxrule=2.2pt,arc=10pt,
        drop shadow=popink,left=11pt,right=11pt,top=10pt,bottom=10pt,height=3.1cm,valign=center]
      \tikz[baseline=-0.7ex]\node[circle,fill=white,draw=popink,line width=1.3pt,minimum size=1.6cm,inner sep=0]{\popdisplay\fontsize{24}{24}\selectfont\color{poppurple}+};%
      \hspace{8pt}\begin{minipage}[c]{0.6\linewidth}
        {\popdisplay\fontsize{11}{12.5}\selectfont\color{white}\MakeUppercase{Passe à OnBuch+}}\par\vspace{4pt}
        {\raggedright\popsemi\fontsize{8.4}{11}\selectfont\color{white}Tous les cours signés, Tuteur IA illimité, fiches PDF — onglet OnBuch+ de l'app\par}
      \end{minipage}
    \end{tcolorbox}
  \end{minipage}
  \vspace{8pt}
  \popcontenu
  \vfill
  \signatureonbuch
  \restoregeometry
  \newpage
}

% Rangée « Ce cours contient » (page de garde)
\newcommand{\poptile}[3]{% #1 couleur #2 gros texte #3 légende
  \begin{tcolorbox}[enhanced,colback=#1,colframe=popink,boxrule=2pt,arc=9pt,shadow={2.5pt}{-2.5pt}{0pt}{popink},
      left=6pt,right=6pt,top=9pt,bottom=9pt,halign=center,valign=center,height=1.75cm,width=\linewidth,nobeforeafter]
    {\popdisplay\fontsize{12.5}{13}\selectfont\color{white}\MakeUppercase{#2}}\par\vspace{3pt}
    {\centering\popsemi\fontsize{7.8}{9.5}\selectfont\color{white}#3\par}
  \end{tcolorbox}}
\newcommand{\popcontenu}{%
  \par\noindent{\popxboldls\fontsize{7.5}{7.5}\selectfont\color{popmuted}CE COURS SIGNÉ CONTIENT}\par\vspace{7pt}
  \noindent\begin{minipage}[t]{0.235\linewidth}\poptile{popblue}{Cours}{complet, illustré, pas à pas}\end{minipage}\hfill
  \begin{minipage}[t]{0.235\linewidth}\poptile{poporange}{Méthodes}{et exercices résolus}\end{minipage}\hfill
  \begin{minipage}[t]{0.235\linewidth}\poptile{poppink}{Exercices}{type examen + corrigés}\end{minipage}\hfill
  \begin{minipage}[t]{0.235\linewidth}\poptile{popgreen}{Fiche bilan}{l'essentiel à retenir}\end{minipage}\par}

% Sommaire
\newtcolorbox{sommaire}{enhanced,colback=white,colframe=popink,boxrule=2.2pt,arc=10pt,
  drop shadow=popink,left=18pt,right=18pt,top=13pt,bottom=13pt,before skip=6pt,after skip=20pt,
  before upper={\poppill{Sommaire}{poppurple}{white}\par\vspace{9pt}\popsemi\fontsize{10.6}{14.5}\selectfont\color{popink}}}

% Signature de fin de cours — poussée tout en bas de la dernière page
\newcommand{\findecours}{%
  \par\vfill\nopagebreak
  \begin{center}\popsquiggle{poporange}\end{center}\vspace{4pt}
  \signatureonbuch
}
\AtBeginDocument{\def\llbracket{\mathopen{[\mkern-3.5mu[}}\def\rrbracket{\mathclose{]\mkern-3.5mu]}}} % intervalles d'entiers (glyphe ⟦ absent de la police)
% Glyphes absents de Fira Math : redéfinis à partir de glyphes disponibles
\AtBeginDocument{%
  \def\setminus{\mathbin{\backslash}}\let\smallsetminus\setminus
  \def\vdots{\mathord{\vbox{\baselineskip3.2pt\lineskiplimit0pt\kern4pt\hbox{.}\hbox{.}\hbox{.}}}}%
  \def\ddots{\mathinner{\mkern1mu\raise6.5pt\hbox{.}\mkern2mu\raise3.6pt\hbox{.}\mkern2mu\raise0.7pt\hbox{.}\mkern1mu}}%
  \def\triangle{\mathord{\scalebox{0.9}{$\symup{\Delta}$}}}%
  \def\nabla{\mathord{\raisebox{\depth}{\scalebox{1}[-1]{$\symup{\Delta}$}}}}%
  \def\checkmark{\popcheck{popdark}}%
  \def\star{\mathbin{\ast}}\def\bigstar{\mathord{\ast}}%
  \def\diamond{\mathbin{\scalebox{0.6}{$\Diamond$}}}%
  \def\bigcup{\mathop{\mathchoice{\raisebox{-0.3em}{\scalebox{1.7}{$\cup$}}}{\scalebox{1.25}{$\cup$}}{\cup}{\cup}}\displaylimits}%
  \def\bigcap{\mathop{\mathchoice{\raisebox{-0.3em}{\scalebox{1.7}{$\cap$}}}{\scalebox{1.25}{$\cap$}}{\cap}{\cap}}\displaylimits}%
  \def\lesseqqgtr{\mathrel{\lessgtr}}\def\bigoplus{\mathop{\oplus}\displaylimits}%
}
\providecommand{\ssi}{si et seulement si} % abréviation fréquente
\AtBeginDocument{\providecommand{\wideparen}{\overparen}} % arc de cercle

%% FIGFIX-BEGIN v4 — correctifs globaux des figures (docs/onbuch-generation/tools/figfix)
\usepackage{adjustbox}\usepackage{etoolbox}
% toute figure TikZ/pgfplots plus large que la ligne est réduite à la largeur disponible
\BeforeBeginEnvironment{tikzpicture}{\begin{adjustbox}{max width=\linewidth}}
\AfterEndEnvironment{tikzpicture}{\end{adjustbox}}
% légendes pgfplots lisibles par-dessus les courbes
\pgfplotsset{every axis/.append style={legend style={fill=white,fill opacity=0.92,text opacity=1,draw=black!15,inner sep=3pt}}}
% étiquettes d'arêtes (syntaxe edge["..."]) avec fond blanc
\tikzset{every edge quotes/.append style={fill=white,inner sep=1.2pt}}
%% FIGFIX-END
\begin{document}

\pagegarde{Polygones réguliers}{Mathématiques}{3ème}{Mathématiques – classe de 3ème}{N13}{3 séances (fiche de progression harmonisée)}{Cette leçon étudie les polygones réguliers, figures dont tous les côtés et tous les angles sont égaux. Elle détaille les propriétés, les constructions et les calculs (angles, apothème, périmètre, aire) pour le triangle équilatéral, le carré, l'hexagone et l'octogone réguliers, avec des applications concrètes à la vie courante.
\vspace{4pt}
\begin{multicols}{2}\small
\begin{itemize}
\item À la fin de cette leçon, je sais définir un polygone régulier et reconnaître si un polygone est régulier
\item je sais calculer la mesure d'un angle au centre et d'un angle intérieur d'un polygone régulier à n côtés
\item je sais construire un triangle équilatéral, un carré, un hexagone régulier et un octogone régulier à la règle et au compas
\item je sais calculer la hauteur d'un triangle équilatéral et la diagonale d'un carré en fonction du côté
\item je sais calculer le périmètre et l'aire d'un polygone régulier à l'aide de l'apothème
\item je sais justifier qu'un polygone régulier est inscriptible dans un cercle
\end{itemize}\end{multicols}}

\begin{sommaire}
\begin{multicols}{2}
\begin{enumerate}
\item Définitions et premières propriétés
\item Angles dans un polygone régulier
\item Triangle équilatéral et carré
\item Hexagone et octogone réguliers
\item Apothème, périmètre et aire d'un polygone régulier
\item Applications concrètes et synthèse
\item Activité d'intégration
\item Méthodes et exercices résolus
\item Exercices
\item Corrigés des exercices
\item Fiche bilan
\end{enumerate}
\end{multicols}
\end{sommaire}

\begin{objectifs}
\begin{itemize}
\item À la fin de cette leçon, je sais définir un polygone régulier et reconnaître si un polygone est régulier
\item je sais calculer la mesure d'un angle au centre et d'un angle intérieur d'un polygone régulier à n côtés
\item je sais construire un triangle équilatéral, un carré, un hexagone régulier et un octogone régulier à la règle et au compas
\item je sais calculer la hauteur d'un triangle équilatéral et la diagonale d'un carré en fonction du côté
\item je sais calculer le périmètre et l'aire d'un polygone régulier à l'aide de l'apothème
\item je sais justifier qu'un polygone régulier est inscriptible dans un cercle
\item je sais résoudre des problèmes concrets (pavage, découpe, aménagement) mettant en jeu des polygones réguliers
\item je sais rédiger et justifier une démarche complète de résolution de problème
\end{itemize}
\end{objectifs}

\begin{prerequis}
\begin{itemize}
\item Théorème de Pythagore et sa réciproque (4ème)
\item Angles au centre et angles inscrits dans un cercle (début d'année de 3ème)
\item Médiatrice d'un segment, bissectrice d'un angle et leurs constructions au compas
\item Périmètres et aires des figures usuelles (triangle, carré, disque)
\item Trigonométrie dans le triangle rectangle : sinus, cosinus, tangente (3ème)
\end{itemize}
\end{prerequis}

\newpage

\coursec{Définitions et premières propriétés}

Maman Aïcha veut carreler la cour de sa maison à Yaoundé avec des pavés hexagonaux, comme ceux qu'elle a vus au marché Mokolo. Le vendeur lui affirme que ces pavés s'emboîtent parfaitement sans laisser de trous, contrairement aux pavés pentagonaux qu'elle avait d'abord choisis. Intriguée, elle se demande : pourquoi certains polygones réguliers permettent-ils de paver une surface plane et pas d'autres ? Comment tracer précisément un hexagone ou un octogone régulier avec seulement une règle et un compas ?

Pour répondre à ces questions concrètes, il faut d'abord bien comprendre ce qu'est un polygone régulier, quelles sont ses propriétés remarquables, et quel vocabulaire utiliser pour le décrire. C'est l'objectif de cette première section : poser des définitions solides et découvrir les premières propriétés qui font des polygones réguliers des figures exceptionnelles en géométrie.

\courssub{Qu'est-ce qu'un polygone ?}

\textbf{Intuition.} Ferme les yeux et imagine la forme d'une feuille de papier, d'un panneau « STOP » à l'entrée d'une ville, ou de la case d'une ruche d'abeille. Ces formes ont un point commun : elles sont limitées par des segments de droite qui se succèdent et se rejoignent bout à bout. Ce sont des \cle{polygones}.

\begin{definition}[Polygone]
Un \cle{polygone} est une figure plane fermée, formée d'une suite de segments appelés \cle{côtés}, telle que deux côtés consécutifs ont une et une seule extrémité commune.
\begin{itemize}
\item Les extrémités communes des côtés sont appelées les \cle{sommets} du polygone.
\item Une \cle{diagonale} est un segment qui relie deux sommets \emph{non consécutifs}.
\item Un polygone est dit \cle{convexe} lorsque toutes ses diagonales sont entièrement situées à l'intérieur du polygone (autrement dit, il n'a pas de « creux »).
\end{itemize}
Un polygone à $n$ côtés possède $n$ sommets et $n$ angles.
\end{definition}

\begin{exemplebox}[Quelques polygones connus]
\begin{itemize}
\item Le \textbf{triangle} : 3 côtés, 3 sommets, aucune diagonale.
\item Le \textbf{quadrilatère} : 4 côtés, 2 diagonales (par exemple $[AC]$ et $[BD]$ dans le quadrilatère $ABCD$).
\item Le \textbf{pentagone} : 5 côtés, 5 diagonales.
\item L'\textbf{hexagone} : 6 côtés, 9 diagonales.
\item L'\textbf{octogone} : 8 côtés, 20 diagonales.
\end{itemize}
Tous ces polygones, dans leur version « régulière », seront évoqués dans cette leçon ; le programme de 3ème se concentre particulièrement sur le triangle équilatéral, le carré, l'hexagone et l'octogone réguliers.
\end{exemplebox}

\begin{attention}[Polygone convexe ou non ?]
Dans toute cette leçon, on ne considère que des polygones \textbf{convexes}. Un polygone qui a un « creux » (un angle rentrant, supérieur à $180^\circ$) n'est pas convexe : certaines de ses diagonales passent à l'extérieur de la figure. Vérifie toujours que ta figure n'a pas de creux avant d'appliquer les propriétés de cette leçon.
\end{attention}

\courssub{Définition d'un polygone régulier}

\textbf{Intuition.} Parmi tous les quadrilatères, le carré semble « parfait » : ses quatre côtés ont la même longueur et ses quatre angles ont la même mesure. C'est cette double régularité — sur les longueurs \emph{et} sur les angles — qui définit un polygone régulier.

\begin{definition}[Polygone régulier]
Un polygone est dit \cle{régulier} lorsque :
\begin{itemize}
\item tous ses \textbf{côtés ont la même longueur} ;
\item tous ses \textbf{angles ont la même mesure}.
\end{itemize}
Ces \textbf{deux conditions sont indispensables} : une seule ne suffit pas.
\end{definition}

\begin{attention}[Erreur fréquente : « des côtés égaux suffisent »]
Beaucoup d'élèves croient qu'il suffit que tous les côtés soient égaux pour qu'un polygone soit régulier. C'est \textbf{faux} !
\begin{itemize}
\item Un \textbf{losange qui n'est pas un carré} a ses quatre côtés de même longueur, mais ses angles ne sont pas tous égaux (deux angles aigus et deux angles obtus) : il n'est \textbf{pas régulier}.
\item Réciproquement, un \textbf{rectangle qui n'est pas un carré} a ses quatre angles égaux à $90^\circ$, mais ses côtés n'ont pas tous la même longueur : il n'est \textbf{pas régulier} non plus.
\end{itemize}
Retiens bien : il faut \textbf{à la fois} l'égalité des côtés \textbf{et} l'égalité des angles.
\end{attention}

\begin{popfigure}
\begin{center}
\begin{tikzpicture}[scale=0.9]
% Losange non régulier
\begin{scope}[xshift=-4cm]
\coordinate (A) at (0,1.4);
\coordinate (B) at (2,0);
\coordinate (C) at (0,-1.4);
\coordinate (D) at (-2,0);
\draw[thick,popblue] (A)--(B)--(C)--(D)--cycle;
\fill[popblue] (A) circle (1.5pt) node[above]{\small $A$};
\fill[popblue] (B) circle (1.5pt) node[right]{\small $B$};
\fill[popblue] (C) circle (1.5pt) node[below]{\small $C$};
\fill[popblue] (D) circle (1.5pt) node[left]{\small $D$};
\node[popdark] at (0,0.55) {\small $110^\circ$};
\node[popdark] at (0,-0.55) {\small $110^\circ$};
\node[popdark] at (1.35,0) {\small $70^\circ$};
\node[popdark] at (-1.35,0) {\small $70^\circ$};
\node[popink] at (0.95,0.85) {\small $4$ cm};
\node[popink] at (-0.95,0.85) {\small $4$ cm};
\node[popink] at (0.95,-0.85) {\small $4$ cm};
\node[popink] at (-0.95,-0.85) {\small $4$ cm};
\node[below,popdark,align=center] at (0,-1.9) {\small Losange : côtés égaux, angles inégaux};
\node[below,popdark,align=center] at (0,-2.4) {\small $\Rightarrow$ \textbf{non régulier}};
\end{scope}
% Rectangle non régulier
\begin{scope}[xshift=4cm]
\coordinate (E) at (-2,1);
\coordinate (F) at (2,1);
\coordinate (G) at (2,-1);
\coordinate (H) at (-2,-1);
\draw[thick,poporange] (E)--(F)--(G)--(H)--cycle;
\fill[poporange] (E) circle (1.5pt) node[above left]{\small $E$};
\fill[poporange] (F) circle (1.5pt) node[above right]{\small $F$};
\fill[poporange] (G) circle (1.5pt) node[below right]{\small $G$};
\fill[poporange] (H) circle (1.5pt) node[below left]{\small $H$};
\node[popdark] at (0,1.35) {\small $6$ cm};
\node[popdark] at (0,-1.35) {\small $6$ cm};
\node[popdark] at (2.45,0) {\small $3$ cm};
\node[popdark] at (-2.45,0) {\small $3$ cm};
\node[popdark] at (-1.55,0.65) {\small $90^\circ$};
\node[popdark] at (1.55,0.65) {\small $90^\circ$};
\node[popdark] at (1.55,-0.65) {\small $90^\circ$};
\node[popdark] at (-1.55,-0.65) {\small $90^\circ$};
\node[below,popdark,align=center] at (0,-1.9) {\small Rectangle : angles égaux, côtés inégaux};
\node[below,popdark,align=center] at (0,-2.4) {\small $\Rightarrow$ \textbf{non régulier}};
\end{scope}
\end{tikzpicture}
\end{center}
\legende{Deux contre-exemples : ni le losange non carré, ni le rectangle non carré ne sont des polygones réguliers.}
\end{popfigure}

\begin{exoresolu}[Reconnaître un polygone régulier]
Énoncé. Un quadrilatère $ABCD$ a quatre côtés de longueur $5$ cm et un angle $\widehat{A} = 90^\circ$. Peut-on affirmer que $ABCD$ est un polygone régulier ?
\tcblower
Solution détaillée. Un quadrilatère dont les quatre côtés sont égaux est un losange. Or un losange qui possède un angle droit a ses quatre angles droits : en effet, dans un losange, les angles opposés sont égaux et les angles consécutifs sont supplémentaires, donc si $\widehat{A} = 90^\circ$, alors $\widehat{C} = 90^\circ$ et $\widehat{B} = \widehat{D} = 180^\circ - 90^\circ = 90^\circ$. Le quadrilatère $ABCD$ est donc un \textbf{carré} : ses quatre côtés sont égaux ($5$ cm) et ses quatre angles sont égaux ($90^\circ$). Conclusion : $ABCD$ \textbf{est un polygone régulier} (un carré).
\end{exoresolu}

\courssub{Centre, rayon et apothème d'un polygone régulier}

\textbf{Intuition.} Prends un hexagone régulier, comme un écrou de boulon. Si tu plantes une pointe de compas en son milieu, tu peux tracer un cercle qui passe exactement par les six sommets. Ce point privilégié s'appelle le \emph{centre} du polygone, et il va nous permettre de décrire toutes les dimensions de la figure.

\begin{definition}[Centre, rayon, apothème]
Soit un polygone régulier à $n$ côtés.
\begin{itemize}
\item Le \cle{centre} $O$ du polygone régulier est le point équidistant de tous ses sommets.
\item Le \cle{rayon} du polygone régulier est le rayon $R$ du cercle qui passe par tous ses sommets : c'est la distance $OA$ entre le centre et un sommet.
\item L'\cle{apothème} est la distance $a$ du centre $O$ à un côté : c'est la longueur $OH$, où $H$ est le pied de la perpendiculaire abaissée de $O$ sur un côté. Le point $H$ est le \textbf{milieu} de ce côté.
\end{itemize}
\end{definition}

\begin{popfigure}
\begin{center}
\begin{tikzpicture}[scale=1.15]
\coordinate (O) at (0,0);
\draw[popline] (0,0) circle (2);
\foreach \i in {0,1,2,3,4,5}{
  \coordinate (S\i) at (90-60*\i:2);
  \fill[popblue] (S\i) circle (1.6pt);
}
\draw[thick,popblue] (S0)--(S1)--(S2)--(S3)--(S4)--(S5)--cycle;
\fill[popink] (0,0) circle (1.8pt) node[below left=1pt]{\small $O$};
\draw[thick,popgreen] (0,0)--(S0) node[midway,above right=-1pt]{\small $R$};
\coordinate (H) at ($(S0)!0.5!(S1)$);
\draw[thick,poporange] (0,0)--(H) node[midway,right=1pt]{\small $a$};
\fill[poporange] (H) circle (1.5pt) node[above right=1pt]{\small $H$};
% Marque d'angle droit correcte à H
\draw[poporange] ($(H)!0.15!(O)$) -- ($(H)!0.15!(O)!0.15!90:(S1)$) -- ($(H)!0.15!(S1)$);
\node[popblue] at (90:2.3) {\small $A$};
\node[popblue] at (30:2.3) {\small $B$};
\node[popblue] at (-30:2.3) {\small $C$};
\node[popblue] at (-90:2.3) {\small $D$};
\node[popblue] at (-150:2.3) {\small $E$};
\node[popblue] at (150:2.3) {\small $F$};
\end{tikzpicture}
\end{center}
\legende{Hexagone régulier $ABCDEF$ inscrit dans son cercle circonscrit de centre $O$ : le rayon $R = OA$ et l'apothème $a = OH$, où $H$ est le milieu du côté $[AB]$.}
\end{popfigure}

\courssub{Tout polygone régulier est inscriptible dans un cercle}

Voici la propriété fondamentale qui explique pourquoi le compas est l'outil roi pour construire les polygones réguliers.

\begin{propriete}[Inscriptibilité (admise)]
Tout polygone régulier est \cle{inscriptible} dans un cercle : il existe un cercle, appelé \cle{cercle circonscrit}, qui passe par tous ses sommets. Le centre $O$ de ce cercle est le point d'intersection des \textbf{médiatrices des côtés} du polygone.

Réciproquement : si l'on partage un cercle en $n$ arcs de même longueur et que l'on joint les points de partage consécutifs, on obtient un polygone régulier à $n$ côtés.
\end{propriete}

Cette propriété est admise dans le cas général, mais on peut la comprendre sur deux cas simples.

\begin{experience}[Justification sur le triangle équilatéral et le carré]
\textbf{Cas du triangle équilatéral $ABC$.} Soit $O$ le point d'intersection des médiatrices de $[AB]$ et de $[BC]$.
\begin{enumerate}
\item $O$ est sur la médiatrice de $[AB]$, donc $OA = OB$.
\item $O$ est sur la médiatrice de $[BC]$, donc $OB = OC$.
\item Par conséquent $OA = OB = OC$ : le cercle de centre $O$ et de rayon $OA$ passe par les trois sommets $A$, $B$ et $C$.
\end{enumerate}
\textbf{Cas du carré $ABCD$.} Les diagonales $[AC]$ et $[BD]$ se coupent en leur milieu $O$, et l'on sait que $OA = OB = OC = OD$ (les diagonales d'un carré sont égales et se coupent en leur milieu). Le cercle de centre $O$ passant par $A$ passe donc aussi par $B$, $C$ et $D$. De plus, les diagonales du carré sont les médiatrices les unes des autres, et les médiatrices des côtés passent aussi par $O$.
\end{experience}

\begin{methode}[Construire un polygone régulier à n côtés]
Pour construire un polygone régulier à $n$ côtés avec la règle et le compas :
\begin{enumerate}
\item Tracer un cercle de centre $O$ et de rayon $R$ choisi.
\item Calculer l'angle au centre : $\dfrac{360^\circ}{n}$. Par exemple, pour un hexagone : $\dfrac{360^\circ}{6} = 60^\circ$.
\item Placer un premier sommet $A$ sur le cercle, puis reporter l'angle au centre (ou la corde correspondante) pour obtenir les sommets suivants.
\item Relier les sommets consécutifs à la règle.
\end{enumerate}
Pour l'hexagone, il y a une astuce célèbre : la longueur du côté est \textbf{égale au rayon} $R$ ; il suffit donc de reporter l'ouverture du compas six fois sur le cercle !
\end{methode}

\courssub{Symétries d'un polygone régulier}

\textbf{Intuition.} Plie une feuille découpée en forme de carré le long d'une diagonale : les deux moitiés se superposent exactement. Plie-la le long de la médiatrice d'un côté : même résultat. Le carré possède quatre « plis magiques » : ce sont ses axes de symétrie. Cette richesse est commune à tous les polygones réguliers.

\begin{propriete}[Axes de symétrie]
Un polygone régulier à $n$ côtés possède exactement $n$ axes de symétrie, qui passent tous par son centre $O$ :
\begin{itemize}
\item si $n$ est \textbf{impair} : chaque axe passe par un sommet et par le milieu du côté opposé ;
\item si $n$ est \textbf{pair} : il y a $\dfrac{n}{2}$ axes passant par deux sommets opposés, et $\dfrac{n}{2}$ axes passant par les milieux de deux côtés opposés.
\end{itemize}
\end{propriete}

\begin{exemplebox}[Comptons les axes de symétrie]
\begin{itemize}
\item Triangle équilatéral ($n = 3$, impair) : 3 axes, chacun passant par un sommet et le milieu du côté opposé (ce sont aussi les médianes et les hauteurs).
\item Carré ($n = 4$, pair) : 4 axes — les 2 diagonales et les 2 médiatrices des côtés.
\item Pentagone régulier ($n = 5$, impair) : 5 axes.
\item Hexagone régulier ($n = 6$, pair) : 6 axes — 3 « grandes diagonales » et 3 droites joignant les milieux de côtés opposés.
\end{itemize}
\end{exemplebox}

\begin{popfigure}
\begin{center}
\begin{tikzpicture}[scale=1.1]
\foreach \i in {0,1,2,3,4}{
  \coordinate (P\i) at (90-72*\i:2);
  \fill[popblue] (P\i) circle (1.6pt);
}
\draw[thick,popblue] (P0)--(P1)--(P2)--(P3)--(P4)--cycle;
\fill[popink] (0,0) circle (1.8pt) node[below right=1pt]{\small $O$};
\foreach \i in {0,1,2,3,4}{
  \draw[dashed,poppurple] (90-72*\i:2.6)--(90-72*\i+180:1.35);
}
\node[popblue] at (90:2.25) {\small $A$};
\node[popblue] at (18:2.25) {\small $B$};
\node[popblue] at (-54:2.25) {\small $C$};
\node[popblue] at (-126:2.25) {\small $D$};
\node[popblue] at (162:2.25) {\small $E$};
\end{tikzpicture}
\end{center}
\legende{Pentagone régulier $ABCDE$ et ses 5 axes de symétrie (en pointillés) : chacun passe par un sommet et par le milieu du côté opposé.}
\end{popfigure}

\begin{propriete}[Invariance par rotation]
Un polygone régulier à $n$ côtés et de centre $O$ est \textbf{invariant} par la rotation de centre $O$ et d'angle $\dfrac{360^\circ}{n}$ : cette rotation transforme le polygone en lui-même, chaque sommet étant envoyé sur le sommet suivant.
\end{propriete}

\begin{exemplebox}[Rotation de l'hexagone régulier]
Pour l'hexagone régulier $ABCDEF$ de centre $O$, l'angle de la rotation est $\dfrac{360^\circ}{6} = 60^\circ$. La rotation de centre $O$ et d'angle $60^\circ$ envoie $A$ sur $B$, $B$ sur $C$, \dots, $F$ sur $A$ : l'hexagone se superpose exactement à lui-même. C'est d'ailleurs cette propriété qui permet de le construire au compas en reportant six fois la même ouverture.
\end{exemplebox}

\begin{savaistu}[Le génie des abeilles]
Les cases des ruches d'abeilles sont des \textbf{hexagones réguliers}. Pourquoi ce choix ? Parmi tous les polygones réguliers, seuls le triangle équilatéral, le carré et l'hexagone régulier permettent de paver le plan sans laisser de trous. Et parmi ces trois formes, l'hexagone est celle qui, pour une aire donnée, demande le \textbf{moins de cire} pour les parois : c'est la forme la plus économique ! Les abeilles ont donc « résolu » un problème d'optimisation que les mathématiciens n'ont démontré rigoureusement qu'en 1999 (théorème du nid d'abeille, dû à Thomas Hales). C'est aussi pour cette raison que Maman Aïcha peut paver sa cour avec des pavés hexagonaux, mais pas avec des pentagones réguliers : nous y reviendrons dans la suite de la leçon.
\end{savaistu}

\begin{exercice}[Pour vérifier tes acquis]
\begin{enumerate}
\item Un polygone régulier possède 8 côtés. Combien a-t-il de sommets ? d'axes de symétrie ? Quel est l'angle de la rotation qui le laisse invariant ?
\item Un quadrilatère a quatre angles égaux à $90^\circ$ et deux côtés de $7$ cm et deux côtés de $4$ cm. Est-il régulier ? Justifie.
\item Sur un cercle de centre $O$ et de rayon $5$ cm, place six points également espacés. Quelle est la nature du polygone obtenu en les reliant ? Quelle est la longueur de ses côtés ?
\end{enumerate}
\end{exercice}

\begin{corrige}[Exercice 1]
\begin{enumerate}
\item Un octogone régulier a 8 sommets et 8 axes de symétrie (4 passant par des sommets opposés, 4 par les milieux de côtés opposés). L'angle de rotation est $\dfrac{360^\circ}{8} = 45^\circ$.
\item Ce quadrilatère est un rectangle non carré : ses angles sont égaux mais pas ses côtés. Il n'est \textbf{pas régulier}.
\item Les six points partagent le cercle en 6 arcs égaux : le polygone obtenu est un \textbf{hexagone régulier}. Son côté est égal au rayon, soit $5$ cm.
\end{enumerate}
\end{corrige}

\begin{aretenir}[L'essentiel de la section]
\begin{itemize}
\item Un \textbf{polygone régulier} a tous ses côtés de même longueur \textbf{et} tous ses angles de même mesure : les deux conditions sont nécessaires (pense au losange et au rectangle !).
\item Tout polygone régulier est \textbf{inscriptible dans un cercle} dont le centre $O$ est l'intersection des médiatrices des côtés ; le \textbf{rayon} $R$ relie $O$ à un sommet, l'\textbf{apothème} $a$ relie $O$ au milieu d'un côté.
\item Un polygone régulier à $n$ côtés possède $n$ \textbf{axes de symétrie} et est invariant par la rotation de centre $O$ et d'angle $\dfrac{360^\circ}{n}$.
\end{itemize}
\end{aretenir}

\coursec{Angles dans un polygone régulier}

Dans la section précédente, nous avons défini les polygones réguliers et découvert leurs premières propriétés : tous leurs côtés ont la même longueur, tous leurs angles ont la même mesure, et chacun d'eux est inscrit dans un cercle dont le centre est appelé \cle{centre} du polygone. Mais une question reste en suspens : \emph{combien mesurent ces angles ?} C'est toute l'ambition de cette section. Tu vas voir qu'une seule donnée — le nombre $n$ de côtés — suffit à calculer tous les angles d'un polygone régulier. C'est ce qui rend ces figures si précieuses pour les artisans, les maçons et les décorateurs : avec une formule et un rapporteur, on construit n'importe quel motif régulier.

\courssub{L'angle au centre : partager le tour complet}

\textbf{Intuition.} Imagine une roue de bicyclette à $n$ rayons régulièrement espacés, ou un gâteau découpé en $n$ parts identiques. Le tour complet de la roue (ou du gâteau) représente un angle de $360°$. Si le partage est régulier, chaque part intercepte exactement le même angle au centre.

\begin{definition}[Angle au centre d'un polygone régulier]
Soit un polygone régulier à $n$ côtés, de centre $O$, et soient $A$ et $B$ deux sommets \textbf{consécutifs}. L'\cle{angle au centre} du polygone est l'angle $\widehat{AOB}$.
\end{definition}

\begin{propriete}[Mesure de l'angle au centre]
Dans un polygone régulier à $n$ côtés, l'angle au centre mesure :
\[ \frac{360°}{n} \]
\end{propriete}

\begin{experience}[Pourquoi $360°/n$ ? Justification pas à pas]
Considérons un polygone régulier à $n$ côtés inscrit dans un cercle de centre $O$, de sommets $A_1, A_2, \ldots, A_n$.
\begin{enumerate}
\item Les triangles $OA_1A_2$, $OA_2A_3$, \ldots, $OA_nA_1$ sont tous \textbf{superposables} : ils ont deux côtés égaux au rayon du cercle et un troisième côté égal au côté du polygone (tous les côtés d'un polygone régulier sont égaux).
\item Les $n$ angles au centre $\widehat{A_1OA_2}, \widehat{A_2OA_3}, \ldots$ sont donc tous \textbf{égaux}.
\item Ces $n$ angles, mis bout à bout autour du point $O$, forment un \textbf{tour complet}, c'est-à-dire un angle plein de $360°$.
\item Chacun mesure donc $\dfrac{360°}{n}$. \hfill $\square$
\end{enumerate}
\end{experience}

\begin{exemplebox}[Premiers calculs d'angles au centre]
\begin{itemize}
\item Triangle équilatéral ($n = 3$) : $\dfrac{360°}{3} = 120°$.
\item Carré ($n = 4$) : $\dfrac{360°}{4} = 90°$.
\item Hexagone régulier ($n = 6$) : $\dfrac{360°}{6} = 60°$.
\item Octogone régulier ($n = 8$) : $\dfrac{360°}{8} = 45°$.
\end{itemize}
\end{exemplebox}

\begin{savaistu}[L'angle au centre, outil de construction]
C'est grâce à l'angle au centre qu'on construit un polygone régulier au rapporteur : on trace un cercle, on reporte $n$ fois l'angle $\dfrac{360°}{n}$ depuis le centre, et on relie les points obtenus. Les tailleurs de pierre et les fabricants de tables octogonales procèdent exactement ainsi !
\end{savaistu}

\courssub{Somme des angles intérieurs d'un polygone convexe}

\textbf{Intuition.} Tu sais depuis la classe de 5ème que la somme des angles d'un triangle vaut $180°$. L'idée géniale consiste à découper n'importe quel polygone en triangles : la somme de ses angles s'obtiendra alors en comptant les triangles.

\begin{propriete}[Somme des angles d'un polygone à $n$ côtés — démonstration exigible]
La somme des angles intérieurs d'un polygone convexe à $n$ côtés est :
\[ (n - 2) \times 180° \]
\end{propriete}

\begin{experience}[Démonstration guidée : le découpage en triangles]
Soit un polygone convexe à $n$ côtés et $A$ l'un de ses sommets.
\begin{enumerate}
\item Depuis le sommet $A$, traçons toutes les diagonales issues de $A$. Elles joignent $A$ aux sommets qui ne lui sont ni confondus ni adjacents : il y en a $n - 3$.
\item Ces diagonales découpent le polygone en triangles. Comptons-les : chaque côté du polygone, sauf les deux côtés adjacents à $A$, sert de base à un triangle. Il y a donc exactement $n - 2$ triangles.
\item La somme des angles du polygone est égale à la somme des angles de tous ces triangles (les angles des triangles en $A$ se recomposent en l'angle du polygone en $A$).
\item Chaque triangle contribue pour $180°$. La somme totale vaut donc $(n-2) \times 180°$. \hfill $\square$
\end{enumerate}
\textbf{Vérification :} pour le quadrilatère ($n = 4$) : $(4-2)\times 180° = 360°$, résultat bien connu. Pour le pentagone ($n = 5$) : $3 \times 180° = 540°$.
\end{experience}

\begin{attention}[Convexité et comptage]
La formule $(n-2)\times 180°$ est valable pour les polygones \textbf{convexes} (ceux dont tous les angles intérieurs sont inférieurs à $180°$), ce qui est toujours le cas des polygones réguliers. Autre piège : le nombre de triangles est $n-2$, pas $n$ ! Retiens : « deux côtés de moins que de sommets ».
\end{attention}

\courssub{L'angle intérieur d'un polygone régulier}

Dans un polygone \textbf{régulier}, tous les angles intérieurs sont égaux. Il suffit donc de partager la somme en $n$ parts égales.

\begin{propriete}[Angle intérieur d'un polygone régulier]
L'angle intérieur d'un polygone régulier à $n$ côtés mesure :
\[ \frac{(n-2)\times 180°}{n} \qquad \text{ou, de façon équivalente,} \qquad 180° - \frac{360°}{n} \]
\end{propriete}

\begin{experience}[Pourquoi ces deux formules sont-elles égales ?]
Partons de la première et transformons-la :
\begin{align*}
\frac{(n-2)\times 180°}{n} &= \frac{180° \times n - 360°}{n} = 180° - \frac{360°}{n}.
\end{align*}
\textbf{Interprétation géométrique :} dans le triangle $OAB$ formé par le centre et deux sommets consécutifs, les angles à la base valent chacun $\dfrac{180° - \frac{360°}{n}}{2}$ (triangle isocèle en $O$). Or l'angle intérieur du polygone est formé de \emph{deux} de ces angles à la base, donc il vaut $180° - \dfrac{360°}{n}$. On retrouve la formule !
\end{experience}

\begin{propriete}[Angle au centre et angle intérieur sont supplémentaires]
Dans un polygone régulier, l'angle au centre et l'angle intérieur sont \cle{supplémentaires} :
\[ \text{angle au centre} + \text{angle intérieur} = \frac{360°}{n} + \left(180° - \frac{360°}{n}\right) = 180° \]
\end{propriete}

\begin{methode}[Calculer les angles d'un polygone régulier à $n$ côtés]
\begin{enumerate}
\item Identifier le nombre $n$ de côtés.
\item Calculer l'\textbf{angle au centre} : $\dfrac{360°}{n}$.
\item En déduire l'\textbf{angle intérieur} par supplémentarité : $180° - \dfrac{360°}{n}$.
\item Vérifier avec la somme : $n \times \text{angle intérieur}$ doit être égal à $(n-2)\times 180°$.
\end{enumerate}
\end{methode}

\begin{exemplebox}[Exemple chiffré détaillé : l'octogone régulier]
Pour l'octogone régulier, $n = 8$.
\begin{itemize}
\item Angle au centre : $\dfrac{360°}{8} = 45°$.
\item Angle intérieur : $180° - 45° = 135°$.
\item Vérification : $8 \times 135° = 1080°$ et $(8-2)\times 180° = 6 \times 180° = 1080°$. \checkmark
\end{itemize}
\end{exemplebox}

\begin{popfigure}
\begin{center}
\begin{tikzpicture}[scale=2.1]
\foreach \i in {0,...,7} {
  \coordinate (S\i) at ({22.5+45*\i}:1.4);
}
\draw[thick,popblue] (S0) \foreach \i in {1,...,7} { -- (S\i) } -- cycle;
\coordinate (O) at (0,0);
\fill[popdark] (O) circle (0.02) node[below left=1pt] {\small $O$};
\draw[popmuted] (O) -- (S0);
\draw[popmuted] (O) -- (S1);
\draw[popmuted] (O) -- (S2);
\draw[poporange,thick] (0.45,0) arc (0:45:0.45);
\node[poporange] at ({22.5}:0.62) {\small $45°$};
\draw[popgreen,thick] (S1) ++(-67.5:0.4) arc (-67.5:67.5:0.4);
\node[popgreen] at ($(S1)+(180:0.62)$) {\small $135°$};
\foreach \i in {0,1,2} {
  \fill[popblue] (S\i) circle (0.025);
}
\node[popblue] at ({22.5}:1.58) {\small $A$};
\node[popblue] at ({67.5}:1.58) {\small $B$};
\node[popblue] at ({112.5}:1.55) {\small $C$};
\end{tikzpicture}
\end{center}
\legende{Octogone régulier : angle au centre $\widehat{AOB} = 45°$ (orange) et angle intérieur $\widehat{ABC} = 135°$ (vert). Ils sont supplémentaires : $45° + 135° = 180°$.}
\end{popfigure}

\begin{attention}[Ne pas confondre angle au centre et angle intérieur]
C'est l'erreur la plus fréquente dans les exercices de construction !
\begin{itemize}
\item Pour \textbf{placer les sommets sur le cercle} au rapporteur, on utilise l'\textbf{angle au centre} $\dfrac{360°}{n}$, dont le sommet est le centre $O$.
\item L'\textbf{angle intérieur} $180° - \dfrac{360°}{n}$ a son sommet \emph{sur un sommet du polygone} ; il sert à vérifier la figure ou à construire le polygone sans cercle (règle + rapporteur).
\end{itemize}
Exemple du piège : pour l'hexagone, l'angle au centre vaut $60°$ et l'angle intérieur $120°$. Reporter $120°$ au centre donnerait... un triangle !
\end{attention}

\courssub{Tableau récapitulatif et application au pavage}

\begin{aretenir}[Angles des polygones réguliers particuliers]
\begin{center}
\begin{tabularx}{\linewidth}{|l|c|c|c|X|}
\hline
\rowcolor{popblueL} \textbf{Polygone} & $n$ & \textbf{Angle au centre} $\frac{360°}{n}$ & \textbf{Angle intérieur} $180° - \frac{360°}{n}$ & \textbf{Somme des angles} \\
\hline
Triangle équilatéral & 3 & $120°$ & $60°$ & $180°$ \\
\hline
Carré & 4 & $90°$ & $90°$ & $360°$ \\
\hline
Hexagone régulier & 6 & $60°$ & $120°$ & $720°$ \\
\hline
Octogone régulier & 8 & $45°$ & $135°$ & $1080°$ \\
\hline
\end{tabularx}
\end{center}
\end{aretenir}

\textbf{Le problème du pavage.} Un carreleur de Douala veut recouvrir une cour avec des dalles toutes identiques, en forme de polygone régulier, \emph{sans trou ni recouvrement}. Autour de chaque sommet du pavage, plusieurs dalles se rencontrent : la somme de leurs angles doit faire exactement $360°$. Il faut donc que l'angle intérieur \textbf{divise} $360°$.

\begin{experience}[Démonstration guidée : quels polygones réguliers pavent le plan seuls ?]
Notons $a = 180° - \dfrac{360°}{n}$ l'angle intérieur. S'il y a $k$ dalles autour de chaque sommet, alors $k \times a = 360°$, donc $k = \dfrac{360°}{a}$ doit être un \textbf{nombre entier} (supérieur ou égal à 3, car il faut au moins 3 dalles autour d'un point).
\begin{enumerate}
\item $n = 3$ : $a = 60°$, $k = \dfrac{360°}{60°} = 6$ : entier. Le triangle équilatéral pave. \checkmark
\item $n = 4$ : $a = 90°$, $k = 4$ : entier. Le carré pave. \checkmark
\item $n = 5$ : $a = 108°$, $k = \dfrac{360°}{108°} = 3{,}33\ldots$ : \textbf{pas entier}. Le pentagone ne pave pas : il reste des trous !
\item $n = 6$ : $a = 120°$, $k = 3$ : entier. L'hexagone pave. \checkmark
\item $n \geq 7$ : $a$ est strictement compris entre $120°$ et $180°$, donc $k$ est strictement compris entre 2 et 3 : jamais entier. Aucun ne pave.
\end{enumerate}
\textbf{Conclusion :} seuls le triangle équilatéral, le carré et l'hexagone régulier pavent le plan à eux seuls. \hfill $\square$
\end{experience}

\begin{popfigure}
\begin{center}
\begin{tikzpicture}[scale=0.62]
% Pavage hexagonal
\begin{scope}
\foreach \col in {0,1,2} {
  \foreach \row in {0,1} {
    \begin{scope}[shift={({\col*1.5},{\row*1.732 + \col*0.866})}]
      \draw[thick,popblue,fill=popblueL] (0:0.58) \foreach \a in {60,120,180,240,300} { -- (\a:0.58) } -- cycle;
    \end{scope}
  }
}
\node at (2.2,-1.1) {\small Pavage hexagonal : sans trou};
\end{scope}
% Tentative pentagone
\begin{scope}[shift={(7.2,0)}]
\draw[thick,poporange,fill=poporangeL] (90:0.62) \foreach \a in {162,234,306,378} { -- (\a:0.62) } -- cycle;
\begin{scope}[shift={(18:1.19)}]
  \draw[thick,poporange,fill=poporangeL] (162:0.62) \foreach \a in {234,306,378,450} { -- (\a:0.62) } -- cycle;
\end{scope}
\begin{scope}[shift={(306:1.19)}]
  \draw[thick,poporange,fill=poporangeL] (18:0.62) \foreach \a in {90,162,234,306} { -- (\a:0.62) } -- cycle;
\end{scope}
\draw[thick,popgreen,fill=popgreenL] (0.55,-0.18) -- (0.99,0.14) -- (0.94,-0.62) -- cycle;
\node[popgreen] at (1.6,-0.9) {\small trou !};
\node at (0.9,-1.5) {\small Pentagones : des trous apparaissent};
\end{scope}
\end{tikzpicture}
\end{center}
\legende{À gauche, les hexagones s'assemblent parfaitement ($3 \times 120° = 360°$) : c'est le nid d'abeille. À droite, trois pentagones laissent un trou ($3 \times 108° = 324° < 360°$).}
\end{popfigure}

\begin{savaistu}[Le génie des abeilles]
Les abeilles construisent leurs rayons en hexagones réguliers. Ce n'est pas un hasard : parmi les trois polygones qui pavent le plan, l'hexagone est celui qui, pour une même aire de cellule, demande le \emph{moins de cire} (le plus petit périmètre). Les abeilles optimisent leur matériau mieux que bien des ingénieurs ! Les pavages hexagonaux se retrouvent aussi dans les dalles de trottoir et certains tissus traditionnels.
\end{savaistu}

\begin{exoresolu}[Les dalles du marché]
Un artisan veut fabriquer des dalles en forme de polygone régulier à 12 côtés (dodécagone) pour paver une terrasse à Yaoundé, en utilisant un seul type de dalle.
\begin{enumerate}
\item Calcule l'angle au centre et l'angle intérieur d'un tel polygone.
\item Ce pavage est-il possible avec ce seul type de dalle ? Justifie.
\end{enumerate}
\tcblower
\textbf{Solution.}
\begin{enumerate}
\item Angle au centre : $\dfrac{360°}{12} = 30°$. Angle intérieur : $180° - 30° = 150°$.
\item Autour d'un sommet du pavage, il faudrait $k$ dalles avec $k \times 150° = 360°$, soit $k = \dfrac{360°}{150°} = 2{,}4$. Ce n'est \textbf{pas un nombre entier} : deux dalles laissent un trou de $360° - 2 \times 150° = 60°$, trois dalles se chevauchent de $3 \times 150° - 360° = 90°$. Le pavage est donc \textbf{impossible} avec des dodécagones seuls. (En revanche, en mélangeant dodécagones et triangles équilatéraux, on obtient un beau pavage mixte !)
\end{enumerate}
\end{exoresolu}

\begin{exercice}[À toi de jouer]
\begin{enumerate}
\item Un polygone régulier a un angle au centre de $40°$. Combien a-t-il de côtés ? Quelle est la mesure de son angle intérieur ?
\item Un polygone régulier a un angle intérieur de $156°$. Détermine son nombre de côtés.
\item Peut-on paver une salle de classe uniquement avec des dalles en forme de ce polygone ? Justifie ta réponse.
\end{enumerate}
\end{exercice}

\begin{corrige}[Exercice : À toi de jouer]
\begin{enumerate}
\item $n = \dfrac{360°}{40°} = 9$ côtés (ennéagone). Angle intérieur : $180° - 40° = 140°$.
\item L'angle au centre vaut $180° - 156° = 24°$, donc $n = \dfrac{360°}{24°} = 15$ côtés.
\item Il faudrait $k = \dfrac{360°}{156°} \approx 2{,}3$ dalles par sommet : ce n'est pas entier, donc le pavage est impossible avec ce seul polygone.
\end{enumerate}
\end{corrige}

\begin{aretenir}[L'essentiel de la section]
\begin{itemize}
\item Angle au centre d'un polygone régulier à $n$ côtés : $\dfrac{360°}{n}$.
\item Somme des angles d'un polygone convexe à $n$ côtés : $(n-2)\times 180°$ (découpage en $n-2$ triangles).
\item Angle intérieur : $180° - \dfrac{360°}{n}$ ; il est \textbf{supplémentaire} de l'angle au centre.
\item Un polygone régulier pave le plan seul si et seulement si son angle intérieur divise $360°$ : seuls le triangle équilatéral, le carré et l'hexagone conviennent.
\end{itemize}
\end{aretenir}

\coursec{Triangle équilatéral et carré}

Dans la section précédente, nous avons appris à calculer les angles d'un polygone régulier : l'angle au centre vaut $\dfrac{360°}{n}$ et l'angle intérieur vaut $\dfrac{(n-2)\times 180°}{n}$, où $n$ est le nombre de côtés. Nous allons maintenant appliquer ces résultats aux deux polygones réguliers les plus simples et les plus utiles de la vie courante : le \cle{triangle équilatéral} ($n=3$) et le \cle{carré} ($n=4$). Tu les rencontres partout : panneaux de signalisation, carreaux de faïence, terrains de sport, toitures des cases. Allons-y !

\courssub{Le triangle équilatéral}

\textbf{Définition et propriétés}

\begin{definition}[Triangle équilatéral]
Un \cle{triangle équilatéral} est un triangle dont les trois côtés ont la même longueur. C'est le polygone régulier à $3$ côtés.
\end{definition}

Appliquons la formule des angles vue à la section précédente. L'angle intérieur d'un polygone régulier à $3$ côtés vaut :
\[ \frac{(3-2)\times 180°}{3} = \frac{180°}{3} = 60°. \]
L'angle au centre vaut $\dfrac{360°}{3} = 120°$.

\begin{propriete}[Propriétés du triangle équilatéral]
Dans un triangle équilatéral $ABC$ :
\begin{itemize}
\item les trois angles sont égaux et mesurent chacun $60°$ ;
\item le triangle possède \cle{trois axes de symétrie} : les trois médiatrices des côtés, qui sont aussi les hauteurs, les médianes et les bissectrices ;
\item ces trois droites se coupent en un même point $O$, centre du cercle circonscrit et du cercle inscrit.
\end{itemize}
\end{propriete}

\begin{attention}[Triangle équilatéral et axes de symétrie]
Ne confonds pas : un triangle équilatéral a \textbf{trois} axes de symétrie, pas un seul. Chacun passe par un sommet et par le milieu du côté opposé. En revanche, un triangle équilatéral n'a \textbf{pas} de centre de symétrie.
\end{attention}

\textbf{Construction à la règle et au compas}

\begin{methode}[Construire un triangle équilatéral de côté donné]
Pour construire un triangle équilatéral $ABC$ connaissant la longueur $a$ du côté :
\begin{enumerate}
\item Trace le segment $[AB]$ de longueur $a$.
\item Trace l'arc de cercle de centre $A$ et de rayon $a$.
\item Trace l'arc de cercle de centre $B$ et de même rayon $a$.
\item Les deux arcs se coupent en $C$. Trace $[AC]$ et $[BC]$.
\end{enumerate}
Par construction, $AC = AB = a$ et $BC = AB = a$ : le triangle est bien équilatéral.
\end{methode}

\begin{popfigure}
\begin{center}
\begin{tikzpicture}[scale=1.6]
\coordinate (A) at (0,0);
\coordinate (B) at (3,0);
\coordinate (C) at (1.5,2.598);
\draw[popmuted,dashed] (A) ++(0:3) arc (0:70:3);
\draw[popmuted,dashed] (B) ++(180:3) arc (180:110:3);
\draw[thick,popblue] (A) -- (B) -- (C) -- cycle;
\fill[popblue] (A) circle (1.2pt) node[below left]{$A$};
\fill[popblue] (B) circle (1.2pt) node[below right]{$B$};
\fill[popblue] (C) circle (1.2pt) node[above]{$C$};
\node[popblue,below] at (1.5,0) {$a$};
\end{tikzpicture}
\end{center}
\legende{Construction du triangle équilatéral : deux arcs de même rayon $a$, centrés en $A$ et $B$.}
\end{popfigure}

\begin{methode}[Construire un triangle équilatéral inscrit dans un cercle]
Pour inscrire un triangle équilatéral dans un cercle de centre $O$ et de rayon $R$ :
\begin{enumerate}
\item Place un premier point $A$ sur le cercle.
\item L'angle au centre vaut $120°$ : avec le rapporteur, construis les angles $\widehat{AOB} = 120°$ et $\widehat{BOC} = 120°$ pour obtenir les points $B$ puis $C$ sur le cercle.
\item Relie $A$, $B$ et $C$.
\end{enumerate}
Astuce de vérification : le côté du triangle équilatéral inscrit vaut $a = R\sqrt{3}$, où $R$ est le rayon du cercle. Tu peux donc contrôler ta construction en mesurant le côté obtenu.
\end{methode}

\textbf{Hauteur et aire du triangle équilatéral}

Voici le grand résultat de cette sous-section. Soit $ABC$ un triangle équilatéral de côté $a$, et $H$ le pied de la hauteur issue de $C$. Comme la hauteur est aussi médiane, $H$ est le milieu de $[AB]$, donc $AH = HB = \dfrac{a}{2}$.

\begin{propriete}[Hauteur et aire du triangle équilatéral]
Dans un triangle équilatéral de côté $a$, la hauteur vaut
\[ h = \frac{a\sqrt{3}}{2} \]
et l'aire vaut
\[ \mathcal{A} = \frac{a^2\sqrt{3}}{4}. \]
\emph{Démonstration exigible au BEPC (théorème de Pythagore).}
\end{propriete}

\begin{experience}[Démonstration : la hauteur par Pythagore]
Le triangle $CBH$ est rectangle en $H$, avec $CB = a$ (hypoténuse) et $BH = \dfrac{a}{2}$. D'après le théorème de Pythagore :
\begin{align*}
CB^2 &= CH^2 + BH^2\\
a^2 &= h^2 + \left(\frac{a}{2}\right)^2\\
h^2 &= a^2 - \frac{a^2}{4} = \frac{4a^2 - a^2}{4} = \frac{3a^2}{4}\\
h &= \sqrt{\frac{3a^2}{4}} = \frac{a\sqrt{3}}{2}.
\end{align*}
L'aire s'en déduit immédiatement :
\[ \mathcal{A} = \frac{\text{base}\times\text{hauteur}}{2} = \frac{a \times \dfrac{a\sqrt{3}}{2}}{2} = \frac{a^2\sqrt{3}}{4}. \]
\end{experience}

\begin{popfigure}
\begin{center}
\begin{tikzpicture}[scale=1.8]
\coordinate (A) at (0,0);
\coordinate (B) at (3,0);
\coordinate (C) at (1.5,2.598);
\coordinate (H) at (1.5,0);
\draw[thick,popblue] (A) -- (B) -- (C) -- cycle;
\draw[thick,poporange] (C) -- (H);
\draw[popdark] (H) ++(0,0.18) -- ++(-0.18,0.18) -- ++(-0.18,0);
\fill[popblue] (A) circle (1pt) node[below left]{$A$};
\fill[popblue] (B) circle (1pt) node[below right]{$B$};
\fill[popblue] (C) circle (1pt) node[above]{$C$};
\fill[poporange] (H) circle (1pt) node[below]{$H$};
\node[popblue,below left] at (0.75,0) {$\frac{a}{2}$};
\node[popblue,below right] at (2.25,0) {$\frac{a}{2}$};
\node[poporange,right] at (1.55,1.3) {$h=\frac{a\sqrt{3}}{2}$};
\draw[popgreen] (B) ++(180:0.7) arc (180:120:0.7);
\node[popgreen] at (2.3,0.28) {$60°$};
\end{tikzpicture}
\end{center}
\legende{Triangle équilatéral de côté $a$ : la hauteur $[CH]$ coupe $[AB]$ en son milieu.}
\end{popfigure}

\begin{attention}[L'erreur classique : oublier le dénominateur]
Beaucoup d'élèves écrivent $h = a\sqrt{3}$ au lieu de $h = \dfrac{a\sqrt{3}}{2}$. C'est faux ! Vérifie toujours la cohérence : la hauteur doit être \textbf{plus petite} que le côté $a$ (un triangle « aplati » n'existe pas). Or $\dfrac{\sqrt{3}}{2} \approx 0,87 < 1$, donc $h \approx 0,87\,a < a$ : c'est cohérent. De même, l'aire est $\dfrac{a^2\sqrt{3}}{4}$, pas $a^2\sqrt{3}$.
\end{attention}

\begin{exoresolu}[Panneau triangulaire]
Un panneau de signalisation a la forme d'un triangle équilatéral de côté $a = 90$ cm. Calcule sa hauteur et son aire (donne la valeur exacte puis une valeur approchée à $0,1$ près).
\tcblower
\textbf{Hauteur :}
\[ h = \frac{a\sqrt{3}}{2} = \frac{90\sqrt{3}}{2} = 45\sqrt{3} \approx 77,9 \text{ cm}. \]
\textbf{Aire :}
\[ \mathcal{A} = \frac{a^2\sqrt{3}}{4} = \frac{90^2\sqrt{3}}{4} = \frac{8100\sqrt{3}}{4} = 2025\sqrt{3} \approx 3507,4 \text{ cm}^2. \]
Le panneau mesure environ $77,9$ cm de haut et sa surface est d'environ $3507,4$ cm$^2$.
\end{exoresolu}

\courssub{Le carré}

\textbf{Définition et propriétés}

\begin{definition}[Carré]
Un \cle{carré} est un quadrilatère dont les quatre côtés ont la même longueur et dont les quatre angles sont droits ($90°$). C'est le polygone régulier à $4$ côtés.
\end{definition}

Vérifions avec la formule des angles : l'angle intérieur vaut $\dfrac{(4-2)\times 180°}{4} = \dfrac{360°}{4} = 90°$, et l'angle au centre vaut $\dfrac{360°}{4} = 90°$. Tout concorde.

\begin{propriete}[Diagonales du carré]
Dans un carré, les deux diagonales :
\begin{itemize}
\item ont la \cle{même longueur} ;
\item sont \cle{perpendiculaires} ;
\item se coupent en leur \cle{milieu} (chacune est la médiatrice de l'autre) ;
\item sont des axes de symétrie du carré (qui en possède quatre au total : les deux diagonales et les deux médiatrices des côtés).
\end{itemize}
Le point d'intersection des diagonales est le \cle{centre de symétrie} du carré.
\end{propriete}

\textbf{Diagonale du carré}

\begin{propriete}[Diagonale d'un carré de côté $a$]
La diagonale $d$ d'un carré de côté $a$ vaut
\[ d = a\sqrt{2}. \]
\emph{Démonstration exigible au BEPC (théorème de Pythagore).}
\end{propriete}

\begin{experience}[Démonstration : la diagonale par Pythagore]
Soit $ABCD$ un carré de côté $a$. La diagonale $[AC]$ partage le carré en deux triangles rectangles. Dans le triangle $ABC$, rectangle en $B$ :
\begin{align*}
AC^2 &= AB^2 + BC^2\\
d^2 &= a^2 + a^2 = 2a^2\\
d &= \sqrt{2a^2} = a\sqrt{2}.
\end{align*}
Comme $\sqrt{2} \approx 1,414$, la diagonale est environ $1,4$ fois plus longue que le côté : c'est cohérent, la diagonale est le « chemin en ligne droite » à travers le carré.
\end{experience}

\begin{exemplebox}[Terrain carré à Bafoussam]
Un terrain carré à Bafoussam a un côté de $20$ m.
\begin{itemize}
\item \textbf{Diagonale :} $d = 20\sqrt{2} \approx 28,3$ m. C'est la longueur du chemin qui traverse le terrain d'un coin à l'autre.
\item \textbf{Aire :} $\mathcal{A} = a^2 = 20^2 = 400$ m$^2$.
\item \textbf{Périmètre :} $P = 4a = 80$ m : c'est la longueur de grillage nécessaire pour clôturer le terrain.
\end{itemize}
\end{exemplebox}

\textbf{Carré inscrit dans un cercle}

\begin{methode}[Construire un carré inscrit dans un cercle]
Pour inscrire un carré dans un cercle de centre $O$ :
\begin{enumerate}
\item Trace un premier diamètre $[AC]$.
\item Trace la perpendiculaire à $(AC)$ passant par $O$ : elle coupe le cercle en $B$ et $D$ (deuxième diamètre).
\item Relie les quatre points $A$, $B$, $C$, $D$ dans l'ordre : $ABCD$ est un carré.
\end{enumerate}
Justification : les diagonales $[AC]$ et $[BD]$ sont des diamètres (donc de même longueur), elles se coupent en leur milieu $O$ et sont perpendiculaires. Un quadrilatère dont les diagonales ont ces trois propriétés est un carré.
\end{methode}

\begin{popfigure}
\begin{center}
\begin{tikzpicture}[scale=1.5]
\coordinate (O) at (0,0);
\draw[popline] (O) circle (2);
\coordinate (A) at (45:2);
\coordinate (B) at (135:2);
\coordinate (C) at (225:2);
\coordinate (D) at (315:2);
\draw[thick,popblue] (A) -- (B) -- (C) -- (D) -- cycle;
\draw[thick,poporange] (A) -- (C);
\draw[thick,popgreen] (B) -- (D);
\draw[popdark] (0.177,0.177) -- (0,0.354) -- (-0.177,0.177);
\fill[popblue] (A) circle (1.2pt) node[above right]{$A$};
\fill[popblue] (B) circle (1.2pt) node[above left]{$B$};
\fill[popblue] (C) circle (1.2pt) node[below left]{$C$};
\fill[popblue] (D) circle (1.2pt) node[below right]{$D$};
\fill[popdark] (O) circle (1.2pt) node[below left=1pt]{$O$};
\node[poporange] at (0.75,-0.55) {$d = 2R$};
\node[popblue] at (0,1.62) {$a$};
\end{tikzpicture}
\end{center}
\legende{Carré inscrit dans un cercle : deux diamètres perpendiculaires donnent les quatre sommets.}
\end{popfigure}

\textbf{Lien entre le côté et le rayon du cercle circonscrit}

Dans le carré inscrit, la diagonale est un diamètre : $d = 2R$. Comme $d = a\sqrt{2}$, on obtient :
\[ 2R = a\sqrt{2} \quad\Longrightarrow\quad R = \frac{a\sqrt{2}}{2}. \]

Pour le triangle équilatéral, on admet le résultat suivant (tu le justifieras en exercice à l'aide de la hauteur) :

\begin{aretenir}[Rayon du cercle circonscrit]
\begin{itemize}
\item Carré de côté $a$ : $R = \dfrac{a\sqrt{2}}{2}$.
\item Triangle équilatéral de côté $a$ : $R = \dfrac{a\sqrt{3}}{3}$ (admis).
\end{itemize}
Remarque : pour le triangle équilatéral, le centre $O$ est situé aux deux tiers de la hauteur à partir du sommet : $R = \dfrac{2}{3}h = \dfrac{2}{3}\times\dfrac{a\sqrt{3}}{2} = \dfrac{a\sqrt{3}}{3}$.
\end{aretenir}

\begin{exoresolu}[Table ronde et nappe carrée]
Mama Fotso veut recouvrir une table ronde de rayon $R = 60$ cm par une nappe carrée inscrite exactement dans le cercle. Quelle est la longueur du côté de la nappe ?
\tcblower
On utilise $R = \dfrac{a\sqrt{2}}{2}$, donc :
\[ a = \frac{2R}{\sqrt{2}} = \frac{2R\sqrt{2}}{\sqrt{2}\times\sqrt{2}} = \frac{2R\sqrt{2}}{2} = R\sqrt{2} = 60\sqrt{2} \approx 84,9 \text{ cm}. \]
La nappe carrée doit avoir un côté d'environ $84,9$ cm.
\end{exoresolu}

\begin{attention}[Bien choisir la bonne formule]
Ne mélange pas les deux formules du rayon : $R = \dfrac{a\sqrt{2}}{2}$ pour le \textbf{carré}, $R = \dfrac{a\sqrt{3}}{3}$ pour le \textbf{triangle équilatéral}. Un moyen de vérifier : dans les deux cas, le côté $a$ doit être plus grand que le rayon $R$ (le cercle passe par les sommets, et le polygone est à l'intérieur du cercle). Or $\dfrac{\sqrt{2}}{2}\approx 0,71$ et $\dfrac{\sqrt{3}}{3}\approx 0,58$, tous deux inférieurs à $1$ : cohérent.
\end{attention}

\begin{savaistu}[Le carré dans la vie camerounaise]
Les carreaux de faïence des maisons, les damiers des jeux de dames très populaires dans les quartiers de Yaoundé et de Douala, et même les plans de nombreux terrains urbains sont carrés. Quand un maçon vérifie qu'une pièce est bien « d'équerre », il mesure les deux diagonales : si elles sont égales, les angles sont droits. C'est exactement la propriété des diagonales du carré (et du rectangle) utilisée sur le chantier !
\end{savaistu}

\begin{exercice}[Pour t'entraîner]
\begin{enumerate}
\item Un triangle équilatéral a un côté de $12$ cm. Calcule sa hauteur, son aire et le rayon de son cercle circonscrit.
\item La diagonale d'un carré mesure $10$ cm. Calcule son côté, puis son aire.
\item Construis un cercle de rayon $4$ cm, puis inscris-y un carré et un triangle équilatéral. Mesure les côtés obtenus et compare avec les formules du cours.
\end{enumerate}
\end{exercice}

\begin{corrige}[Exercice]
\begin{enumerate}
\item Hauteur : $h = \dfrac{12\sqrt{3}}{2} = 6\sqrt{3} \approx 10,4$ cm. Aire : $\mathcal{A} = \dfrac{12^2\sqrt{3}}{4} = \dfrac{144\sqrt{3}}{4} = 36\sqrt{3} \approx 62,4$ cm$^2$. Rayon : $R = \dfrac{12\sqrt{3}}{3} = 4\sqrt{3} \approx 6,9$ cm.
\item $d = a\sqrt{2}$ donc $a = \dfrac{10}{\sqrt{2}} = \dfrac{10\sqrt{2}}{2} = 5\sqrt{2} \approx 7,1$ cm. Aire : $\mathcal{A} = a^2 = (5\sqrt{2})^2 = 25\times 2 = 50$ cm$^2$.
\item Carré : $a = R\sqrt{2} = 4\sqrt{2} \approx 5,7$ cm. Triangle équilatéral : $a = R\sqrt{3} = 4\sqrt{3} \approx 6,9$ cm. Les mesures à la règle doivent être proches de ces valeurs.
\end{enumerate}
\end{corrige}

\begin{aretenir}[L'essentiel de la section]
\begin{itemize}
\item Triangle équilatéral : trois côtés égaux, trois angles de $60°$, trois axes de symétrie.
\item Hauteur : $h = \dfrac{a\sqrt{3}}{2}$ ; aire : $\mathcal{A} = \dfrac{a^2\sqrt{3}}{4}$ (par Pythagore).
\item Carré : quatre côtés égaux, quatre angles droits ; diagonales égales, perpendiculaires et médiatrices l'une de l'autre.
\item Diagonale du carré : $d = a\sqrt{2}$ (par Pythagore).
\item Carré inscrit dans un cercle : tracer deux diamètres perpendiculaires.
\item Rayons du cercle circonscrit : $R = \dfrac{a\sqrt{2}}{2}$ (carré) et $R = \dfrac{a\sqrt{3}}{3}$ (triangle équilatéral).
\end{itemize}
\end{aretenir}

\coursec{Hexagone et octogone réguliers}

Dans la section précédente, tu as découvert les deux polygones réguliers les plus simples : le \cle{triangle équilatéral} (3 côtés) et le \cle{carré} (4 côtés). Tu as vu que chacun possède un angle au centre ($120°$ puis $90°$) et un angle intérieur ($60°$ puis $90°$). Continuons notre voyage : que se passe-t-il avec 6 côtés, puis avec 8 côtés ? Tu vas découvrir que l'\cle{hexagone régulier} cache une propriété étonnante — son côté est exactement égal au rayon de son cercle circonscrit — et que l'\cle{octogone régulier} se construit très élégamment à partir d'un carré. Ces deux figures sont partout autour de toi : les écrous des vélos, les panneaux STOP, les tables de salon, les alvéoles des abeilles.

\courssub{L'hexagone régulier : définition et angles}

\begin{definition}[Hexagone régulier]
Un \cle{hexagone régulier} est un polygone régulier à $6$ côtés : ses six côtés ont la même longueur et ses six angles intérieurs ont la même mesure. Ses six sommets sont situés sur un même cercle, appelé \cle{cercle circonscrit}, dont le centre $O$ est le \cle{centre} de l'hexagone.
\end{definition}

Appliquons les formules vues dans la section 2 avec $n = 6$ :
\begin{itemize}
\item \textbf{Angle au centre} : $\dfrac{360°}{6} = 60°$. Chaque angle $\widehat{AOB}$ formé par deux rayons vers deux sommets consécutifs mesure $60°$.
\item \textbf{Angle intérieur} : $\dfrac{(6-2)\times 180°}{6} = \dfrac{720°}{6} = 120°$.
\end{itemize}

\begin{aretenir}[Angles de l'hexagone régulier]
Dans un hexagone régulier :
\[ \text{angle au centre} = 60° \qquad \text{angle intérieur} = 120° \]
On remarque que $60° + 120° = 180°$ : l'angle au centre et l'angle intérieur sont \emph{supplémentaires}.
\end{aretenir}

\courssub{Propriété fondamentale : le côté est égal au rayon}

Voici la propriété qui rend l'hexagone si spécial — et si facile à construire.

\begin{propriete}[Côté de l'hexagone régulier]
Dans un hexagone régulier inscrit dans un cercle de centre $O$ et de rayon $R$, la longueur du côté est égale au rayon :
\[ a = R. \]
\end{propriete}

\begin{experience}[Démonstration : pourquoi le côté est égal au rayon]
Soit $A$ et $B$ deux sommets consécutifs d'un hexagone régulier inscrit dans un cercle de centre $O$ et de rayon $R$.

\textbf{Étape 1.} Le triangle $OAB$ est \cle{isocèle} en $O$, car $OA = OB = R$ (rayons du même cercle).

\textbf{Étape 2.} L'angle au centre vaut $\widehat{AOB} = \dfrac{360°}{6} = 60°$.

\textbf{Étape 3.} Dans un triangle isocèle, les angles à la base sont égaux :
\[ \widehat{OAB} = \widehat{OBA} = \frac{180° - 60°}{2} = \frac{120°}{2} = 60°. \]

\textbf{Étape 4.} Les trois angles du triangle $OAB$ valent donc $60°$ : le triangle $OAB$ est \cle{équilatéral}.

\textbf{Conclusion.} $AB = OA = OB = R$. Le côté de l'hexagone régulier est égal au rayon du cercle circonscrit. $\blacksquare$
\end{experience}

\begin{popfigure}
\begin{center}
\begin{tikzpicture}[scale=1.1]
\draw[popmuted] (0,0) circle (2);
\foreach \i/\name in {0/A,60/B,120/C,180/D,240/E,300/F}
  \coordinate (\name) at (\i:2);
\foreach \i/\col in {0/popblueL,60/poporangeL,120/popgreenL,180/poppurpleL,240/poppinkL,300/popgoldL}
  \fill[\col,opacity=0.75] (0,0) -- (\i:2) -- (\i+60:2) -- cycle;
\draw[very thick,popdark] (A)--(B)--(C)--(D)--(E)--(F)--cycle;
\foreach \p in {A,B,C,D,E,F} \fill[popdark] (\p) circle (1.5pt);
\fill[popdark] (0,0) circle (1.5pt) node[below left=1pt] {$O$};
\draw[popink,thick] (0,0) -- (A) node[midway,below right] {$R$};
\draw[poporange,thick] (A) -- (B) node[midway,right] {$a$};
\node[popink] at (0.75,0.28) {$60°$};
\node[anchor=west] at (A) {$A$};
\node[anchor=west] at (B) {$B$};
\node[anchor=south] at (C) {$C$};
\node[anchor=east] at (D) {$D$};
\node[anchor=east] at (E) {$E$};
\node[anchor=north] at (F) {$F$};
\end{tikzpicture}
\end{center}
\legende{Hexagone régulier inscrit dans un cercle, décomposé en 6 triangles équilatéraux. Chaque triangle a pour côtés $R$, $R$ et $a$, avec $a = R$.}
\end{popfigure}

\begin{attention}[Ne pas confondre avec les autres polygones]
Cette égalité $a = R$ est \textbf{propre à l'hexagone}. Pour le carré inscrit, le côté vaut $R\sqrt{2}$ ; pour le triangle équilatéral inscrit, le côté vaut $R\sqrt{3}$. Ne généralise pas abusivement : seul l'hexagone vérifie $a = R$.
\end{attention}

\courssub{Construction de l'hexagone régulier}

La propriété $a = R$ donne immédiatement une construction au compas d'une simplicité remarquable : il suffit de \emph{reporter le rayon six fois} sur le cercle.

\begin{methode}[Construire un hexagone régulier inscrit dans un cercle]
\begin{enumerate}
\item Trace un cercle de centre $O$ et de rayon $R$. Place un premier sommet $A$ sur le cercle.
\item Sans changer l'écartement du compas (égal à $R$), pique le compas en $A$ et trace un arc qui coupe le cercle en $B$.
\item Pique le compas en $B$, trace un arc qui coupe le cercle en $C$. Recommence pour obtenir $D$, $E$, $F$.
\item Le sixième report retombe exactement sur $A$ : c'est une vérification de précision.
\item Relie les six points consécutifs : $ABCDEF$ est un hexagone régulier.
\end{enumerate}
\textbf{Justification} : chaque corde tracée mesure $R$, donc chaque triangle formé avec le centre est équilatéral, chaque angle au centre vaut $60°$, et $6 \times 60° = 360°$ : le tour complet est exactement couvert.
\end{methode}

\begin{savaistu}[Pourquoi les écrous sont-ils hexagonaux ?]
Regarde un vélo, une moto ou les roues d'une voiture à Douala : les écrous et les boulons ont presque toujours \textbf{six pans}. Pourquoi ? Six faces offrent une excellente prise à la clé : on peut serrer en tournant par petits angles de $60°$, ce qui est précieux dans les endroits exigus où l'on ne peut pas faire un grand mouvement. De plus, l'hexagone reste facile à fabriquer à partir de barres de métal. Les abeilles, elles aussi, ont « choisi » l'hexagone : les alvéoles de leurs ruches sont des prismes à base hexagonale, car l'hexagone permet de paver le plan sans laisser aucun trou, avec un minimum de cire pour un maximum de volume de miel. Les mathématiciens ont démontré que c'est la forme la plus économique pour paver le plan !
\end{savaistu}

\courssub{Aire de l'hexagone régulier}

Puisque l'hexagone régulier se décompose en \textbf{6 triangles équilatéraux} de côté $a$, son aire se déduit de la formule de l'aire du triangle équilatéral vue dans la section précédente : $\mathcal{A}_{\triangle} = \dfrac{a^2\sqrt{3}}{4}$.

\begin{propriete}[Aire de l'hexagone régulier]
L'aire d'un hexagone régulier de côté $a$ est :
\[ \mathcal{A} = 6 \times \frac{a^2\sqrt{3}}{4} = \frac{3\sqrt{3}}{2}\,a^2. \]
Son périmètre est $\mathcal{P} = 6a$.
\end{propriete}

\begin{exoresolu}[Hexagone de côté 4 cm]
Un hexagone régulier a pour côté $a = 4$ cm.
\begin{enumerate}
\item Quel est le rayon de son cercle circonscrit ?
\item Calcule son périmètre.
\item Calcule son aire exacte, puis donne une valeur approchée au dixième de cm².
\end{enumerate}
\tcblower
\textbf{1. Rayon.} Par la propriété fondamentale, $R = a = 4$ cm.

\textbf{2. Périmètre.} $\mathcal{P} = 6a = 6 \times 4 = 24$ cm.

\textbf{3. Aire.}
\[ \mathcal{A} = \frac{3\sqrt{3}}{2}\,a^2 = \frac{3\sqrt{3}}{2} \times 16 = 24\sqrt{3} \text{ cm}^2. \]
Valeur approchée : $\mathcal{A} \approx 24 \times 1,732 \approx 41,6$ cm².

\textbf{Vérification par les triangles} : l'aire d'un triangle équilatéral de côté 4 est $\dfrac{16\sqrt{3}}{4} = 4\sqrt{3}$ cm², et $6 \times 4\sqrt{3} = 24\sqrt{3}$ cm². Cohérent !
\end{exoresolu}

\begin{attention}[Pièges fréquents avec l'hexagone]
\begin{itemize}
\item Oublier de multiplier par 6 : l'aire n'est pas $\dfrac{a^2\sqrt{3}}{4}$ (c'est l'aire d'\emph{un seul} triangle), mais $6 \times \dfrac{a^2\sqrt{3}}{4}$.
\item Confondre le rayon du cercle circonscrit ($R = a$) avec l'\cle{apothème} (distance du centre au milieu d'un côté), qui vaut $\dfrac{a\sqrt{3}}{2}$ — tu verras cela dans la section 5.
\item Croire que l'angle intérieur vaut $60°$ : c'est l'angle \emph{au centre} qui vaut $60°$ ; l'angle intérieur vaut $120°$.
\end{itemize}
\end{attention}

\courssub{L'octogone régulier : définition et angles}

\begin{definition}[Octogone régulier]
Un \cle{octogone régulier} est un polygone régulier à $8$ côtés : huit côtés de même longueur, huit angles intérieurs de même mesure, et huit sommets sur un même cercle circonscrit.
\end{definition}

Avec $n = 8$ :
\begin{itemize}
\item \textbf{Angle au centre} : $\dfrac{360°}{8} = 45°$.
\item \textbf{Angle intérieur} : $\dfrac{(8-2)\times 180°}{8} = \dfrac{1080°}{8} = 135°$.
\end{itemize}

\begin{aretenir}[Angles de l'octogone régulier]
\[ \text{angle au centre} = 45° \qquad \text{angle intérieur} = 135° \]
Là encore, $45° + 135° = 180°$ : angle au centre et angle intérieur sont supplémentaires. Remarque aussi que $135° = 90° + 45°$ : l'angle intérieur de l'octogone, c'est un angle droit plus un demi-angle droit.
\end{aretenir}

\courssub{Construction de l'octogone régulier inscrit}

Contrairement à l'hexagone, le côté de l'octogone n'est pas égal au rayon : on ne peut pas le construire par simple report du rayon. L'astuce consiste à \textbf{partir d'un carré inscrit} et à \textbf{doubler le nombre de côtés} grâce aux bissectrices.

\begin{methode}[Construire un octogone régulier à partir d'un carré inscrit]
\begin{enumerate}
\item Trace un cercle de centre $O$. Trace deux diamètres \emph{perpendiculaires} $[AC]$ et $[BD]$ : les quatre extrémités $A$, $B$, $C$, $D$ forment un \textbf{carré inscrit} (en pointillés sur la figure).
\item Trace la \textbf{bissectrice} de l'angle au centre $\widehat{AOB}$ (ou, ce qui revient au même, la \textbf{médiatrice} du côté $[AB]$) : elle coupe le cercle en un nouveau point, milieu de l'arc $\wideparen{AB}$.
\item Fais de même pour les trois autres angles au centre du carré : tu obtiens quatre nouveaux points sur le cercle.
\item Relie les huit points dans l'ordre : tu obtiens un octogone régulier.
\end{enumerate}
\textbf{Justification} : la bissectrice partage chaque angle de $90°$ en deux angles de $45°$. On obtient donc $8$ angles au centre de $45°$, et $8 \times 45° = 360°$ : les huit arcs sont égaux, donc les huit cordes (côtés) sont égales : l'octogone est régulier.
\end{methode}

\begin{popfigure}
\begin{center}
\begin{tikzpicture}[scale=1.05]
\draw[popmuted] (0,0) circle (2);
\coordinate (A) at (90:2); \coordinate (B) at (180:2);
\coordinate (C) at (270:2); \coordinate (D) at (0:2);
\draw[dashed,popblue,thick] (A)--(B)--(C)--(D)--cycle;
\foreach \i in {45,135,225,315}
  \draw[poporange,dashed] (0,0) -- (\i:2.15);
\foreach \i/\name in {90/P1,135/P2,180/P3,225/P4,270/P5,315/P6,0/P7,45/P8}
  \coordinate (\name) at (\i:2);
\draw[very thick,popink] (P1)--(P8)--(P7)--(P6)--(P5)--(P4)--(P3)--(P2)--cycle;
\foreach \p in {P1,P2,P3,P4,P5,P6,P7,P8} \fill[popdark] (\p) circle (1.5pt);
\fill[popdark] (0,0) circle (1.5pt) node[below left=1pt] {$O$};
\node[poporange] at (67.5:1.15) {$45°$};
\node[popblue] at (135:2.35) {\footnotesize carré inscrit};
\node[popink,anchor=west] at (22.5:2.05) {\footnotesize octogone};
\end{tikzpicture}
\end{center}
\legende{Construction de l'octogone régulier : le carré inscrit (pointillés bleus) et les bissectrices des angles au centre (pointillés oranges) fournissent les 8 sommets de l'octogone (trait gras).}
\end{popfigure}

\begin{attention}[Bissectrice ou médiatrice : les deux marchent]
La médiatrice d'un côté du carré inscrit passe par le centre $O$ (car le triangle $OAB$ est isocèle en $O$) et coupe le cercle au milieu de l'arc : c'est exactement la même droite que la bissectrice de l'angle au centre. Utilise celle qui te arrange avec ton matériel, mais \textbf{justifie} ton choix dans la rédaction.
\end{attention}

\courssub{L'octogone dans la vie courante}

\begin{popfigure}
\begin{center}
\begin{tikzpicture}[scale=1.0]
\foreach \i in {0,45,...,315}
  \coordinate (S\i) at (\i+22.5:1.7);
\fill[popink] (S0)--(S45)--(S90)--(S135)--(S180)--(S225)--(S270)--(S315)--cycle;
\draw[white,very thick] (22.5:1.35)--(67.5:1.35)--(112.5:1.35)--(157.5:1.35)--(202.5:1.35)--(247.5:1.35)--(292.5:1.35)--(337.5:1.35)--cycle;
\node[white] at (0,0) {\Large\bfseries STOP};
\draw[<->,popdark] (22.5:1.95) -- (67.5:1.95) node[midway,above,sloped] {\footnotesize $a$};
\draw[popdark,fill=popdark] (0,-1.7) rectangle (0.12,-3.2);
\node[popmuted] at (0,-3.5) {\footnotesize poteau};
\end{tikzpicture}
\end{center}
\legende{Le panneau STOP est un octogone régulier : sa forme à 8 côtés est reconnaissable de loin, même de dos ou dans le brouillard.}
\end{popfigure}

L'octogone régulier est l'un des polygones les plus présents dans notre environnement :
\begin{itemize}
\item Le \textbf{panneau STOP} : c'est le seul panneau routier octogonal au monde. Sa forme unique permet de l'identifier instantanément, même s'il est retourné, sale ou partiellement caché.
\item Les \textbf{tables} et plateaux : beaucoup de tables de salon et de plateaux de service sont octogonaux, un compromis élégant entre le carré et le cercle.
\item Les \textbf{pièces de monnaie} et médailles : certaines pièces ont un contour octogonal pour être distinguées au toucher.
\item L'\textbf{architecture} : tours, coupoles, bassins et pavillons octogonaux (par exemple les fonts baptismaux de nombreuses églises) exploitent la solidité et la symétrie de cette figure.
\end{itemize}

\begin{savaistu}[Pour aller plus loin — hors exigibilité au BEPC]
Quelle est la longueur du côté d'un octogone régulier inscrit dans un cercle de rayon $R$ ? En utilisant la trigonométrie (que tu approfondiras au lycée), on montre que :
\[ a = R\sqrt{2 - \sqrt{2}} \approx 0,765\,R. \]
Par exemple, pour $R = 10$ cm, le côté vaut environ $7,65$ cm. Cette formule n'est \textbf{pas exigible} au BEPC, mais elle te sera utile au lycée — et elle explique pourquoi la construction géométrique par les bissectrices est si précieuse : elle évite tout calcul !
\end{savaistu}

\courssub{Exercice de synthèse}

\begin{exoresolu}[De l'hexagone à l'octogone]
Un artisan de Yaoundé fabrique deux plateaux en bois, tous deux inscrits dans des cercles de rayon $R = 30$ cm : l'un est hexagonal, l'autre octogonal.
\begin{enumerate}
\item Calcule le périmètre du plateau hexagonal.
\item Calcule l'aire exacte du plateau hexagonal.
\item Quel est l'angle intérieur du plateau octogonal ?
\end{enumerate}
\tcblower
\textbf{1.} Pour l'hexagone régulier, $a = R = 30$ cm, donc $\mathcal{P} = 6a = 6 \times 30 = 180$ cm.

\textbf{2.} $\mathcal{A} = \dfrac{3\sqrt{3}}{2}a^2 = \dfrac{3\sqrt{3}}{2} \times 900 = 1350\sqrt{3}$ cm², soit environ $2338$ cm².

\textbf{3.} L'angle intérieur de l'octogone régulier vaut $\dfrac{(8-2)\times 180°}{8} = 135°$.
\end{exoresolu}

\begin{exercice}[À toi de jouer]
\begin{enumerate}
\item Construis un cercle de rayon 5 cm, puis un hexagone régulier inscrit en reportant le rayon. Calcule son périmètre et son aire exacte.
\item Dans un cercle de centre $O$, trace un carré inscrit $ABCD$, puis construis un octogone régulier en traçant les médiatrices des côtés du carré. Justifie que les huit angles au centre obtenus valent $45°$.
\item Un panneau STOP a la forme d'un octogone régulier. Calcule la mesure de chacun de ses angles intérieurs, puis la somme de ses huit angles intérieurs.
\end{enumerate}
\end{exercice}

\begin{corrige}[Exercice : À toi de jouer]
\textbf{1.} L'hexagone a pour côté $a = R = 5$ cm. Périmètre : $\mathcal{P} = 6 \times 5 = 30$ cm. Aire : $\mathcal{A} = \dfrac{3\sqrt{3}}{2} \times 25 = \dfrac{75\sqrt{3}}{2}$ cm² $\approx 65$ cm².

\textbf{2.} Les angles au centre du carré valent $90°$. La médiatrice d'un côté $[AB]$ passe par $O$ (triangle $OAB$ isocèle en $O$) et partage l'angle $\widehat{AOB}$ en deux angles égaux de $\dfrac{90°}{2} = 45°$. On obtient ainsi 8 angles au centre de $45°$, et $8 \times 45° = 360°$ : l'octogone est régulier.

\textbf{3.} Angle intérieur : $\dfrac{(8-2)\times 180°}{8} = 135°$. Somme : $8 \times 135° = 1080°$ (ce qui redonne bien $(8-2)\times 180°$).
\end{corrige}

Tu maîtrises maintenant les quatre polygones réguliers du programme : triangle équilatéral, carré, hexagone et octogone. Dans la section suivante, nous verrons comment calculer l'\cle{apothème}, le périmètre et l'aire de \emph{n'importe quel} polygone régulier grâce à une formule générale.

\coursec{Apothème, périmètre et aire d'un polygone régulier}

Dans la section précédente, nous avons appris à construire et à reconnaître les polygones réguliers classiques : le triangle équilatéral, le carré, l'hexagone et l'octogone réguliers. Nous savons désormais qu'un polygone régulier possède un \cle{centre} $O$, que tous ses sommets sont sur un même cercle de rayon $R$ (le \cle{cercle circonscrit}), et que tous ses côtés ont la même longueur. Il reste une question très concrète, celle que se posent l'artisan qui découpe une table octogonale, le carreleur qui achète des tomettes hexagonales ou le mécanicien qui dessine un écrou : \emph{quelle surface occupe un polygone régulier, et comment la calculer ?} C'est toute l'ambition de cette section : découvrir une formule simple et élégante, $\mathcal{A} = \dfrac{\text{périmètre} \times \text{apothème}}{2}$, et apprendre à l'utiliser dans des situations réelles.

\courssub{L'apothème : la distance du centre à un côté}

\textbf{Une intuition pour commencer.} Place un hexagone régulier à plat sur ta table. Son centre $O$ est à la même distance de chaque sommet : c'est le rayon $R$ du cercle circonscrit. Mais $O$ est aussi à la même distance de chaque \emph{côté}. Si tu abaisses la perpendiculaire depuis $O$ sur un côté $[AB]$, elle le rencontre en son milieu $H$ (car le triangle $OAB$ est isocèle en $O$, et dans un triangle isocèle la hauteur issue du sommet principal est aussi médiane). La longueur $OH$ est ce qu'on appelle l'\cle{apothème}.

\begin{definition}[Apothème d'un polygone régulier]
Soit un polygone régulier de centre $O$. On appelle \cle{apothème} du polygone la distance du centre $O$ à l'un quelconque de ses côtés, c'est-à-dire la longueur $OH$, où $H$ est le pied de la perpendiculaire abaissée de $O$ sur ce côté.\\
Comme le triangle $OAB$ formé par le centre et un côté $[AB]$ est isocèle en $O$, le point $H$ est le \textbf{milieu} de $[AB]$, et l'apothème est aussi la \textbf{hauteur} de ce triangle isocèle élémentaire. Tous les côtés donnent la même longueur $OH$ : on parle donc de \emph{l'}apothème du polygone.
\end{definition}

\begin{attention}[Apothème ou rayon ? Ne les confonds plus !]
C'est l'erreur la plus fréquente dans cette leçon :
\begin{itemize}
\item le \cle{rayon} $R$ va du \textbf{centre à un sommet} : $R = OA = OB$ ;
\item l'\cle{apothème} $a_p$ va du \textbf{centre au milieu d'un côté} : $a_p = OH$, perpendiculairement au côté.
\end{itemize}
Dans le triangle $OHA$ rectangle en $H$, l'hypoténuse est $OA = R$. Or l'hypoténuse est le plus grand côté d'un triangle rectangle : on a donc toujours $a_p < R$, l'apothème est \emph{strictement plus petit} que le rayon. Si ton calcul donne un apothème plus grand que le rayon, c'est qu'il y a une erreur !
\end{attention}

\begin{popfigure}
\begin{center}
\begin{tikzpicture}[scale=1.15]
\coordinate (O) at (0,0);
\foreach \i in {0,...,5} {
  \coordinate (A\i) at ({60*\i+30}:2.6);
}
\draw[popmuted, thin] (0,0) circle (2.6);
\draw[popblue, thick] (A0) -- (A1) -- (A2) -- (A3) -- (A4) -- (A5) -- cycle;
\draw[poporange, thick, fill=poporangeL] (O) -- (A0) -- (A1) -- cycle;
\coordinate (H) at ($(A0)!0.5!(A1)$);
\draw[red, very thick] (O) -- (H);
\draw[popdark] (O) -- (A0);
\draw[popdark] (O) -- (A1);
\fill (O) circle (1.5pt) node[below left] {$O$};
\fill (A0) circle (1.5pt) node[right] {$A$};
\fill (A1) circle (1.5pt) node[above right] {$B$};
\fill[red] (H) circle (1.5pt) node[above right, red] {$H$};
\node[popdark] at (0.85,0.55) {$R$};
\node[red] at (0.45,1.15) {$a_p$};
\draw[red] ($(H)+(-0.13,-0.075)$) -- ($(H)+(-0.26,0)$) -- ($(H)+(-0.13,0.075)$);
\end{tikzpicture}
\end{center}
\legende{Hexagone régulier de centre $O$ : le triangle élémentaire $OAB$ est mis en évidence, et l'apothème $OH$ (en rouge) est la hauteur de ce triangle, perpendiculaire au côté $[AB]$ en son milieu.}
\end{popfigure}

\courssub{Décomposition en triangles isocèles et formule de l'aire}

\textbf{L'idée directrice.} Un polygone régulier à $n$ côtés, de centre $O$, peut être « découpé en parts de gâteau » : en joignant $O$ à chacun des $n$ sommets, on obtient exactement $n$ triangles isocèles \textbf{tous superposables} (mêmes longueurs : deux côtés égaux à $R$, une base égale au côté $c$ du polygone). L'aire du polygone est alors simplement la somme des aires de ces $n$ triangles identiques. Cette idée simple va nous donner l'une des plus belles formules de la géométrie du collège.

\begin{experience}[Démonstration guidée : l'aire d'un polygone régulier]
Considérons un polygone régulier à $n$ côtés, de centre $O$, de côté $c$ et d'apothème $a_p$. Suivons le raisonnement pas à pas.
\begin{enumerate}
\item \textbf{Découpage.} Les segments joignant $O$ aux sommets partagent le polygone en $n$ triangles isocèles superposables, comme $OAB$ sur la figure.
\item \textbf{Aire d'un triangle.} Dans le triangle $OAB$, prenons $[AB]$ comme base : $AB = c$. La hauteur correspondante est $OH = a_p$ (l'apothème, perpendiculaire à $[AB]$). Donc :
\[ \mathcal{A}_{OAB} = \frac{\text{base} \times \text{hauteur}}{2} = \frac{c \times a_p}{2}. \]
\item \textbf{Somme des $n$ triangles.} Les $n$ triangles ont la même aire, donc :
\[ \mathcal{A} = n \times \frac{c \times a_p}{2} = \frac{(n \times c) \times a_p}{2}. \]
\item \textbf{Conclusion.} Or $n \times c$ est la somme des longueurs des $n$ côtés, c'est-à-dire le \textbf{périmètre} $P$ du polygone. D'où la formule annoncée.
\end{enumerate}
\end{experience}

\begin{propriete}[Aire d'un polygone régulier]
L'aire $\mathcal{A}$ d'un polygone régulier de périmètre $P$ et d'apothème $a_p$ est :
\[ \mathcal{A} = \frac{P \times a_p}{2} \qquad \text{c'est-à-dire} \qquad \mathcal{A} = \frac{n \times c \times a_p}{2}, \]
où $n$ est le nombre de côtés et $c$ la longueur d'un côté. Cette démonstration est formatrice et peut être demandée au BEPC sous forme guidée.
\end{propriete}

\begin{aretenir}[À retenir absolument]
\begin{itemize}
\item Apothème $a_p$ : distance du centre au \textbf{milieu d'un côté}.
\item Aire : $\mathcal{A} = \dfrac{\text{périmètre} \times \text{apothème}}{2}$.
\item Le périmètre, lui, est immédiat : $P = n \times c$.
\end{itemize}
\end{aretenir}

\courssub{Calculer l'apothème et le côté : Pythagore et trigonométrie}

La formule de l'aire est belle, mais encore faut-il connaître l'apothème ! Tout se joue dans le \textbf{triangle rectangle} $OHA$ :
\begin{itemize}
\item l'hypoténuse est $OA = R$ (le rayon) ;
\item $AH = \dfrac{c}{2}$ (car $H$ est le milieu de $[AB]$) ;
\item l'angle au centre du polygone vaut $\dfrac{360°}{n}$, donc $\widehat{AOH} = \dfrac{1}{2} \times \dfrac{360°}{n} = \dfrac{180°}{n}$ (la hauteur $[OH]$ du triangle isocèle $OAB$ est aussi bissectrice de l'angle $\widehat{AOB}$).
\end{itemize}

\begin{propriete}[Apothème et côté en fonction du rayon]
Pour un polygone régulier à $n$ côtés inscrit dans un cercle de rayon $R$ :
\[ a_p = R \cos\left(\frac{180°}{n}\right) \qquad \text{et} \qquad c = 2R\sin\left(\frac{180°}{n}\right). \]
En effet, dans le triangle $OHA$ rectangle en $H$ : $\cos(\widehat{AOH}) = \dfrac{OH}{OA} = \dfrac{a_p}{R}$ et $\sin(\widehat{AOH}) = \dfrac{AH}{OA} = \dfrac{c/2}{R}$.
\end{propriete}

Quand on connaît le rayon $R$ et le côté $c$, on peut aussi utiliser le \textbf{théorème de Pythagore} dans le triangle $OHA$ rectangle en $H$ : $OH^2 + AH^2 = OA^2$, c'est-à-dire
\[ a_p = \sqrt{R^2 - \left(\frac{c}{2}\right)^2}. \]

\begin{methode}[Calculer l'aire d'un polygone régulier]
\begin{enumerate}
\item \textbf{Identifier les données} : nombre de côtés $n$, longueur du côté $c$ ou rayon $R$ du cercle circonscrit.
\item \textbf{Calculer l'apothème} : soit par la trigonométrie $a_p = R\cos\left(\dfrac{180°}{n}\right)$, soit par le théorème de Pythagore dans le triangle $OHA$.
\item \textbf{Calculer le périmètre} : $P = n \times c$.
\item \textbf{Appliquer la formule} : $\mathcal{A} = \dfrac{P \times a_p}{2}$, puis conclure avec l'unité d'aire et une phrase de réponse.
\end{enumerate}
\end{methode}

\begin{exoresolu}[Apothème et aire d'un hexagone régulier de côté 6 cm]
Une tomette (carreau de sol) a la forme d'un hexagone régulier de côté $c = 6$~cm. Calculer son apothème, puis son aire (arrondir au dixième).
\tcblower
\textbf{Solution détaillée.}
\begin{enumerate}
\item \textbf{Rayon.} Propriété remarquable de l'hexagone régulier : son côté est égal au rayon du cercle circonscrit, donc $R = c = 6$~cm.
\item \textbf{Apothème.} L'angle $\widehat{AOH}$ vaut $\dfrac{180°}{6} = 30°$. Donc :
\[ a_p = R\cos(30°) = 6 \times \frac{\sqrt{3}}{2} = 3\sqrt{3} \approx 5{,}2~\text{cm}. \]
\emph{Vérification par Pythagore :} $a_p = \sqrt{6^2 - 3^2} = \sqrt{36 - 9} = \sqrt{27} = 3\sqrt{3}$. Les deux méthodes concordent.
\item \textbf{Périmètre.} $P = 6 \times 6 = 36$~cm.
\item \textbf{Aire.}
\[ \mathcal{A} = \frac{P \times a_p}{2} = \frac{36 \times 3\sqrt{3}}{2} = 54\sqrt{3} \approx 93{,}5~\text{cm}^2. \]
\end{enumerate}
\textbf{Conclusion :} l'apothème vaut $3\sqrt{3} \approx 5{,}2$~cm et la tomette a une aire d'environ $93{,}5$~cm$^2$.
\end{exoresolu}

\begin{exemplebox}[Octogone régulier inscrit dans un cercle de rayon 10 cm]
Un artisan de Foumban veut fabriquer un plateau octogonal régulier inscrit dans un cercle de rayon $R = 10$~cm. Calculons le côté, l'apothème et l'aire du plateau.
\begin{enumerate}
\item \textbf{Angle} : $\widehat{AOH} = \dfrac{180°}{8} = 22{,}5°$.
\item \textbf{Côté} : $c = 2R\sin(22{,}5°) \approx 20 \times 0{,}3827 \approx 7{,}65$~cm.
\item \textbf{Apothème} : $a_p = R\cos(22{,}5°) \approx 10 \times 0{,}9239 \approx 9{,}24$~cm.
\item \textbf{Périmètre} : $P = 8 \times 7{,}65 \approx 61{,}2$~cm.
\item \textbf{Aire} : $\mathcal{A} = \dfrac{P \times a_p}{2} \approx \dfrac{61{,}2 \times 9{,}24}{2} \approx 282{,}8$~cm$^2$.
\end{enumerate}
Le plateau a une aire d'environ $282{,}8$~cm$^2$. Remarque au passage : l'aire du disque de rayon 10~cm est $\pi R^2 \approx 314{,}2$~cm$^2$ : l'octogone occupe déjà environ $90\,\%$ du disque !
\end{exemplebox}

\begin{attention}[Pièges de rédaction et de calcul]
\begin{itemize}
\item \textbf{Ne pas oublier le facteur $\dfrac{1}{2}$} dans $\mathcal{A} = \dfrac{P \times a_p}{2}$ : beaucoup d'élèves écrivent $P \times a_p$ et doublent l'aire !
\item \textbf{Unités} : le périmètre est en cm, l'apothème en cm, donc l'aire est en cm$^2$. Précise toujours l'unité dans la conclusion.
\item \textbf{Calculatrice en degrés} : les angles $22{,}5°$, $30°$ sont en degrés ; vérifie le réglage de ta calculatrice.
\item \textbf{Ne pas arrondir trop tôt} : garde les valeurs exactes (comme $3\sqrt{3}$) ou les valeurs approchées avec au moins 4 chiffres jusqu'au calcul final, et n'arrondis qu'à la toute fin.
\end{itemize}
\end{attention}

\courssub{À périmètre fixé : le polygone se rapproche du disque}

Terminons par une observation fascinante. Prenons une ficelle de longueur fixe, disons 12~cm, et formons successivement un triangle équilatéral, un carré, un hexagone, un octogone, un dodécagone (12 côtés) régulier... Le périmètre est toujours le même ($P = 12$~cm), mais qu'en est-il de l'aire ?

\begin{popfigure}
\begin{center}
\begin{tikzpicture}
\begin{axis}[
  width=0.85\linewidth, height=6cm,
  xlabel={Nombre de côtés $n$},
  ylabel={Aire (cm$^2$)},
  xmin=2, xmax=13, ymin=6, ymax=12.2,
  xtick={3,4,6,8,12},
  grid=both, grid style={popline!40},
  axis line style={popdark},
]
\addplot[popblue, dashed, thick, domain=2:13] {11.459};
\addplot[poporange, thick, smooth, mark=*, mark size=2.5pt] coordinates {
  (3,6.928) (4,9) (6,10.392) (8,10.863) (12,11.196)
};
\node[popblue, anchor=south west] at (axis cs:2.2,11.459) {\small Aire du disque de même périmètre};
\end{axis}
\end{tikzpicture}
\end{center}
\legende{Pour un périmètre fixé à 12 cm, l'aire du polygone régulier augmente avec le nombre de côtés et se rapproche de l'aire du disque de même périmètre (ligne pointillée).}
\end{popfigure}

\begin{savaistu}[Pourquoi les bulles et les gouttes sont rondes]
Le graphique ci-dessus illustre un fait célèbre : \textbf{à périmètre fixé, plus un polygone régulier a de côtés, plus son aire est grande}, et elle se rapproche de l'aire du disque de même périmètre, sans jamais la dépasser. Le cercle est le « champion » des formes : c'est la figure qui enferme la plus grande aire pour un périmètre donné. C'est pour cela que les bulles de savon, les gouttes d'eau et même les cernes dans l'eau adoptent une forme circulaire : la nature « économise » la surface ! Les Grecs anciens connaissaient déjà ce problème, appelé \emph{problème isopérimétrique}, et Archimède utilisait des polygones à 96 côtés pour encadrer le nombre $\pi$.
\end{savaistu}

\begin{exercice}[À toi de jouer]
Un panneau de signalisation « STOP » a la forme d'un octogone régulier de côté $c = 30$~cm.
\begin{enumerate}
\item Calculer le rayon $R$ du cercle circonscrit (on donne $\sin(22{,}5°) \approx 0{,}383$).
\item Calculer l'apothème du panneau (on donne $\cos(22{,}5°) \approx 0{,}924$).
\item En déduire l'aire de la face visible du panneau, arrondie au cm$^2$.
\end{enumerate}
\end{exercice}

\begin{corrige}[Exercice : le panneau STOP]
\begin{enumerate}
\item De $c = 2R\sin(22{,}5°)$, on tire $R = \dfrac{c}{2\sin(22{,}5°)} \approx \dfrac{30}{2 \times 0{,}383} = \dfrac{30}{0{,}766} \approx 39{,}2$~cm.
\item $a_p = R\cos(22{,}5°) \approx 39{,}2 \times 0{,}924 \approx 36{,}2$~cm.
\item Périmètre : $P = 8 \times 30 = 240$~cm. Aire : $\mathcal{A} = \dfrac{P \times a_p}{2} \approx \dfrac{240 \times 36{,}2}{2} = 120 \times 36{,}2 = 4344$~cm$^2$. La face visible du panneau a une aire d'environ $4344$~cm$^2$.
\end{enumerate}
\end{corrige}

En résumé, cette section t'a donné un outil puissant et très concret : dès qu'un polygone régulier apparaît (tomette, plateau, panneau, écrou), pense au découpage en triangles isocèles, calcule l'apothème grâce au triangle rectangle $OHA$ (par la trigonométrie ou par Pythagore), puis applique $\mathcal{A} = \dfrac{P \times a_p}{2}$. La section suivante te montrera comment mobiliser toutes ces notions dans des problèmes de synthèse issus de la vie courante.

\coursec{Applications concrètes et synthèse}

Dans la section précédente, nous avons appris à calculer l'\cle{apothème}, le \cle{périmètre} et l'\cle{aire} d'un polygone régulier. Ces formules ne sont pas de simples exercices de style : elles servent chaque jour aux maçons qui pavent une cour à Douala, aux menuisiers qui découpent une table à Bafoussam, ou aux propriétaires qui clôturent un champ à Mbouda. Dans cette dernière section, nous mettons tout cela en action à travers des problèmes concrets, puis nous dresserons le bilan complet des quatre polygones réguliers étudiés dans ce chapitre.

\courssub{Méthode générale de résolution d'un problème concret}

Face à un problème de la vie courante (pavage, découpe, clôture), le plus difficile n'est pas le calcul : c'est de \emph{traduire la situation en mathématiques}. Voici la démarche à suivre systématiquement, celle que l'on attend de toi au BEPC.

\begin{methode}[Résoudre un problème concret en 4 étapes]
\begin{enumerate}
\item \textbf{Modéliser par un schéma coté.} Dessine la situation à main levée, reporte toutes les dimensions de l'énoncé sur le schéma, et convertis-les dans la \textbf{même unité}.
\item \textbf{Identifier le polygone et les grandeurs utiles.} De quel polygone régulier s'agit-il ? Cherche-t-on une aire, un périmètre, un nombre d'objets ?
\item \textbf{Calculer.} Écris d'abord le \emph{calcul littéral} (avec les lettres ou les mots), puis remplace par les valeurs numériques.
\item \textbf{Conclure avec l'unité.} Rédige une phrase de conclusion complète, avec l'unité adaptée, et vérifie que le résultat est réaliste.
\end{enumerate}
\end{methode}

\begin{attention}[Le piège n° 1 : les unités]
Avant \textbf{tout} calcul d'aire ou de périmètre, convertis toutes les longueurs dans la même unité. Une cour de $6$ m pavée avec des pavés de $30$ cm de côté impose de tout passer en mètres ($30$ cm $= 0,3$ m) ou tout en centimètres ($6$ m $= 600$ cm). Si tu mélanges mètres et centimètres, ton aire sera fausse d'un facteur $10\,000$ ! Rappel : $1~\text{m}^2 = 10\,000~\text{cm}^2$.
\end{attention}

\courssub{Problème de pavage : combien de pavés pour une cour ?}

\textbf{Situation.} Une famille de Yaoundé veut paver sa cour rectangulaire de $6$ m de long sur $4,5$ m de large avec des pavés hexagonaux réguliers de $30$ cm de côté. Combien de pavés faut-il prévoir ?

\begin{popfigure}
\begin{center}
\begin{tikzpicture}[scale=0.75]
% Cour rectangulaire
\draw[fill=popcream,draw=popdark,thick] (0,0) rectangle (6,4.5);
% Cotes
\draw[<->,>=stealth,popblue] (0,-0.5) -- (6,-0.5) node[midway,below]{$6$ m};
\draw[<->,>=stealth,popblue] (-0.5,0) -- (-0.5,4.5) node[midway,left]{$4{,}5$ m};
% Quelques pavés hexagonaux
\begin{scope}
\foreach \x/\y in {0.9/0.9, 1.9/0.9, 2.9/0.9, 1.4/1.75, 2.4/1.75}{
\draw[fill=poporangeL,draw=poporange] (\x,\y) ++(0:0.3) \foreach \a in {60,120,180,240,300}{ -- ++(\a:0.3) } -- cycle;
}
\end{scope}
% Pavé légende avec cote
\begin{scope}[shift={(4.6,3.2)}]
\draw[fill=poporangeL,draw=poporange,thick] (0:0.45) \foreach \a in {60,120,180,240,300}{ -- ++(\a:0.45) } -- cycle;
\draw[<->,>=stealth,popblue] (0:0.45) -- (60:0.45) node[midway,right,font=\small]{$30$ cm};
\end{scope}
\node[font=\small\itshape,popmuted] at (3,4.1) {Cour rectangulaire à paver};
\end{tikzpicture}
\end{center}
\legende{Plan coté de la cour et d'un pavé hexagonal de $30$ cm de côté.}
\end{popfigure}

\textbf{Étape 1 et 2 : modélisation.} L'hexagone régulier est un excellent pavé : ses angles de $120°$ permettent aux pavés de s'emboîter sans laisser de trous (trois hexagones se rejoignent en chaque sommet : $3 \times 120° = 360°$). Pour estimer le nombre de pavés, on divise l'aire de la cour par l'aire d'un pavé.

\textbf{Étape 3 : calculs.} Aire de la cour :
\[ \mathcal{A}_{\text{cour}} = 6 \times 4{,}5 = 27~\text{m}^2. \]
Aire d'un pavé hexagonal de côté $c = 0{,}3$ m. On utilise la formule vue dans la section précédente : l'hexagone régulier se décompose en $6$ triangles équilatéraux de côté $c$, donc
\[ \mathcal{A}_{\text{pavé}} = 6 \times \frac{c^2\sqrt{3}}{4} = \frac{3\sqrt{3}}{2}c^2 = \frac{3\sqrt{3}}{2}\times 0{,}09 \approx 0{,}234~\text{m}^2. \]
Nombre de pavés :
\[ N = \frac{27}{0{,}234} \approx 115{,}4. \]

\textbf{Étape 4 : conclusion.} Il faut donc environ $116$ pavés. Mais les pavés au bord de la cour devront être découpés, ce qui produit des chutes : un maçon prudent commande une marge de $5$ \% à $10$ \%, soit environ \textbf{125 pavés}.

\begin{exemplebox}[Et avec des pavés carrés ?]
Avec des pavés carrés de $30$ cm de côté : aire d'un pavé $= 0{,}3^2 = 0{,}09~\text{m}^2$, donc $N = \dfrac{27}{0{,}09} = 300$ pavés exactement. Le carré pave parfaitement un rectangle quand les dimensions sont des multiples du côté : ici $6 \div 0{,}3 = 20$ pavés en longueur et $4{,}5 \div 0{,}3 = 15$ en largeur, soit $20 \times 15 = 300$. Aucune découpe !
\end{exemplebox}

\courssub{Problème de découpe : une table octogonale dans un disque}

\textbf{Situation.} Un menuisier découpe une table octogonale régulière dans une plaque circulaire de contreplaqué de rayon $R = 60$ cm. Quelle est l'aire des chutes de matériau ?

\textbf{Modélisation.} La table est un octogone régulier inscrit dans le cercle de rayon $R = 60$ cm. Les chutes sont la partie du disque non recouverte par l'octogone :
\[ \mathcal{A}_{\text{chutes}} = \mathcal{A}_{\text{disque}} - \mathcal{A}_{\text{octogone}}. \]

\textbf{Calculs.} Aire du disque :
\[ \mathcal{A}_{\text{disque}} = \pi R^2 = \pi \times 60^2 = 3600\pi \approx 11\,310~\text{cm}^2. \]
Pour l'octogone, on le décompose en $8$ triangles isocèles de sommet le centre, d'angle au centre $\dfrac{360°}{8} = 45°$. L'aire de chaque triangle est $\dfrac{1}{2}R^2 \sin(45°)$, donc :
\[ \mathcal{A}_{\text{octogone}} = 8 \times \frac{1}{2}R^2 \sin(45°) = 4 \times 3600 \times \frac{\sqrt{2}}{2} = 7200\sqrt{2} \approx 10\,182~\text{cm}^2. \]
D'où :
\[ \mathcal{A}_{\text{chutes}} \approx 11\,310 - 10\,182 = 1\,128~\text{cm}^2. \]

\textbf{Conclusion.} Les chutes représentent environ $1\,128~\text{cm}^2$, soit environ $10$ \% de la plaque. Le menuisier pourra les réutiliser pour de petites pièces.

\begin{attention}[Bien identifier la donnée : côté ou rayon ?]
Dans un problème de découpe, vérifie si l'énoncé donne le \textbf{côté} du polygone ou le \textbf{rayon} du cercle circonscrit. Les formules ne sont pas les mêmes ! Pour l'octogone inscrit dans un cercle de rayon $R$, l'aire est $2\sqrt{2}\,R^2$ ; si l'on te donne le côté $c$, il faut d'abord relier $c$ à $R$ (par exemple avec l'angle au centre et la trigonométrie).
\end{attention}

\courssub{Problème d'optimisation : quel enclos pour le même périmètre ?}

\textbf{Situation.} Un éleveur de l'Adamaoua dispose de $120$ m de fil de clôture. Vaut-il mieux construire un enclos carré ou un enclos hexagonal régulier pour offrir le plus d'espace à ses animaux ?

\begin{exoresolu}[Enclos carré contre enclos hexagonal]
Un éleveur possède $120$ m de clôture. Comparer l'aire d'un enclos carré et celle d'un enclos hexagonal régulier de même périmètre.
\tcblower
\textbf{Enclos carré.} Le périmètre est $4c = 120$ m, donc $c = 30$ m.
\[ \mathcal{A}_{\text{carré}} = c^2 = 30^2 = 900~\text{m}^2. \]
\textbf{Enclos hexagonal régulier.} Le périmètre est $6c = 120$ m, donc $c = 20$ m.
\[ \mathcal{A}_{\text{hexagone}} = \frac{3\sqrt{3}}{2}c^2 = \frac{3\sqrt{3}}{2}\times 400 = 600\sqrt{3} \approx 1\,039~\text{m}^2. \]
\textbf{Conclusion.} À périmètre égal, l'enclos hexagonal offre environ $1\,039 - 900 = 139~\text{m}^2$ de plus que l'enclos carré, soit un gain d'environ $15$ \%. Plus un polygone régulier a de côtés, plus son aire est grande à périmètre fixé.
\end{exoresolu}

\begin{savaistu}[La forme idéale : le cercle]
À périmètre égal, c'est le \textbf{cercle} qui possède la plus grande aire. Vérifions : un cercle de périmètre $120$ m a un rayon $R = \dfrac{120}{2\pi} \approx 19{,}1$ m, donc une aire $\pi R^2 \approx 1\,146~\text{m}^2$, encore supérieure à celle de l'hexagone ! C'est pourquoi les abeilles construisent leurs alvéoles en \textbf{hexagones réguliers} : l'hexagone est le polygone qui s'approche le mieux du cercle idéal \emph{tout en s'emboîtant parfaitement} avec ses voisins, sans laisser aucun vide. Un remarquable compromis entre économie de cire et espace de stockage, que la nature a « découvert » bien avant les mathématiciens !
\end{savaistu}

\courssub{Lecture et exploitation de plans}

Au BEPC, de nombreux problèmes partent d'un plan : piscine hexagonale, kiosque octogonal, dalles de terrasse. La clé est toujours la même : repérer le polygone régulier, relever ses dimensions sur le plan, puis appliquer la bonne formule.

\begin{exemplebox}[Piscine hexagonale et kiosque octogonal]
\textbf{Piscine.} Une piscine a la forme d'un hexagone régulier de côté $2$ m. Son emprise au sol vaut $\mathcal{A} = \dfrac{3\sqrt{3}}{2}\times 2^2 = 6\sqrt{3} \approx 10{,}4~\text{m}^2$. Son périmètre ($6 \times 2 = 12$ m) donne la longueur de la margelle à poser.

\textbf{Kiosque.} Un kiosque octogonal régulier a un côté de $1{,}5$ m. Son périmètre est $8 \times 1{,}5 = 12$ m : c'est la longueur de rampe nécessaire pour en faire le tour. Pour son aire, on peut le décomposer en $8$ triangles isocèles d'angle au centre $45°$, ou utiliser l'apothème : $\mathcal{A} = \dfrac{\text{périmètre} \times \text{apothème}}{2}$.
\end{exemplebox}

\begin{exercice}[Dalles triangulaires]
Une terrasse rectangulaire de $3{,}6$ m sur $2{,}4$ m est recouverte de dalles triangulaires équilatérales de $40$ cm de côté. (1) Calculer l'aire de la terrasse en $\text{m}^2$. (2) Calculer l'aire d'une dalle en $\text{m}^2$ (on prendra $\sqrt{3} \approx 1{,}73$). (3) En déduire le nombre de dalles nécessaires, arrondi à l'entier supérieur.
\end{exercice}

\begin{corrige}[Exercice : dalles triangulaires]
(1) $\mathcal{A}_{\text{terrasse}} = 3{,}6 \times 2{,}4 = 8{,}64~\text{m}^2$.

(2) Conversion : $40$ cm $= 0{,}4$ m. Aire d'une dalle :
\[ \mathcal{A}_{\text{dalle}} = \frac{c^2\sqrt{3}}{4} = \frac{0{,}16 \times 1{,}73}{4} \approx 0{,}069~\text{m}^2. \]

(3) Nombre de dalles : $N = \dfrac{8{,}64}{0{,}069} \approx 125{,}2$. Il faut donc prévoir au moins $126$ dalles, plus une marge pour les découpes.
\end{corrige}

\courssub{Bilan de synthèse : les quatre polygones réguliers du chapitre}

\begin{popfigure}
\begin{center}
\begin{tikzpicture}[scale=0.62]
% Triangle équilatéral
\begin{scope}[shift={(0,0)}]
\draw[popmuted] (0,0) circle (1.1);
\draw[fill=popblueL,draw=popblue,thick] (90:1.1) -- (210:1.1) -- (330:1.1) -- cycle;
\draw[poporange,->,>=stealth] (0.55,0) arc (0:120:0.55);
\node[poporange,font=\small] at (0.75,0.55) {$120°$};
\node[font=\small\bfseries,popdark] at (0,-1.7) {Triangle équilatéral};
\node[font=\small,popdark] at (0,-2.25) {$\mathcal{A}=\dfrac{c^2\sqrt{3}}{4}$};
\end{scope}
% Carré
\begin{scope}[shift={(4,0)}]
\draw[popmuted] (0,0) circle (1.1);
\draw[fill=popgreenL,draw=popgreen,thick] (45:1.1) -- (135:1.1) -- (225:1.1) -- (315:1.1) -- cycle;
\draw[poporange,->,>=stealth] (0.55,0) arc (0:90:0.55);
\node[poporange,font=\small] at (0.8,0.5) {$90°$};
\node[font=\small\bfseries,popdark] at (0,-1.7) {Carré};
\node[font=\small,popdark] at (0,-2.25) {$\mathcal{A}=c^2$};
\end{scope}
% Hexagone
\begin{scope}[shift={(8,0)}]
\draw[popmuted] (0,0) circle (1.1);
\draw[fill=poporangeL,draw=poporange,thick] (0:1.1) \foreach \a in {60,120,180,240,300} { -- (\a:1.1) } -- cycle;
\draw[popblue,->,>=stealth] (0.55,0) arc (0:60:0.55);
\node[popblue,font=\small] at (0.85,0.4) {$60°$};
\node[font=\small\bfseries,popdark] at (0,-1.7) {Hexagone régulier};
\node[font=\small,popdark] at (0,-2.25) {$\mathcal{A}=\dfrac{3\sqrt{3}}{2}c^2$};
\end{scope}
% Octogone
\begin{scope}[shift={(12,0)}]
\draw[popmuted] (0,0) circle (1.1);
\draw[fill=poppurpleL,draw=poppurple,thick] (0:1.1) \foreach \a in {45,90,135,180,225,270,315} { -- (\a:1.1) } -- cycle;
\draw[popblue,->,>=stealth] (0.55,0) arc (0:45:0.55);
\node[popblue,font=\small] at (0.85,0.35) {$45°$};
\node[font=\small\bfseries,popdark] at (0,-1.7) {Octogone régulier};
\node[font=\small,popdark] at (0,-2.25) {$\mathcal{A}=\dfrac{P \times a}{2}$};
\end{scope}
\end{tikzpicture}
\end{center}
\legende{Les quatre polygones réguliers étudiés, inscrits dans leur cercle circonscrit, avec leur angle au centre et leur formule d'aire ($c$ = côté, $P$ = périmètre, $a$ = apothème).}
\end{popfigure}

\begin{aretenir}[Tableau récapitulatif à connaître par cœur]
\begin{center}
\begin{tabularx}{\linewidth}{|l|X|X|X|X|}
\hline
\rowcolor{popblueL} & \textbf{Triangle équilatéral} & \textbf{Carré} & \textbf{Hexagone} & \textbf{Octogone} \\
\hline Nombre de côtés & $3$ & $4$ & $6$ & $8$ \\
\hline Angle au centre & $120°$ & $90°$ & $60°$ & $45°$ \\
\hline Angle intérieur & $60°$ & $90°$ & $120°$ & $135°$ \\
\hline Périmètre & $3c$ & $4c$ & $6c$ & $8c$ \\
\hline Aire & $\dfrac{c^2\sqrt{3}}{4}$ & $c^2$ & $\dfrac{3\sqrt{3}}{2}c^2$ & $\dfrac{P \times a}{2}$ \\
\hline Construction & $3$ arcs de rayon $R$ & $2$ diamètres perpendiculaires & $6$ reports du rayon & Bissectrices du carré \\
\hline
\end{tabularx}
\end{center}
Formule universelle valable pour tous : $\mathcal{A} = \dfrac{\text{périmètre} \times \text{apothème}}{2}$.
\end{aretenir}

\begin{methode}[Réflexes finaux pour le BEPC]
\begin{enumerate}
\item Je dessine un schéma coté et je convertis toutes les longueurs dans la même unité.
\item J'identifie le polygone régulier et je choisis la formule adaptée (ou la formule universelle avec l'apothème).
\item J'écris le calcul littéral, puis le calcul numérique avec une valeur approchée proprement arrondie.
\item Je conclus par une phrase avec l'unité, et je vérifie l'ordre de grandeur de ma réponse.
\end{enumerate}
\end{methode}

Tu disposes maintenant de toutes les clés du chapitre : définitions, angles, constructions à la règle et au compas, apothème, aires et problèmes concrets. Entraîne-toi sur les exercices de synthèse de l'unité : les polygones réguliers sont un classique du BEPC, et une copie bien rédigée — schéma, calcul littéral, conclusion avec unité — y fait toujours la différence.

\newpage

\coursec{Activité d'intégration}

\courssub{Objectifs de l'activité}

Cette activité d'intégration clôt le chapitre sur les \cle{polygones réguliers}. En travaillant par groupes de 3 ou 4 élèves, tu vas mobiliser l'ensemble des savoirs et savoir-faire de la leçon pour résoudre un problème authentique de la vie courante. À l'issue de cette activité, tu sauras :

\begin{itemize}
\item calculer un \cle{angle intérieur} d'un polygone régulier et vérifier qu'un pavage est possible ;
\item calculer l'\cle{aire} d'un triangle équilatéral, d'un carré et d'un hexagone régulier ;
\item résoudre un problème concret de quantité et de coût (nombre de pavés, budget) ;
\item construire un motif de pavage à la règle et au compas ;
\item rédiger et justifier une démarche complète, puis argumenter un choix.
\end{itemize}

\courssub{Situation}

\textbf{Le pavage de la cour de Maman Aïcha.}

Maman Aïcha, commerçante au quartier Nkol-Eton à Yaoundé, vient de terminer la construction de sa maison. Sa cour, de forme rectangulaire, mesure \SI{6}{m} de long sur \SI{4,5}{m} de large. Fatiguée de la boue en saison des pluies, elle décide de paver entièrement cette cour.

Chez le fournisseur de matériaux de construction, on lui propose trois modèles de pavés en forme de \cle{polygones réguliers}, tous de côté $c = \SI{30}{cm}$ :

\begin{center}
\begin{tabularx}{\linewidth}{|l|X|X|}\hline
\rowcolor{popblueL} \textbf{Modèle} & \textbf{Forme du pavé} & \textbf{Prix au m\textsuperscript{2}} \\ \hline
Modèle T & Triangle équilatéral de côté \SI{30}{cm} & \num{12000} F \\ \hline
Modèle C & Carré de côté \SI{30}{cm} & \num{10000} F \\ \hline
Modèle H & Hexagone régulier de côté \SI{30}{cm} & \num{15000} F \\ \hline
\end{tabularx}
\end{center}

Le fournisseur affirme que ces trois modèles permettent de paver la cour « sans trous ni chevauchements », mais que des pavés en forme de \emph{pentagone régulier}, eux, ne conviendraient pas. Il conseille en outre d'acheter $10\,\%$ de pavés en plus pour compenser les découpes le long des bords.

Maman Aïcha, qui veut dépenser intelligemment, sollicite ta classe. Ta mission, en groupe : vérifier les affirmations du fournisseur, comparer les trois options et lui rédiger un conseil argumenté.

\courssub{Consignes}

\begin{enumerate}
\item \textbf{Étude des angles.} Calculer l'angle intérieur d'un triangle équilatéral, d'un carré, d'un hexagone régulier et d'un pentagone régulier. En déduire, en expliquant la condition qu'un pavage doit vérifier autour de chaque sommet, que les trois modèles proposés peuvent paver la cour sans trous, et que le pentagone régulier ne convient pas.
\item \textbf{Aire d'un pavé.} Calculer l'aire exacte d'un pavé de chaque type, en m\textsuperscript{2}. Pour le triangle équilatéral, on calculera d'abord sa hauteur ; pour l'hexagone, on le décomposera en 6 triangles équilatéraux.
\item \textbf{Quantité et coût.} Calculer l'aire de la cour. Pour chaque modèle, estimer le nombre de pavés nécessaires (arrondir à l'entier supérieur), ajouter $10\,\%$ de marge pour les découpes, puis déterminer le coût total en utilisant le prix au m\textsuperscript{2} (on comptera l'aire réellement achetée, marge comprise).
\item \textbf{Construction.} À l'échelle $1/50$, construire à la règle et au compas un motif du pavage hexagonal (au moins 7 hexagones : un central et 6 autour).
\item \textbf{Conseil rédigé.} Rédiger un texte de 10 à 15 lignes conseillant Maman Aïcha : quel modèle choisir ? Justifier par le coût, mais aussi par le nombre de découpes et l'esthétique. Chaque groupe présentera sa démarche au tableau.
\end{enumerate}

\courssub{Ressources à mobiliser}

\begin{aretenir}[Boîte à outils de la leçon]
\begin{itemize}
\item La somme des angles intérieurs d'un polygone à $n$ côtés vaut $(n-2)\times 180°$.
\item Chaque angle intérieur d'un polygone régulier à $n$ côtés mesure $\dfrac{(n-2)\times 180°}{n}$.
\item \textbf{Condition de pavage} : autour d'un sommet commun, la somme des angles des pavés doit être exactement $360°$.
\item Hauteur d'un triangle équilatéral de côté $c$ : $h = \dfrac{c\sqrt{3}}{2}$ ; aire : $\mathcal{A} = \dfrac{c \times h}{2} = \dfrac{c^2\sqrt{3}}{4}$.
\item Aire d'un carré de côté $c$ : $\mathcal{A} = c^2$.
\item Un hexagone régulier de côté $c$ se décompose en 6 triangles équilatéraux de côté $c$ : $\mathcal{A}_{\text{hex}} = \dfrac{3c^2\sqrt{3}}{2}$.
\item Aire d'un rectangle : $\mathcal{A} = L \times \ell$.
\end{itemize}
\end{aretenir}

\begin{popfigure}
\begin{center}
\begin{tikzpicture}[scale=0.9]
% Pavage hexagonal : 1 hexagone central + 6 voisins (côté 1 cm)
\foreach \cx/\cy in {0/0, 1.732/0, -1.732/0, 0.866/1.5, -0.866/1.5, 0.866/-1.5, -0.866/-1.5}{
  \draw[popblue, thick, fill=popblueL] (\cx,\cy)
    \foreach \a in {0,60,...,300}{ -- ++(\a:1) } -- cycle;
}
\draw[poporange, very thick] (0,0) circle (1);
\fill[popdark] (0,0) circle (1.5pt) node[below left=2pt, popdark] {$O$};
\draw[poporange, thick] (0,0) -- (1,0) node[midway, below, popdark] {\small $R$};
\draw[poporange, thick] (1,0) -- (0.5,0.866) node[midway, right, popdark] {\small $c$};
\end{tikzpicture}
\end{center}
\legende{Motif de pavage hexagonal : un hexagone central entouré de 6 hexagones. Pour l'hexagone régulier, le côté $c$ est égal au rayon $R$ du cercle circonscrit : c'est la clé de la construction au compas.}
\end{popfigure}

\courssub{Démarche et corrigé commenté}

\begin{exoresolu}[Le pavage de la cour de Maman Aïcha]
Reprendre les cinq tâches de la situation : angles et pavage, aires des pavés, quantités et coûts, construction, conseil à Maman Aïcha.
\tcblower
\textbf{Étape 1 : les angles et la possibilité du pavage.}

On applique la formule $\widehat{A} = \dfrac{(n-2)\times 180°}{n}$ :
\begin{align*}
\text{Triangle }(n=3) :&\quad \widehat{A} = \frac{1\times 180°}{3} = 60° ;\\
\text{Carré }(n=4) :&\quad \widehat{A} = \frac{2\times 180°}{4} = 90° ;\\
\text{Hexagone }(n=6) :&\quad \widehat{A} = \frac{4\times 180°}{6} = 120° ;\\
\text{Pentagone }(n=5) :&\quad \widehat{A} = \frac{3\times 180°}{5} = 108°.
\end{align*}
Un pavage sans trous ni chevauchements exige qu'autour de chaque sommet, les angles des pavés totalisent exactement $360°$ (un tour complet). Or :
\[ 6\times 60° = 360° \;;\qquad 4\times 90° = 360° \;;\qquad 3\times 120° = 360°. \]
Donc 6 triangles, 4 carrés ou 3 hexagones se raccordent parfaitement autour d'un sommet : le fournisseur dit vrai pour ses trois modèles. En revanche, pour le pentagone : $3\times 108° = 324° < 360°$ (il reste un trou de $36°$) et $4\times 108° = 432° > 360°$ (les pavés se chevauchent). Comme $360$ n'est pas un multiple de $108$, \textbf{le pentagone régulier ne peut pas paver la cour}.

\textbf{Étape 2 : aire d'un pavé de chaque type.}

Le côté commun est $c = \SI{30}{cm} = \SI{0,3}{m}$.

\emph{Triangle équilatéral.} La hauteur issue d'un sommet est aussi médiane : elle partage la base en deux segments de \SI{15}{cm}. Le théorème de Pythagore dans le demi-triangle rectangle donne :
\[ h^2 = 30^2 - 15^2 = 900 - 225 = 675, \quad \text{donc } h = \sqrt{675} = 15\sqrt{3} \approx \SI{25,98}{cm}. \]
\[ \mathcal{A}_T = \frac{30 \times 15\sqrt{3}}{2} = 225\sqrt{3} \approx \SI{389,7}{cm^2} \approx \SI{0,0390}{m^2}. \]

\emph{Carré.} $\mathcal{A}_C = 30^2 = \SI{900}{cm^2} = \SI{0,09}{m^2}$.

\emph{Hexagone régulier.} Il se décompose en 6 triangles équilatéraux de côté \SI{30}{cm} (reliés au centre) :
\[ \mathcal{A}_H = 6 \times 225\sqrt{3} = 1350\sqrt{3} \approx \SI{2338,3}{cm^2} \approx \SI{0,2338}{m^2}. \]

\textbf{Étape 3 : nombre de pavés et coût total.}

Aire de la cour : $\mathcal{A} = 6 \times 4,5 = \SI{27}{m^2}$.

\begin{center}
\begin{tabularx}{\linewidth}{|l|X|X|X|}\hline
\rowcolor{popblueL} & \textbf{Triangle} & \textbf{Carré} & \textbf{Hexagone} \\ \hline
Aire d'un pavé (m\textsuperscript{2}) & $0,0390$ & $0,09$ & $0,2338$ \\ \hline
Nombre théorique $27 \div \mathcal{A}$ & $693$ & $300$ & $116$ \\ \hline
Avec $10\,\%$ de marge & $763$ & $330$ & $128$ \\ \hline
Aire achetée (m\textsuperscript{2}) & $29,7$ & $29,7$ & $29,7$ \\ \hline
Coût total & $29,7\times\num{12000} = \num{356400}$ F & $29,7\times\num{10000} = \num{297000}$ F & $29,7\times\num{15000} = \num{445500}$ F \\ \hline
\end{tabularx}
\end{center}

Détail des arrondis : $27 \div 0,0390 \approx 692,3$ donc $693$ pavés ; $693 \times 1,1 = 762,3$ donc $763$ pavés. De même, $27 \div 0,2338 \approx 115,5$ donc $116$ pavés ; $116 \times 1,1 = 127,6$ donc $128$ pavés. Avec la marge de $10\,\%$, l'aire achetée est dans les trois cas $27 \times 1,1 = \SI{29,7}{m^2}$.

\textbf{Étape 4 : construction du motif hexagonal à l'échelle $1/50$.}

À l'échelle $1/50$, le côté de \SI{30}{cm} est représenté par $30 \div 50 = \SI{0,6}{cm} = \SI{6}{mm}$. On trace un cercle de centre $O$ et de rayon \SI{6}{mm} ; en reportant le rayon six fois de proche en proche sur le cercle, on obtient les six sommets de l'hexagone régulier central (propriété fondamentale : le côté de l'hexagone régulier est égal au rayon du cercle circonscrit). On construit ensuite chacun des 6 hexagones voisins en traçant, autour de chaque côté de l'hexagone central, un nouvel hexagone de même côté, toujours par reports du rayon au compas.

\textbf{Étape 5 : conseil rédigé (exemple de production attendue).}

« Maman Aïcha, après calculs, nous vous conseillons le \textbf{modèle carré}. C'est le moins cher : environ \num{297000} F, contre \num{356400} F pour le triangulaire et \num{445500} F pour l'hexagonal. De plus, les bords de votre cour étant droits, les pavés carrés exigent le moins de découpes, donc moins de pertes et une pose plus rapide. L'hexagone est le plus beau et le plus solide en emboîtement, mais il coûte $50\,\%$ de plus que le carré. Si votre budget le permet et que l'esthétique compte, choisissez-le ; sinon, le carré reste le choix raisonnable. »
\end{exoresolu}

\begin{attention}[Pièges fréquents dans cette activité]
\begin{itemize}
\item Ne pas oublier de convertir les cm en m (ou les cm\textsuperscript{2} en m\textsuperscript{2} : $1\,\text{m}^2 = \num{10000}\,\text{cm}^2$) avant de calculer les coûts.
\item La marge de $10\,\%$ s'applique \emph{après} le calcul du nombre théorique de pavés, et on arrondit toujours à l'entier supérieur (on ne peut pas acheter une fraction de pavé).
\item Dire « l'angle du pentagone est $108°$ » ne suffit pas : il faut justifier l'impossibilité du pavage en montrant qu'aucun multiple entier de $108°$ ne vaut $360°$.
\item Dans le calcul de la hauteur du triangle équilatéral, bien préciser que la hauteur est aussi médiane avant d'appliquer le théorème de Pythagore : c'est cette justification qui vaut des points au BEPC.
\end{itemize}
\end{attention}

\courssub{Bilan de l'activité}

Cette activité a permis de mettre en évidence que :

\begin{itemize}
\item les \textbf{angles intérieurs} des polygones réguliers ne sont pas qu'une formule : ils expliquent concrètement pourquoi seuls le triangle équilatéral, le carré et l'hexagone régulier pavent le plan (ce sont les trois seuls pavages réguliers possibles !) ;
\item la \textbf{décomposition} d'un polygone régulier en triangles est une stratégie puissante pour calculer des aires ;
\item les mathématiques de la leçon (angles, Pythagore, aires, proportionnalité) servent à prendre une \textbf{décision économique argumentée}, compétence utile bien au-delà de la classe ;
\item la \textbf{rédaction d'une démarche justifiée} est aussi importante que le résultat numérique : au BEPC comme dans la vie, une réponse non justifiée ne convainc personne.
\end{itemize}

\begin{savaistu}[Les abeilles, championnes du pavage]
Les alvéoles des ruches d'abeilles sont des hexagones réguliers ! Ce n'est pas un hasard : parmi les trois polygones qui pavent le plan, l'hexagone est celui qui, pour une aire donnée, utilise le \emph{moins de cire} (le plus petit périmètre). Les abeilles optimisent leur matériau depuis des millions d'années... Les pavés hexagonaux que l'on voit sur certaines rues de Yaoundé et de Douala utilisent la même idée d'emboîtement parfait.
\end{savaistu}

\newpage

\coursec{Méthodes et exercices résolus}

\courssub{Reconnaître et justifier qu'un polygone est régulier}

\begin{definition}[Polygone régulier]
Un \cle{polygone régulier} est un polygone dont :
\begin{itemize}
\item tous les \textbf{côtés} ont la même longueur ;
\item tous les \textbf{angles} ont la même mesure.
\end{itemize}
Les deux conditions sont \textbf{indispensables} et simultanées.
Un polygone régulier est toujours \cle{inscriptible} dans un cercle : il existe un cercle, appelé \cle{cercle circonscrit}, qui passe par tous ses sommets. Le centre $O$ de ce cercle est appelé le \cle{centre} du polygone régulier, et son rayon $R$ est le \cle{rayon} du polygone.
\end{definition}

\begin{exemplebox}[Les polygones réguliers à connaître en 3ème]
\begin{itemize}
\item Le \cle{triangle équilatéral} est le seul triangle régulier ($n = 3$ côtés) ;
\item le \cle{carré} est le seul quadrilatère régulier ($n = 4$ côtés) ;
\item l'\cle{hexagone régulier} ($n = 6$ côtés) ;
\item l'\cle{octogone régulier} ($n = 8$ côtés).
\end{itemize}
\end{exemplebox}

\begin{popfigure}
\begin{center}
\begin{tikzpicture}[scale=1.1]
\draw[popline] (0,0) circle (2);
\coordinate (O) at (0,0);
\coordinate (A) at (2,0);
\coordinate (B) at (1,1.732);
\coordinate (C) at (-1,1.732);
\coordinate (D) at (-2,0);
\coordinate (E) at (-1,-1.732);
\coordinate (F) at (1,-1.732);
\draw[thick,popblue] (A)--(B)--(C)--(D)--(E)--(F)--cycle;
\draw[poporange,thick] (O)--(A);
\draw[poporange,thick] (O)--(B);
\fill[popdark] (O) circle (1.5pt) node[below left] {$O$};
\foreach \p/\pos in {A/right,B/above right,C/above left,D/left,E/below left,F/below right}
  \fill[popblue] (\p) circle (1.5pt) node[\pos] {$\p$};
\node[poporange] at (0.95,0.28) {\small $R$};
\draw[popgreen,->] (0.55,0) arc (0:60:0.55);
\node[popgreen] at (0.85,0.62) {\small $60^\circ$};
\end{tikzpicture}
\end{center}
\legende{Hexagone régulier inscrit dans son cercle circonscrit : l'angle au centre $\widehat{AOB}$ vaut $\frac{360^\circ}{6} = 60^\circ$.}
\end{popfigure}

\begin{methode}[Montrer qu'un polygone est régulier]
Pour prouver qu'un polygone est régulier, on peut utiliser l'une des deux démarches suivantes :
\begin{enumerate}
\item \textbf{Par la définition} : montrer que tous les côtés ont la même longueur \textbf{et} que tous les angles ont la même mesure. Les deux conditions sont indispensables.
\item \textbf{Par le cercle circonscrit} : montrer que les sommets du polygone partagent un cercle en arcs de même longueur (angles au centre égaux à $\dfrac{360^\circ}{n}$).
\end{enumerate}
Réflexes à retenir :
\begin{itemize}
\item Un losange n'est pas forcément régulier (côtés égaux mais angles inégaux) ;
\item Un rectangle n'est pas forcément régulier (angles égaux mais côtés inégaux) ;
\item Le seul quadrilatère régulier est le \cle{carré}, le seul triangle régulier est le \cle{triangle équilatéral}.
\end{itemize}
\end{methode}

\begin{exoresolu}[Un polygone régulier ou non ?]
On considère un quadrilatère $ABCD$ inscrit dans un cercle de centre $O$ tel que :
$\widehat{AOB} = \widehat{BOC} = \widehat{COD} = \widehat{DOA} = 90^\circ$.
\begin{enumerate}
\item Montrer que $ABCD$ est un polygone régulier. Lequel ?
\item Le quadrilatère $EFGH$ est un losange dont un angle mesure $70^\circ$. Est-il régulier ? Justifier.
\end{enumerate}
\tcblower
\textbf{1.} Les quatre angles au centre $\widehat{AOB}$, $\widehat{BOC}$, $\widehat{COD}$, $\widehat{DOA}$ sont égaux à $90^\circ = \dfrac{360^\circ}{4}$. Les arcs $\overset{\frown}{AB}$, $\overset{\frown}{BC}$, $\overset{\frown}{CD}$, $\overset{\frown}{DA}$ ont donc la même longueur, et les cordes correspondantes aussi :
\[ AB = BC = CD = DA. \]
De plus, chaque côté est vu du centre $O$ sous l'angle $90^\circ$ : les triangles $OAB$, $OBC$, $OCD$, $ODA$ sont isocèles en $O$ (car $OA = OB = OC = OD$, rayon) et superposables. Leurs angles à la base valent $\dfrac{180^\circ - 90^\circ}{2} = 45^\circ$. Chaque angle du quadrilatère vaut donc $45^\circ + 45^\circ = 90^\circ$.
Le quadrilatère $ABCD$ a ses quatre côtés égaux et ses quatre angles égaux : c'est un polygone régulier à $4$ côtés, c'est-à-dire un \textbf{carré}.

\textbf{2.} Le losange $EFGH$ a bien ses quatre côtés égaux, mais un de ses angles mesure $70^\circ$, donc l'angle opposé mesure $70^\circ$ et les deux autres mesurent $180^\circ - 70^\circ = 110^\circ$ (deux angles consécutifs d'un parallélogramme sont supplémentaires). Les angles ne sont pas tous égaux : $EFGH$ \textbf{n'est pas régulier}.

\textbf{Conclusion :} $ABCD$ est un carré (polygone régulier) ; $EFGH$ est un losange non régulier, car l'égalité des côtés ne suffit pas.
\end{exoresolu}

\courssub{Calculer les angles d'un polygone régulier}

\begin{propriete}[Angles d'un polygone régulier]
Soit un polygone régulier à $n$ côtés, de centre $O$, et $A$, $B$, $C$ trois sommets consécutifs.
\begin{itemize}
\item \textbf{Angle au centre} : $\widehat{AOB} = \dfrac{360^\circ}{n}$ (le tour complet est partagé en $n$ angles égaux).
\item \textbf{Angle intérieur} : $\widehat{ABC} = 180^\circ - \dfrac{360^\circ}{n}$.
\item \textbf{Somme des angles} d'un polygone à $n$ côtés : $(n-2)\times 180^\circ$.
\end{itemize}
\end{propriete}

\begin{experience}[D'où vient la formule de l'angle intérieur ?]
Considérons le triangle $OAB$ formé par le centre et deux sommets consécutifs.
\begin{enumerate}
\item $OA = OB = R$ (rayons du cercle circonscrit), donc le triangle $OAB$ est \textbf{isocèle en $O$} : ses angles à la base sont égaux.
\item L'angle au sommet est $\widehat{AOB} = \dfrac{360^\circ}{n}$. La somme des angles d'un triangle valant $180^\circ$, chaque angle à la base vaut :
\[ \widehat{OAB} = \widehat{OBA} = \frac{180^\circ - \frac{360^\circ}{n}}{2} = 90^\circ - \frac{180^\circ}{n}. \]
\item L'angle intérieur $\widehat{ABC}$ du polygone est formé de deux tels angles à la base (ceux des triangles $OAB$ et $OBC$) :
\[ \widehat{ABC} = 2\left(90^\circ - \frac{180^\circ}{n}\right) = 180^\circ - \frac{360^\circ}{n}. \]
\item \textbf{Vérification} : la somme des $n$ angles intérieurs vaut $n\left(180^\circ - \frac{360^\circ}{n}\right) = 180n - 360 = (n-2)\times 180^\circ$ : on retrouve bien la somme des angles d'un polygone à $n$ côtés.
\end{enumerate}
\end{experience}

\begin{aretenir}[Valeurs à connaître par cœur]
\begin{center}
\begin{tabular}{|l|c|c|c|}
\hline
\rowcolor{popblueL} Polygone & $n$ & Angle au centre & Angle intérieur \\
\hline
Triangle équilatéral & $3$ & $120^\circ$ & $60^\circ$ \\
\hline
Carré & $4$ & $90^\circ$ & $90^\circ$ \\
\hline
Hexagone régulier & $6$ & $60^\circ$ & $120^\circ$ \\
\hline
Octogone régulier & $8$ & $45^\circ$ & $135^\circ$ \\
\hline
\end{tabular}
\end{center}
Remarque : l'angle au centre et l'angle intérieur sont \textbf{supplémentaires} (leur somme vaut $180^\circ$).
\end{aretenir}

\begin{exoresolu}[Angles d'un octogone régulier]
$ABCDEFGH$ est un octogone régulier de centre $O$.
\begin{enumerate}
\item Calculer la mesure de l'angle au centre $\widehat{AOB}$.
\item En déduire la mesure de l'angle intérieur $\widehat{ABC}$.
\item Vérifier avec la somme des angles d'un polygone à $8$ côtés.
\end{enumerate}
\tcblower
\textbf{1.} L'octogone régulier a $n = 8$ côtés, donc :
\[ \widehat{AOB} = \frac{360^\circ}{8} = 45^\circ. \]

\textbf{2.} Le triangle $OAB$ est isocèle en $O$ car $OA = OB$ (rayons du cercle circonscrit). Ses angles à la base sont égaux :
\[ \widehat{OAB} = \widehat{OBA} = \frac{180^\circ - 45^\circ}{2} = \frac{135^\circ}{2} = 67{,}5^\circ. \]
L'angle intérieur de l'octogone est formé de deux tels angles :
\[ \widehat{ABC} = \widehat{ABO} + \widehat{OBC} = 67{,}5^\circ + 67{,}5^\circ = 135^\circ. \]
On retrouve la formule : $180^\circ - \dfrac{360^\circ}{8} = 180^\circ - 45^\circ = 135^\circ$.

\textbf{3.} La somme des angles d'un polygone à $8$ côtés vaut :
\[ (8-2)\times 180^\circ = 6 \times 180^\circ = 1080^\circ. \]
Or $8 \times 135^\circ = 1080^\circ$ : la vérification est concluante.

\textbf{Conclusion :} $\widehat{AOB} = 45^\circ$ et $\widehat{ABC} = 135^\circ$.
\end{exoresolu}

\courssub{Exploiter le triangle équilatéral et le carré}

\begin{propriete}[Triangle équilatéral de côté $a$]
\begin{itemize}
\item Chaque hauteur est aussi médiane et médiatrice : elle partage le côté opposé en deux segments de longueur $\dfrac{a}{2}$.
\item \textbf{Hauteur} : $h = \dfrac{a\sqrt{3}}{2}$.
\item \textbf{Aire} : $\mathcal{A} = \dfrac{a^2\sqrt{3}}{4}$.
\end{itemize}
\emph{Justification de la hauteur (exigible) :} dans le triangle rectangle formé par un sommet, le pied $H$ de la hauteur et un sommet adjacent, le théorème de Pythagore donne :
\[ h^2 = a^2 - \left(\frac{a}{2}\right)^2 = a^2 - \frac{a^2}{4} = \frac{3a^2}{4}, \quad \text{donc} \quad h = \frac{a\sqrt{3}}{2}. \]
\end{propriete}

\begin{propriete}[Carré de côté $a$]
\begin{itemize}
\item \textbf{Diagonale} : $d = a\sqrt{2}$ (théorème de Pythagore dans le triangle rectangle formé par deux côtés consécutifs et une diagonale : $d^2 = a^2 + a^2 = 2a^2$).
\item \textbf{Périmètre} : $\mathcal{P} = 4a$ ; \textbf{aire} : $\mathcal{A} = a^2$.
\item Les diagonales se coupent en leur milieu, sont perpendiculaires et de même longueur ; leur point d'intersection est le centre du cercle circonscrit, de rayon $R = \dfrac{a\sqrt{2}}{2}$.
\end{itemize}
\end{propriete}

\begin{exoresolu}[Terrain triangulaire à Mbalmayo]
Un terrain a la forme d'un triangle équilatéral $ABC$ de côté $a = 60$ m. On veut installer une clôture le long de la hauteur issue de $A$ pour partager le terrain en deux parcelles.
\begin{enumerate}
\item Calculer la longueur exacte de cette clôture, puis donner une valeur approchée au décimètre près.
\item Calculer l'aire totale du terrain (valeur approchée au m$^2$ près).
\end{enumerate}
On donne $\sqrt{3} \approx 1{,}732$.
\tcblower
\textbf{1.} Soit $H$ le pied de la hauteur issue de $A$. Dans un triangle équilatéral, la hauteur est aussi médiane, donc $H$ est le milieu de $[BC]$ et $BH = \dfrac{60}{2} = 30$ m.
Le triangle $ABH$ est rectangle en $H$. D'après le théorème de Pythagore :
\begin{align*}
AH^2 + BH^2 &= AB^2 \\
AH^2 &= 60^2 - 30^2 = 3600 - 900 = 2700 \\
AH &= \sqrt{2700} = \sqrt{900 \times 3} = 30\sqrt{3} \text{ m}.
\end{align*}
Valeur approchée : $AH \approx 30 \times 1{,}732 = 51{,}96$, soit $AH \approx 52{,}0$ m au décimètre près.

\textbf{2.} L'aire du terrain est :
\[ \mathcal{A} = \frac{BC \times AH}{2} = \frac{60 \times 30\sqrt{3}}{2} = 900\sqrt{3} \approx 900 \times 1{,}732 = 1558{,}8. \]
L'aire du terrain est d'environ $\SI{1559}{m^2}$.

\textbf{Conclusion :} la clôture mesure $30\sqrt{3}$ m $\approx 52{,}0$ m et le terrain a une aire d'environ $\SI{1559}{m^2}$.
\end{exoresolu}

\courssub{Exploiter l'hexagone régulier}

\begin{propriete}[La propriété-clé de l'hexagone régulier]
Dans un hexagone régulier $ABCDEF$ de centre $O$ :
\begin{itemize}
\item l'angle au centre vaut $\dfrac{360^\circ}{6} = 60^\circ$ ;
\item le triangle $OAB$ est isocèle en $O$ avec un angle au sommet de $60^\circ$ : ses angles à la base valent $\dfrac{180^\circ - 60^\circ}{2} = 60^\circ$, il est donc \textbf{équilatéral} ;
\item conséquence fondamentale : \cle{le côté de l'hexagone régulier est égal au rayon du cercle circonscrit} : $AB = R$ ;
\item l'hexagone régulier se décompose en $6$ triangles équilatéraux de côté $R$.
\end{itemize}
Réflexe : pour construire un hexagone régulier, on reporte le rayon six fois sur le cercle.
\end{propriete}

\begin{exoresolu}[Côté et aire d'un hexagone régulier]
Un hexagone régulier $ABCDEF$ est inscrit dans un cercle de centre $O$ et de rayon $R = 10$ cm.
\begin{enumerate}
\item Justifier que le triangle $OAB$ est équilatéral et donner la longueur du côté de l'hexagone.
\item Calculer le périmètre de l'hexagone.
\item Calculer l'aire de l'hexagone (on donnera la valeur exacte, puis une valeur approchée au cm$^2$ près avec $\sqrt{3} \approx 1{,}732$).
\end{enumerate}
\tcblower
\textbf{1.} Dans le cercle de centre $O$, l'angle au centre interceptant le côté $[AB]$ vaut :
\[ \widehat{AOB} = \frac{360^\circ}{6} = 60^\circ. \]
Le triangle $OAB$ est isocèle en $O$ car $OA = OB = R$. Ses angles à la base sont donc égaux :
\[ \widehat{OAB} = \widehat{OBA} = \frac{180^\circ - 60^\circ}{2} = 60^\circ. \]
Les trois angles du triangle $OAB$ valent $60^\circ$ : le triangle $OAB$ est \textbf{équilatéral}, donc $AB = OA = OB = 10$ cm.
Le côté de l'hexagone mesure $10$ cm, égal au rayon.

\textbf{2.} Le périmètre est :
\[ \mathcal{P} = 6 \times 10 = 60 \text{ cm}. \]

\textbf{3.} L'hexagone est formé de $6$ triangles équilatéraux de côté $10$ cm. La hauteur d'un tel triangle vaut :
\[ h = \frac{10\sqrt{3}}{2} = 5\sqrt{3} \text{ cm}. \]
Aire d'un triangle : $\mathcal{A}_1 = \dfrac{10 \times 5\sqrt{3}}{2} = 25\sqrt{3}$ cm$^2$.
Aire de l'hexagone :
\[ \mathcal{A} = 6 \times 25\sqrt{3} = 150\sqrt{3} \approx 150 \times 1{,}732 = 259{,}8. \]

\textbf{Conclusion :} le côté mesure $10$ cm, le périmètre $60$ cm, et l'aire vaut $150\sqrt{3}$ cm$^2 \approx 260$ cm$^2$.
\end{exoresolu}

\courssub{Apothème, périmètre et aire d'un polygone régulier}

\begin{definition}[Apothème]
L'\cle{apothème} d'un polygone régulier est la distance du centre $O$ à un côté. Si $H$ est le milieu du côté $[AB]$, alors l'apothème est la longueur $OH$ ; c'est aussi la hauteur issue de $O$ dans le triangle isocèle $OAB$.
\end{definition}

\begin{propriete}[Aire d'un polygone régulier]
Un polygone régulier à $n$ côtés, de côté $c$ et d'apothème $a$, se décompose en $n$ triangles isocèles identiques de base $c$ et de hauteur $a$. Son aire est :
\[ \mathcal{A} = n \times \frac{c \times a}{2} = \frac{\mathcal{P} \times a}{2}, \]
où $\mathcal{P} = n \times c$ est le périmètre du polygone.
\end{propriete}

\begin{methode}[Calculer l'apothème avec Pythagore]
\begin{enumerate}
\item Justifier l'angle droit : dans le triangle isocèle $OAB$ ($OA = OB$), la médiane $(OH)$ issue du sommet principal $O$ est aussi hauteur, donc $(OH) \perp (AB)$.
\item Le triangle $OHA$ est rectangle en $H$, avec $OA = R$ (rayon), $HA = \dfrac{c}{2}$ (moitié du côté) et $OH = a$ (apothème).
\item Le théorème de Pythagore donne : $a^2 + \left(\dfrac{c}{2}\right)^2 = R^2$, d'où $a = \sqrt{R^2 - \dfrac{c^2}{4}}$.
\end{enumerate}
\end{methode}

\begin{exoresolu}[Apothème et aire d'un octogone régulier]
Un panneau de signalisation a la forme d'un octogone régulier de côté $c = 30$ cm, inscrit dans un cercle de centre $O$ et de rayon $R \approx 39{,}2$ cm. Soit $H$ le milieu d'un côté $[AB]$.
\begin{enumerate}
\item Calculer l'apothème $OH$ de l'octogone (arrondir au dixième).
\item En déduire l'aire du panneau (arrondir au cm$^2$).
\end{enumerate}
\tcblower
\textbf{1.} La droite $(OH)$ joint le centre au milieu du côté $[AB]$ : dans le triangle isocèle $OAB$ (car $OA = OB = R$), c'est la médiane issue du sommet principal $O$, donc aussi la hauteur : $(OH) \perp (AB)$. Le triangle $OHA$ est rectangle en $H$, avec :
\[ OA = R \approx 39{,}2 \text{ cm} \quad \text{et} \quad HA = \frac{c}{2} = 15 \text{ cm}. \]
D'après le théorème de Pythagore :
\begin{align*}
OH^2 + HA^2 &= OA^2 \\
OH^2 &= 39{,}2^2 - 15^2 = 1536{,}64 - 225 = 1311{,}64 \\
OH &= \sqrt{1311{,}64} \approx 36{,}2 \text{ cm}.
\end{align*}
L'apothème mesure environ $36{,}2$ cm.

\textbf{2.} Le périmètre de l'octogone est $\mathcal{P} = 8 \times 30 = 240$ cm. L'aire est :
\[ \mathcal{A} = \frac{\mathcal{P} \times a}{2} \approx \frac{240 \times 36{,}2}{2} = 120 \times 36{,}2 = 4344 \text{ cm}^2. \]

\textbf{Conclusion :} l'apothème mesure environ $36{,}2$ cm et le panneau a une aire d'environ $\SI{4344}{cm^2}$.
\end{exoresolu}

\courssub{Construire un polygone régulier}

\begin{methode}[Constructions à la règle et au compas]
\textbf{Triangle équilatéral :}
\begin{enumerate}
\item Tracer un segment $[AB]$.
\item Tracer les arcs de cercle de centres $A$ et $B$ et de rayon $AB$ ; leur intersection est le troisième sommet $C$.
\end{enumerate}
\textbf{Carré :}
\begin{enumerate}
\item Tracer un cercle de centre $O$ et deux diamètres perpendiculaires $[AC]$ et $[BD]$.
\item Joindre les quatre extrémités $A$, $B$, $C$, $D$ : $ABCD$ est un carré (quatre angles au centre de $90^\circ$).
\end{enumerate}
\textbf{Hexagone régulier :}
\begin{enumerate}
\item Tracer un cercle de centre $O$ et de rayon $R$.
\item Choisir un point $A$ sur le cercle ; reporter le rayon $R$ au compas de proche en proche : on obtient $B$, $C$, $D$, $E$, $F$.
\item Joindre les six points consécutifs.
\end{enumerate}
\textbf{Octogone régulier :}
\begin{enumerate}
\item Construire un carré $ABCD$ inscrit dans le cercle.
\item Tracer les bissectrices des angles au centre $\widehat{AOB}$, $\widehat{BOC}$, \ldots{} (ou les médiatrices des côtés) : elles coupent le cercle en quatre nouveaux points.
\item Joindre les huit points dans l'ordre : chaque angle au centre vaut alors $\dfrac{90^\circ}{2} = 45^\circ$.
\end{enumerate}
\end{methode}

\begin{exoresolu}[Construction et justification]
\begin{enumerate}
\item Construire un hexagone régulier $ABCDEF$ inscrit dans un cercle de centre $O$ et de rayon $4$ cm.
\item Justifier que la construction donne bien un polygone régulier.
\end{enumerate}
\tcblower
\textbf{1.} On trace le cercle $\mathcal{C}$ de centre $O$ et de rayon $4$ cm. On place un point $A$ sur $\mathcal{C}$. L'écartement du compas étant réglé sur $4$ cm, on reporte cette longueur sur le cercle à partir de $A$ : on obtient successivement $B$, $C$, $D$, $E$, $F$. On trace les segments $[AB]$, $[BC]$, $[CD]$, $[DE]$, $[EF]$, $[FA]$.

\textbf{2.} Par construction, $OA = OB = AB = 4$ cm : le triangle $OAB$ est équilatéral, donc $\widehat{AOB} = 60^\circ$. Il en est de même pour les six triangles $OAB$, $OBC$, \ldots, $OFA$. Les six angles au centre valent chacun $60^\circ$, et $6 \times 60^\circ = 360^\circ$ : le point $F$ retombe exactement à distance $R$ du point de départ, la figure se referme.
Les six côtés sont égaux ($4$ cm) et les six angles au centre sont égaux, donc les arcs et les angles intérieurs sont égaux : chaque angle intérieur vaut $180^\circ - 60^\circ = 120^\circ$.
Le polygone $ABCDEF$ a tous ses côtés et tous ses angles égaux : c'est bien un \textbf{hexagone régulier}.
\end{exoresolu}

\courssub{Problèmes de synthèse et situations concrètes}

\begin{methode}[Résoudre un problème concret avec les polygones réguliers]
\begin{enumerate}
\item \textbf{Modéliser} : identifier le polygone régulier dans la situation (dalle, table, panneau, pavage) et noter les données : nombre de côtés $n$, côté $c$, rayon $R$.
\item \textbf{Choisir l'outil adapté} : angle au centre $\frac{360^\circ}{n}$, Pythagore dans le triangle rectangle $OHA$, aire $\frac{\mathcal{P}\times a}{2}$.
\item \textbf{Calculer} en gardant les valeurs exactes le plus longtemps possible, puis arrondir seulement à la fin, à la précision demandée.
\item \textbf{Conclure} par une phrase répondant à la question posée, avec l'unité.
\end{enumerate}
\end{methode}

\begin{exoresolu}[Pavage d'une cour à Yaoundé]
Une dalle de jardin a la forme d'un hexagone régulier de côté $40$ cm.
\begin{enumerate}
\item Montrer que trois dalles placées autour d'un même point se rejoignent parfaitement, sans laisser de trou.
\item Calculer l'aire d'une dalle (valeur approchée au cm$^2$ près, $\sqrt{3} \approx 1{,}732$).
\item Combien de dalles faut-il environ pour paver une cour de $10$ m$^2$ ?
\end{enumerate}
\tcblower
\textbf{1.} L'angle intérieur d'un hexagone régulier vaut :
\[ 180^\circ - \frac{360^\circ}{6} = 180^\circ - 60^\circ = 120^\circ. \]
Trois angles de $120^\circ$ autour d'un même point donnent $3 \times 120^\circ = 360^\circ$, soit un tour complet : les trois dalles se rejoignent exactement sans trou ni recouvrement. C'est pourquoi l'hexagone régulier permet un pavage parfait (comme les alvéoles des abeilles).

\textbf{2.} L'hexagone régulier de côté $c = 40$ cm se décompose en $6$ triangles équilatéraux de côté $40$ cm. La hauteur de chacun est :
\[ h = \frac{40\sqrt{3}}{2} = 20\sqrt{3} \text{ cm}. \]
Aire de l'hexagone :
\[ \mathcal{A} = 6 \times \frac{40 \times 20\sqrt{3}}{2} = 6 \times 400\sqrt{3} = 2400\sqrt{3} \approx 2400 \times 1{,}732 = 4156{,}8 \text{ cm}^2. \]
Une dalle a une aire d'environ $\SI{4157}{cm^2}$.

\textbf{3.} Convertissons : $10$ m$^2 = 100\,000$ cm$^2$. Le nombre de dalles est :
\[ \frac{100\,000}{4156{,}8} \approx 24{,}06. \]
Il faut donc prévoir environ \textbf{25 dalles} (en arrondissant à l'entier supérieur pour tenir compte des découpes).

\textbf{Conclusion :} les dalles hexagonales pavent parfaitement la cour ; chaque dalle couvre environ $\SI{4157}{cm^2}$ et il faut environ $25$ dalles pour $10$ m$^2$.
\end{exoresolu}

\begin{savaistu}[Pourquoi les abeilles construisent-elles des hexagones ?]
Les alvéoles des ruches sont des hexagones réguliers. Parmi les polygones réguliers, seuls le triangle équilatéral ($6 \times 60^\circ = 360^\circ$), le carré ($4 \times 90^\circ = 360^\circ$) et l'hexagone régulier ($3 \times 120^\circ = 360^\circ$) permettent de paver le plan sans trou. À périmètre égal, c'est l'hexagone qui a la plus grande aire : les abeilles dépensent ainsi le moins de cire possible pour stocker le maximum de miel. Les mathématiciens n'ont démontré rigoureusement ce résultat (le « théorème du nid d'abeille ») qu'en 1999 !
\end{savaistu}

\begin{attention}[Les 5 erreurs qui coûtent des points au BEPC]
\begin{enumerate}
\item \textbf{Croire que les côtés égaux suffisent.} Un losange a quatre côtés égaux mais n'est pas régulier : il faut vérifier \textbf{à la fois} l'égalité des côtés \textbf{et} celle des angles. Réciproquement, un rectangle a des angles égaux mais n'est pas régulier.
\item \textbf{Confondre angle au centre et angle intérieur.} Pour l'hexagone, l'angle au centre vaut $60^\circ$ et l'angle intérieur $120^\circ$ : ne pas les intervertir. Formules : centre $= \dfrac{360^\circ}{n}$, intérieur $= 180^\circ - \dfrac{360^\circ}{n}$.
\item \textbf{Oublier la justification de l'angle droit avant Pythagore.} Avant d'appliquer le théorème de Pythagore dans le triangle $OHA$, il faut justifier que $(OH) \perp (AB)$ : dans un triangle isocèle, la médiane issue du sommet principal est aussi hauteur. Sans cette phrase, le calcul n'est pas valide.
\item \textbf{Se tromper dans la hauteur du triangle équilatéral.} La hauteur est $h = \dfrac{a\sqrt{3}}{2}$, pas $\dfrac{a}{2}$ ni $a\sqrt{3}$. Refaire le calcul par Pythagore en cas de doute : $h^2 = a^2 - \left(\frac{a}{2}\right)^2$.
\item \textbf{Arrondir trop tôt et oublier les unités.} Garde les valeurs exactes (avec $\sqrt{3}$, $\sqrt{2}$) pendant les calculs et n'arrondis qu'à la fin, à la précision demandée. Termine toujours par une phrase de conclusion avec l'unité (cm, cm$^2$, m$^2$) : une réponse sans unité ni conclusion perd des points de présentation.
\end{enumerate}
\end{attention}

\newpage

\coursec{Exercices}

\courssub{Je vérifie mes connaissances}

\begin{exercice}[QCM]
Pour chaque question, une seule réponse est exacte. Recopie la lettre de la bonne réponse.
\begin{enumerate}
\item Un polygone régulier est un polygone dont :
\begin{enumerate}
\item[a)] tous les côtés sont égaux ;
\item[b)] tous les angles sont égaux ;
\item[c)] tous les côtés sont égaux \textbf{et} tous les angles sont égaux.
\end{enumerate}
\item L'angle au centre d'un hexagone régulier mesure :
\begin{enumerate}
\item[a)] $45^{\circ}$ ;\qquad b) $60^{\circ}$ ;\qquad c) $120^{\circ}$.
\end{enumerate}
\item L'angle intérieur d'un octogone régulier mesure :
\begin{enumerate}
\item[a)] $135^{\circ}$ ;\qquad b) $120^{\circ}$ ;\qquad c) $108^{\circ}$.
\end{enumerate}
\item L'apothème d'un polygone régulier est :
\begin{enumerate}
\item[a)] la distance du centre à un sommet ;
\item[b)] la distance du centre au milieu d'un côté ;
\item[c)] la longueur d'un côté.
\end{enumerate}
\end{enumerate}
\end{exercice}

\begin{exercice}[Vrai ou faux]
Réponds par VRAI ou FAUX en justifiant chaque réponse par une propriété du cours ou un contre-exemple.
\begin{enumerate}
\item Tout triangle équilatéral est un polygone régulier.
\item Tout rectangle est un polygone régulier.
\item Dans un hexagone régulier inscrit dans un cercle, le côté est égal au rayon du cercle.
\item La somme des angles intérieurs d'un pentagone convexe vaut $540^{\circ}$.
\end{enumerate}
\end{exercice}

\begin{exercice}[Définitions]
\begin{enumerate}
\item Donne la définition d'un polygone régulier.
\item Qu'appelle-t-on \cle{angle au centre} d'un polygone régulier ? Donne la formule qui permet de le calculer en fonction du nombre $n$ de côtés.
\item Qu'appelle-t-on \cle{apothème} d'un polygone régulier ? Illustre ta réponse par une figure à main levée pour un carré.
\end{enumerate}
\end{exercice}

\begin{exercice}[Calculs directs]
\begin{enumerate}
\item Calcule l'angle au centre d'un polygone régulier à $10$ côtés.
\item Calcule la somme des angles intérieurs d'un polygone convexe à $8$ côtés, puis l'angle intérieur d'un octogone régulier.
\item Un carré a pour côté $\SI{6}{cm}$. Calcule son périmètre, son aire et la longueur de sa diagonale (donner la valeur exacte).
\item Un triangle équilatéral a pour côté $\SI{8}{cm}$. Calcule sa hauteur (valeur exacte) et son aire (valeur exacte puis arrondie au $\text{mm}^2$).
\end{enumerate}
\end{exercice}

\courssub{Je m'entraîne}

\begin{exercice}[Nombre de côtés]
Un polygone régulier a un angle intérieur de $140^{\circ}$.
\begin{enumerate}
\item Calcule la mesure de son angle au centre.
\item En déduire le nombre de côtés de ce polygone.
\item Vérifie le résultat en calculant la somme des angles intérieurs de ce polygone.
\end{enumerate}
\end{exercice}

\begin{exercice}[Hexagone inscrit]
Un hexagone régulier $ABCDEF$ est inscrit dans un cercle de centre $O$ et de rayon $\SI{8}{cm}$.
\begin{enumerate}
\item Justifie que le triangle $OAB$ est équilatéral, puis déduis-en la longueur du côté de l'hexagone.
\item Calcule le périmètre de l'hexagone.
\item Calcule l'apothème de l'hexagone (valeur exacte).
\item Calcule l'aire de l'hexagone (valeur exacte, puis arrondie au $\text{cm}^2$).
\end{enumerate}
\end{exercice}

\begin{exercice}[Construction d'un octogone]
\begin{enumerate}
\item Tracez un cercle de centre $O$ et de rayon $\SI{5}{cm}$. Construis à la règle et au compas un octogone régulier inscrit dans ce cercle (on pourra construire deux diamètres perpendiculaires, puis les bissectrices des angles obtenus).
\item Calcule la mesure de l'angle au centre de cet octogone.
\item Calcule la mesure de chaque angle intérieur de cet octogone.
\end{enumerate}
\end{exercice}

\begin{exercice}[Carré et diagonale]
Une cour a la forme d'un carré dont la diagonale mesure $\SI{20}{m}$.
\begin{enumerate}
\item Calcule la longueur du côté de ce carré (valeur exacte).
\item Calcule le périmètre et l'aire de la cour.
\item On veut entourer la cour d'une clôture coûtant \SI{2500}{F} le mètre. Calcule le coût total de la clôture (on prendra $\sqrt{2} \approx \num{1,414}$).
\end{enumerate}
\end{exercice}

\courssub{Je me prépare au Probatoire}

\begin{exercice}[La table octogonale]
Un menuisier de Yaoundé découpe une table octogonale régulière dans un disque de contreplaqué de diamètre \SI{1,20}{m}. L'octogone est inscrit dans le disque.

On donne : $\sin 22,5^{\circ} \approx \num{0,383}$ et $\cos 22,5^{\circ} \approx \num{0,924}$.
\begin{enumerate}
\item Calcule le rayon $R$ du disque.
\item Calcule la mesure de l'angle au centre de l'octogone, puis la mesure de chaque angle intérieur.
\item On admet que le côté de l'octogone est donné par $c = 2R\sin 22,5^{\circ}$. Calcule le côté $c$ et le périmètre de la table (arrondir au $\text{cm}$).
\item On admet que l'apothème est donné par $a = R\cos 22,5^{\circ}$. Calcule l'apothème $a$.
\item Calcule l'aire de la table (arrondir au $\text{dm}^2$).
\item Calcule l'aire du disque de contreplaqué, puis l'aire des chutes perdues (arrondir au $\text{dm}^2$).
\item Le contreplaqué coûte \SI{4500}{F/m^2}. Calcule le coût des chutes perdues (arrondir au franc).
\end{enumerate}
\end{exercice}

\begin{exercice}[Le terrain de M. Ngoa]
M. Ngoa possède à Obala un terrain formé d'un carré $ABCD$ de côté $\SI{30}{m}$, prolongé sur le côté $[CD]$ par un triangle équilatéral $CDE$ extérieur au carré.
\begin{enumerate}
\item Fais une figure à main levée du terrain et code-la.
\item Calcule le périmètre du terrain (attention : le côté $[CD]$ est intérieur au terrain).
\item Calcule la hauteur du triangle équilatéral $CDE$ (valeur exacte).
\item Calcule l'aire totale du terrain (valeur exacte, puis arrondie au $\text{m}^2$ ; on prendra $\sqrt{3} \approx \num{1,732}$).
\item M. Ngoa veut clôturer son terrain avec un grillage coûtant \SI{2500}{F} le mètre, en prévoyant un portail de $\SI{4}{m}$ qui ne sera pas grillagé. Calcule la dépense prévue.
\item M. Ngoa veut semer de l'herbe sur tout le terrain. Un sac de semence couvre $\SI{50}{m^2}$ et coûte \SI{6000}{F}. Combien de sacs doit-il acheter ? Quelle est la dépense ?
\end{enumerate}
\end{exercice}

\begin{exercice}[Somme des angles et pavage]
\textbf{Partie A : somme des angles intérieurs.}

Soit un polygone convexe à $n$ côtés ($n \geq 3$).
\begin{enumerate}
\item En choisissant un sommet du polygone et en traçant toutes les diagonales issues de ce sommet, explique pourquoi le polygone est partagé en $(n-2)$ triangles.
\item Démontre que la somme des angles intérieurs d'un polygone convexe à $n$ côtés vaut $(n-2) \times 180^{\circ}$.
\item En déduire la mesure de l'angle intérieur d'un polygone \emph{régulier} à $n$ côtés.
\end{enumerate}

\textbf{Partie B : pavage du sol.}

Un carreleur de Douala veut recouvrir un sol avec des carreaux identiques en forme de polygone régulier, joints sans trou autour de chaque sommet.
\begin{enumerate}
\item Calcule l'angle intérieur d'un hexagone régulier.
\item Justifie que trois hexagones réguliers se rejoignent exactement en un point sans laisser de trou.
\item Le carreleur peut-il paver le sol avec des octogones réguliers identiques seuls ? Justifie par un calcul.
\item Cite un autre polygone régulier (autre que l'hexagone) permettant de paver le sol, et justifie ta réponse.
\end{enumerate}
\end{exercice}

\newpage

\coursec{Corrigés des exercices}

\begin{corrige}[Exercice 1 — QCM]
\begin{enumerate}
\item \textbf{c)} Un polygone régulier est un polygone \textbf{convexe} dont tous les côtés sont égaux \textbf{et} tous les angles sont égaux. Les deux conditions sont nécessaires (un losange a côtés égaux mais angles inégaux ; un rectangle a angles égaux mais côtés inégaux).
\item \textbf{b)} L'angle au centre d'un polygone régulier à $n$ côtés vaut $\frac{360^{\circ}}{n}$. Pour un hexagone ($n=6$) : $\frac{360^{\circ}}{6}=60^{\circ}$.
\item \textbf{a)} L'angle intérieur d'un polygone régulier à $n$ côtés vaut $\frac{(n-2)\times 180^{\circ}}{n}$. Pour un octogone ($n=8$) : $\frac{6\times 180^{\circ}}{8}=135^{\circ}$.
\item \textbf{b)} L'apothème est la distance du centre au milieu d'un côté (c'est la hauteur des triangles isocèles formés par les rayons). La distance du centre à un sommet est le \cle{rayon} du cercle circonscrit.
\end{enumerate}
\end{corrige}

\begin{corrige}[Exercice 2 — Vrai ou faux]
\begin{enumerate}
\item \textbf{VRAI.} Un triangle équilatéral a ses trois côtés égaux et ses trois angles égaux (chacun $60^{\circ}$). C'est le polygone régulier à $3$ côtés.
\item \textbf{FAUX.} Un rectangle a quatre angles droits (égaux), mais ses côtés ne sont pas nécessairement tous égaux (seul le carré, cas particulier de rectangle, est un polygone régulier).
\item \textbf{VRAI.} Dans un hexagone régulier inscrit, l'angle au centre vaut $60^{\circ}$. Le triangle $OAB$ est isocèle en $O$ ($OA=OB=R$) avec un angle au sommet de $60^{\circ}$, donc il est équilatéral. Ainsi $AB=OA=R$.
\item \textbf{VRAI.} La somme des angles intérieurs d'un polygone convexe à $n$ côtés vaut $(n-2)\times 180^{\circ}$. Pour $n=5$ : $(5-2)\times 180^{\circ}=540^{\circ}$.
\end{enumerate}
\end{corrige}

\begin{corrige}[Exercice 3 — Définitions]
\begin{enumerate}
\item Un \cle{polygone régulier} est un polygone \textbf{convexe} qui a tous ses côtés de même longueur et tous ses angles de même mesure.
\item L'\cle{angle au centre} d'un polygone régulier est l'angle formé par les deux rayons joignant le centre du polygone à deux sommets consécutifs. Sa mesure est $\dfrac{360^{\circ}}{n}$, où $n$ est le nombre de côtés.
\item L'\cle{apothème} d'un polygone régulier est la distance du centre du polygone au milieu d'un côté. C'est la hauteur (issue du centre) des triangles isocèles formés par deux rayons consécutifs et le côté correspondant.
\begin{popfigure}
\begin{tikzpicture}[scale=1.2]
\coordinate (O) at (0,0);
\coordinate (A) at (-1.5,-1.5);
\coordinate (B) at (1.5,-1.5);
\coordinate (C) at (1.5,1.5);
\coordinate (D) at (-1.5,1.5);
\coordinate (M) at (0,-1.5); % milieu de AB
\draw[popblue, thick] (A)--(B)--(C)--(D)--cycle;
\draw[poporange, dashed] (O)--(M) node[midway, right] {$a$ (apothème)};
\draw[poppurple, dashed] (O)--(A) node[midway, left] {$R$ (rayon)};
\fill (O) circle (1.5pt) node[below left] {$O$};
\fill (M) circle (1.5pt) node[below] {$M$};
\draw[popgreen] (0,-1.3) -- (0.2,-1.3) -- (0.2,-1.5); % angle droit
\end{tikzpicture}
\legende{Carré de centre $O$ : apothème $OM$ et rayon $OA$.}
\end{popfigure}
\end{enumerate}
\end{corrige}

\begin{corrige}[Exercice 4 — Calculs directs]
\begin{enumerate}
\item Polygone à $10$ côtés (décagone régulier) : angle au centre $= \dfrac{360^{\circ}}{10} = 36^{\circ}$.
\item Octogone régulier ($n=8$) :
\begin{align*}
\text{Somme des angles} &= (8-2)\times 180^{\circ} = 1080^{\circ},\\
\text{Angle intérieur} &= \dfrac{1080^{\circ}}{8} = 135^{\circ}.
\end{align*}
\item Carré de côté $c=6\,\text{cm}$ :
\begin{align*}
\text{Périmètre } \mathcal{P} &= 4c = 24\,\text{cm},\\
\text{Aire } \mathcal{A} &= c^2 = 36\,\text{cm}^2,\\
\text{Diagonale } d &= c\sqrt{2} = 6\sqrt{2}\,\text{cm} \quad (\text{valeur exacte}).
\end{align*}
\item Triangle équilatéral de côté $c=8\,\text{cm}$ :
\begin{align*}
\text{Hauteur } h &= \dfrac{c\sqrt{3}}{2} = 4\sqrt{3}\,\text{cm} \quad (\text{valeur exacte}),\\
\text{Aire } \mathcal{A} &= \dfrac{c\times h}{2} = \dfrac{8\times 4\sqrt{3}}{2} = 16\sqrt{3}\,\text{cm}^2 \quad (\text{valeur exacte}).
\end{align*}
Valeur approchée : $\sqrt{3}\approx 1,732$, donc $\mathcal{A}\approx 16\times 1,732 = 27,712\,\text{cm}^2$.
$1\,\text{cm}^2 = 100\,\text{mm}^2$, donc $\mathcal{A}\approx 2771,2\,\text{mm}^2 \approx \boxed{2771\,\text{mm}^2}$ (arrondi au $\text{mm}^2$).
\end{enumerate}
\end{corrige}

\begin{corrige}[Exercice 5 — Nombre de côtés]
Soit un polygone régulier d'angle intérieur $140^{\circ}$.
\begin{enumerate}
\item L'angle au centre est supplémentaire de l'angle intérieur (triangle isocèle sommet au centre) :
$\text{Angle au centre} = 180^{\circ} - 140^{\circ} = 40^{\circ}$.
\item Nombre de côtés $n = \dfrac{360^{\circ}}{\text{angle au centre}} = \dfrac{360^{\circ}}{40^{\circ}} = 9$.
C'est un \cle{nonagone} régulier (ou ennéagone).
\item Vérification : somme des angles intérieurs $= (n-2)\times 180^{\circ} = 7\times 180^{\circ} = 1260^{\circ}$.
Chaque angle intérieur $= \dfrac{1260^{\circ}}{9} = 140^{\circ}$. \checkmark
\end{enumerate}
\end{corrige}

\begin{corrige}[Exercice 6 — Hexagone inscrit]
Données : Hexagone régulier $ABCDEF$, centre $O$, rayon $R=8\,\text{cm}$.
\begin{enumerate}
\item Le triangle $OAB$ est isocèle en $O$ car $OA=OB=R$. L'angle au centre $\widehat{AOB}=\frac{360^{\circ}}{6}=60^{\circ}$. Un triangle isocèle avec un angle de $60^{\circ}$ au sommet est équilatéral. Donc $AB=OA=OB=8\,\text{cm}$.
\item Périmètre $\mathcal{P} = 6 \times AB = 6 \times 8 = 48\,\text{cm}$.
\item L'apothème $a$ est la hauteur du triangle équilatéral $OAB$ de côté $8\,\text{cm}$ :
$a = \dfrac{8\sqrt{3}}{2} = 4\sqrt{3}\,\text{cm}$ (valeur exacte).
\item Aire $\mathcal{A} = \dfrac{\mathcal{P}\times a}{2} = \dfrac{48 \times 4\sqrt{3}}{2} = 96\sqrt{3}\,\text{cm}^2$ (valeur exacte).
Approché : $\sqrt{3}\approx 1,732$, $\mathcal{A}\approx 96\times 1,732 = 166,272\,\text{cm}^2 \approx \boxed{166\,\text{cm}^2}$ (arrondi au $\text{cm}^2$).
\end{enumerate}
\end{corrige}

\begin{corrige}[Exercice 7 — Construction d'un octogone]
\begin{enumerate}
\item \textbf{Construction (règle et compas) :}
\begin{enumerate}
\item Tracer un cercle de centre $O$, rayon $5\,\text{cm}$.
\item Tracer un diamètre $[AA']$.
\item Construire la médiatrice de $[AA']$ (ou un diamètre perpendiculaire) : on obtient $[BB']$. Les points $A,A',B,B'$ sont 4 sommets.
\item Construire les bissectrices des angles $\widehat{AOB}$, $\widehat{BOA'}$, $\widehat{A'OB'}$, $\widehat{B'OA}$. Elles coupent le cercle en 4 nouveaux sommets $C,D,E,F$.
\item Relier les 8 sommets dans l'ordre : $A,C,B,D,A',E,B',F$.
\end{enumerate}
\begin{popfigure}
\begin{tikzpicture}[scale=0.8]
\draw[popblue, thick] (0,0) circle (3);
\foreach \ang/\label in {0/A, 45/C, 90/B, 135/D, 180/A', 225/E, 270/B', 315/F} {
  \coordinate (\label) at (\ang:3);
  \fill (\label) circle (2pt) node[shift=(\ang:0.4)] {\label};
}
\draw[poporange] (A)--(C)--(B)--(D)--(A')--(E)--(B')--(F)--cycle;
\draw[dashed, thin] (A)--(A') (B)--(B');
\draw[dashed, thin, poppurple] (C)--(E) (D)--(F);
\fill (0,0) circle (2pt) node[below left] {$O$};
\end{tikzpicture}
\legende{Octogone régulier inscrit : construction par bissectrices.}
\end{popfigure}
\item Angle au centre $= \dfrac{360^{\circ}}{8} = 45^{\circ}$.
\item Angle intérieur $= 180^{\circ} - 45^{\circ} = 135^{\circ}$ (ou $\frac{6\times 180^{\circ}}{8}=135^{\circ}$).
\end{enumerate}
\end{corrige}

\begin{corrige}[Exercice 8 — Carré et diagonale]
Données : Carré de diagonale $d=20\,\text{m}$.
\begin{enumerate}
\item Côté $c = \dfrac{d}{\sqrt{2}} = \dfrac{20}{\sqrt{2}} = 10\sqrt{2}\,\text{m}$ (valeur exacte).
\item Périmètre $\mathcal{P} = 4c = 40\sqrt{2}\,\text{m}$.
Aire $\mathcal{A} = c^2 = (10\sqrt{2})^2 = 200\,\text{m}^2$.
\item Longueur de clôture nécessaire $= \mathcal{P} - 4 = 40\sqrt{2} - 4\,\text{m}$.
Avec $\sqrt{2}\approx 1,414$ : $40\times 1,414 - 4 = 56,56 - 4 = 52,56\,\text{m}$.
Coût total $= 52,56 \times 2500 = 131\,400\,\text{F}$.
\end{enumerate}
\end{corrige}

\begin{corrige}[Exercice 9 — La table octogonale]
\begin{aretenir}[Barème indicatif (sur 20 pts)]
\begin{enumerate}
\item Q1 : 1 pt ; Q2 : 2 pts ; Q3 : 3 pts ; Q4 : 2 pts ; Q5 : 3 pts ; Q6 : 4 pts ; Q7 : 3 pts.
\item \textbf{Points de vigilance :} Conversion unités (m, cm, dm²) ; utilisation des valeurs approchées données ($\sin, \cos$) ; distinction aire disque / aire octogone ; arrondis demandés.
\end{enumerate}
\end{aretenir}

Données : Disque diamètre $1,20\,\text{m}$, octogone régulier inscrit. $\sin 22,5^{\circ}\approx 0,383$, $\cos 22,5^{\circ}\approx 0,924$.
\begin{enumerate}
\item Rayon $R = \dfrac{1,20}{2} = 0,60\,\text{m} = 60\,\text{cm}$.
\item Angle au centre $= \dfrac{360^{\circ}}{8} = 45^{\circ}$.
Angle intérieur $= 180^{\circ} - 45^{\circ} = 135^{\circ}$.
\item Côté $c = 2R\sin 22,5^{\circ} = 2\times 0,60 \times 0,383 = 0,4596\,\text{m} = 45,96\,\text{cm} \approx \boxed{46\,\text{cm}}$.
Périmètre $\mathcal{P} = 8c = 8 \times 0,4596 = 3,6768\,\text{m} \approx \boxed{368\,\text{cm}}$ (arrondi au cm).
\item Apothème $a = R\cos 22,5^{\circ} = 0,60 \times 0,924 = 0,5544\,\text{m} = \boxed{55,44\,\text{cm}}$.
\item Aire de la table $\mathcal{A}_{\text{table}} = \dfrac{\mathcal{P}\times a}{2} = \dfrac{3,6768 \times 0,5544}{2} = 1,0192\ldots\,\text{m}^2$.
En $\text{dm}^2$ : $1\,\text{m}^2 = 100\,\text{dm}^2$, donc $\mathcal{A}_{\text{table}} \approx 101,92\,\text{dm}^2 \approx \boxed{102\,\text{dm}^2}$.
\item Aire du disque $\mathcal{A}_{\text{disque}} = \pi R^2 = \pi \times 0,60^2 = 0,36\pi \approx 1,1310\,\text{m}^2 = 113,10\,\text{dm}^2$ (avec $\pi\approx 3,1416$).
Aire des chutes $\mathcal{A}_{\text{chutes}} = \mathcal{A}_{\text{disque}} - \mathcal{A}_{\text{table}} \approx 113,10 - 101,92 = 11,18\,\text{dm}^2 \approx \boxed{11\,\text{dm}^2}$.
\item Coût des chutes : Surface en $\text{m}^2 = 0,1118\,\text{m}^2$.
Coût $= 0,1118 \times 4500 = 503,1\,\text{F} \approx \boxed{503\,\text{F}}$ (arrondi au franc).
\end{enumerate}
\end{corrige}

\begin{corrige}[Exercice 10 — Le terrain de M. Ngoa]
\begin{aretenir}[Barème indicatif (sur 20 pts)]
\begin{enumerate}
\item Q1 : 2 pts (figure codée) ; Q2 : 2 pts ; Q3 : 2 pts ; Q4 : 4 pts ; Q5 : 5 pts ; Q6 : 5 pts.
\item \textbf{Points de vigilance :} Périmètre : ne pas compter $[CD]$ (intérieur) ; Aire : somme aire carré + aire triangle ; Unités (m, m²) ; Arrondis ; Calcul du nombre de sacs (arrondi au supérieur).
\end{enumerate}
\end{aretenir}

\begin{enumerate}
\item \textbf{Figure codée :}
\begin{popfigure}
\begin{tikzpicture}[scale=0.5]
\coordinate (A) at (0,0);
\coordinate (B) at (6,0);
\coordinate (C) at (6,6);
\coordinate (D) at (0,6);
\coordinate (E) at (3,6+3*1.732); % hauteur triangle eq 6 -> 3*sqrt(3) ~ 5.196
\draw[popblue, very thick] (A)--(B)--(C)--(E)--(D)--cycle;
\draw[popblue, very thick] (A)--(D); % côté gauche
\draw[popblue, very thick] (A)--(B); % bas
\draw[popblue, very thick] (B)--(C); % droite
\draw[dashed, thin] (C)--(D); % CD intérieur
\foreach \p/\pos in {A/below left, B/below right, C/right, D/left, E/above} {
  \fill (\p) circle (3pt) node[\pos] {\p};
}
% Codage
\draw[|-|] (0,-0.5) -- (6,-0.5) node[midway, below] {30 m};
\draw[|-|] (6.5,0) -- (6.5,6) node[midway, right] {30 m};
\draw[|-|] (0,6.5) -- (3,6.5) node[midway, above] {30 m};
\draw[|-|] (3,6.5) -- (6,6.5) node[midway, above] {30 m};
% traits d'égalité sur les côtés du triangle
\draw[poporange] (3.5, 6+1.5) -- (4.5, 6+1.5);
\draw[poporange] (3.5, 6+3.5) -- (4.5, 6+3.5);
\draw[poporange] (1.5, 6+2.5) -- (2.5, 6+2.5);
\draw[poporange] (1.5, 6+4.5) -- (2.5, 6+4.5);
\end{tikzpicture}
\legende{Terrain : carré $ABCD$ + triangle équilatéral $CDE$. $[CD]$ en pointillés (intérieur).}
\end{popfigure}
Les côtés $[AB], [BC], [CE], [ED], [DA]$ font $30\,\text{m}$. $[CD]$ est intérieur.
\item Périmètre $\mathcal{P} = 5 \times 30 = 150\,\text{m}$.
\item Hauteur du triangle équilatéral $CDE$ de côté $30\,\text{m}$ :
$h = \dfrac{30\sqrt{3}}{2} = 15\sqrt{3}\,\text{m}$ (valeur exacte).
\item Aire totale $\mathcal{A} = \mathcal{A}_{\text{carré}} + \mathcal{A}_{\text{triangle}}$.
$\mathcal{A}_{\text{carré}} = 30^2 = 900\,\text{m}^2$.
$\mathcal{A}_{\text{triangle}} = \dfrac{30 \times 15\sqrt{3}}{2} = 225\sqrt{3}\,\text{m}^2$.
$\mathcal{A} = 900 + 225\sqrt{3}\,\text{m}^2$ (valeur exacte).
Approché : $\sqrt{3}\approx 1,732 \Rightarrow 225\times 1,732 = 389,7$.
$\mathcal{A} \approx 900 + 389,7 = 1289,7\,\text{m}^2 \approx \boxed{1290\,\text{m}^2}$.
\item Longueur de grillage $= \mathcal{P} - \text{portail} = 150 - 4 = 146\,\text{m}$.
Dépense $= 146 \times 2500 = 365\,000\,\text{F}$.
\item Nombre de sacs $= \dfrac{\mathcal{A}}{50} \approx \dfrac{1290}{50} = 25,8 \Rightarrow \boxed{26\,\text{sacs}}$ (arrondi au supérieur).
Dépense semence $= 26 \times 6000 = 156\,000\,\text{F}$.
\end{enumerate}
\end{corrige}

\begin{corrige}[Exercice 11 — Somme des angles et pavage]
\begin{aretenir}[Barème indicatif (sur 20 pts)]
\begin{enumerate}
\item Partie A : Q1 (3 pts), Q2 (3 pts), Q3 (2 pts).
\item Partie B : Q1 (2 pts), Q2 (3 pts), Q3 (3 pts), Q4 (4 pts).
\item \textbf{Points de vigilance :} Justification rigoureuse du découpage en $n-2$ triangles (diagonales non sécantes à l'intérieur) ; condition de pavage : somme des angles autour d'un point $=360^{\circ}$ ; calcul exact des angles intérieurs.
\end{enumerate}
\end{aretenir}

\textbf{Partie A : Somme des angles intérieurs.}
\begin{enumerate}
\item On choisit un sommet $A_1$ d'un polygone convexe à $n$ sommets $A_1, A_2, \dots, A_n$. On trace les diagonales $[A_1A_3], [A_1A_4], \dots, [A_1A_{n-1}]$. Il y a $n-3$ diagonales. Elles partagent le polygone en $n-2$ triangles : $\triangle A_1A_2A_3$, $\triangle A_1A_3A_4$, $\dots$, $\triangle A_1A_{n-1}A_n$. Ces triangles n'ont pas de points intérieurs communs et leur réunion est le polygone.
\item La somme des angles de chaque triangle vaut $180^{\circ}$. La somme des angles intérieurs du polygone est la somme des angles de ces $n-2$ triangles. Donc $\mathcal{S}_n = (n-2)\times 180^{\circ}$.
\item Pour un polygone \emph{régulier}, les $n$ angles intérieurs sont égaux. Chacun vaut $\dfrac{\mathcal{S}_n}{n} = \dfrac{(n-2)\times 180^{\circ}}{n}$.
\end{enumerate}

\textbf{Partie B : Pavage du sol.}
\begin{enumerate}
\item Hexagone régulier ($n=6$) : angle intérieur $= \dfrac{(6-2)\times 180^{\circ}}{6} = \dfrac{720^{\circ}}{6} = 120^{\circ}$.
\item Autour d'un point, on peut placer $k$ hexagones. La somme des angles doit faire $360^{\circ}$ : $k \times 120^{\circ} = 360^{\circ} \Rightarrow k=3$. Trois hexagones réguliers se rejoignent exactement sans trou ni chevauchement. \textbf{Oui, c'est possible.}
\item Octogone régulier ($n=8$) : angle intérieur $= 135^{\circ}$.
On cherche un entier $k$ tel que $k \times 135^{\circ} = 360^{\circ}$.
$k = \dfrac{360}{135} = \dfrac{8}{3} \approx 2,66$. Ce n'est pas un entier. \textbf{Non, on ne peut pas paver le plan avec des octogones réguliers seuls.}
\item \textbf{Triangle équilatéral} ($n=3$) : angle $= 60^{\circ}$. $6 \times 60^{\circ} = 360^{\circ}$. \textbf{Oui.}
\textbf{Carré} ($n=4$) : angle $= 90^{\circ}$. $4 \times 90^{\circ} = 360^{\circ}$. \textbf{Oui.}
(Seuls ces trois polygones réguliers permettent un pavage monoédrique du plan).
\end{enumerate}
\end{corrige}

\newpage

\coursec{Fiche bilan}

\begin{popfigure}
\begin{tikzpicture}[
  mindmap,
  grow cyclic,
  every node/.style={concept, execute at begin node=\hskip0pt},
  concept color=popblue,
  level 1/.append style={level distance=4.5cm, sibling angle=360/6},
  level 2/.append style={level distance=3cm, sibling angle=360/4},
  branch/.style={concept color=popblue!70!popgreen, text=popdark},
  branch2/.style={concept color=poporange!80!poppink, text=popdark},
  branch3/.style={concept color=poppurple!70!popblue, text=popdark},
  branch4/.style={concept color=popgreen!60!popblue, text=popdark},
  branch5/.style={concept color=popgold!70!poporange, text=popdark},
  branch6/.style={concept color=poppink!70!poppurple, text=popdark},
  font=\small\bfseries,
  text width=2.8cm,
  align=center
]
\node[concept, concept color=popblue, text=white, font=\large\bfseries, text width=3.2cm] {Polygones\\ réguliers}
  child[branch] { node {Définition \& Propriétés} 
    child { node {Côtés égaux} }
    child { node {Angles égaux} }
    child { node {Inscriptible} }
    child { node {Centre $O$ commun} } }
  child[branch2] { node {Angles} 
    child { node {Angle au centre $\frac{360^\circ}{n}$} }
    child { node {Angle intérieur $\frac{(n-2)180^\circ}{n}$} }
    child { node {Angle extérieur $\frac{360^\circ}{n}$} }
    child { node {Somme int. $(n-2)180^\circ$} } }
  child[branch3] { node {Cas particuliers} 
    child { node {Triangle équilatéral ($n=3$)} }
    child { node {Carré ($n=4$)} }
    child { node {Hexagone régulier ($n=6$)} }
    child { node {Octogone régulier ($n=8$)} } }
  child[branch4] { node {Apothème $a$} 
    child { node {Hauteur triangle isocèle} }
    child { node {$a = R\cos\frac{180^\circ}{n}$} }
    child { node {$a = \frac{c}{2}\cot\frac{180^\circ}{n}$} }
    child { node {Rayon cercle inscrit} } }
  child[branch5] { node {Périmètre \& Aire} 
    child { node {$P = n \times c$} }
    child { node {$A = \frac{P \times a}{2}$} }
    child { node {$A = n \times \frac{c \times a}{2}$} }
    child { node {$A = \frac{n R^2}{2}\sin\frac{360^\circ}{n}$} } }
  child[branch6] { node {Applications} 
    child { node {Pavage (carrelage)} }
    child { node {Architecture, design} }
    child { node {Résolution problèmes} }
    child { node {Construction à la règle \& compas} } };
\end{tikzpicture}
\legende{Carte mentale : synthèse de la leçon sur les polygones réguliers (3ème).}
\end{popfigure}

\begin{bilan}

\courssub{Définitions et propriétés fondamentales}

\begin{itemize}
  \item \cle{Polygone régulier} : polygone convexe dont tous les \textbf{côtés ont même longueur} et tous les \textbf{angles intérieurs ont même mesure}.
  \item Il est \textbf{inscriptible} dans un cercle (centre $O$, rayon $R$) et \textbf{circconscrit} à un cercle (centre $O$, rayon $a$ = apothème).
  \item Les diagonales issues d'un même sommet divisent l'angle intérieur en angles égaux.
  \item Le centre $O$ est le centre de symétrie (si $n$ pair) et le centre de rotation d'angle $\frac{360^\circ}{n}$.
\end{itemize}

\courssub{Formules à connaître par cœur}

\begin{center}
\begin{tabularx}{\linewidth}{|>{\bfseries}l|X|X|}\hline
\textbf{Notion} & \textbf{Formule} & \textbf{Particuliers ($n$)} \\ \hline
Angle au centre & $\widehat{AOB} = \dfrac{360^\circ}{n}$ & $120^\circ$ ($n=3$), $90^\circ$ ($n=4$), $60^\circ$ ($n=6$), $45^\circ$ ($n=8$) \\ \hline
Angle intérieur & $\widehat{ABC} = \dfrac{(n-2)\times 180^\circ}{n}$ & $60^\circ$, $90^\circ$, $120^\circ$, $135^\circ$ \\ \hline
Angle extérieur & $\widehat{ext} = \dfrac{360^\circ}{n}$ & Supplémentaire de l'intérieur \\ \hline
Apothème $a$ & $a = R\cos\dfrac{180^\circ}{n} = \dfrac{c}{2}\cot\dfrac{180^\circ}{n}$ & Tri eq: $a=\frac{c\sqrt{3}}{6}$ ; Carré: $a=\frac{c}{2}$ ; Hex: $a=\frac{c\sqrt{3}}{2}$ \\ \hline
Périmètre $P$ & $P = n \times c$ & $3c$, $4c$, $6c$, $8c$ \\ \hline
Aire $\mathcal{A}$ & $\mathcal{A} = \dfrac{P \times a}{2} = n \times \dfrac{c \times a}{2}$ & Tri eq: $\frac{c^2\sqrt{3}}{4}$ ; Carré: $c^2$ ; Hex: $\frac{3c^2\sqrt{3}}{2}$ \\ \hline
\end{tabularx}
\end{center}

\courssub{Méthodes express}

\begin{enumerate}
  \item \textbf{Calculer un angle} : Identifier $n$, appliquer la formule (centre, intérieur, extérieur). Vérifier : intérieur + extérieur = $180^\circ$.
  \item \textbf{Calculer l'apothème} : Selon les données : $a = R\cos(180^\circ/n)$ ou $a = \frac{c}{2}\cot(180^\circ/n)$ ou Pythagore dans le triangle rectangle $(O, \text{milieu côté}, \text{sommet})$.
  \item \textbf{Calculer l'aire} : Décomposer en $n$ triangles isocèles de base $c$ et hauteur $a$ $\Rightarrow$ $\mathcal{A} = n \times \frac{c \times a}{2}$. Ou utiliser $\mathcal{A} = \frac{P \times a}{2}$.
  \item \textbf{Construction (règle \& compas)} : Tracer le cercle circonscrit $(O; R)$, reporter la corde de longueur $c = 2R\sin(180^\circ/n)$ (ou construire l'angle au centre $360^\circ/n$) $n$ fois.
\end{enumerate}

\end{bilan}

\begin{aretenir}[Checklist avant le BEPC]
\begin{itemize}
  \item $\square$ Je sais donner la définition complète d'un polygone régulier (côtés égaux, angles égaux, convexe).
  \item $\square$ Je connais par cœur les formules de l'angle au centre, de l'angle intérieur et de l'angle extérieur en fonction de $n$.
  \item $\square$ Je sais calculer les angles pour $n=3, 4, 6, 8$ sans calculette ($60^\circ, 90^\circ, 120^\circ, 135^\circ$ etc.).
  \item $\square$ Je sais identifier le centre $O$, le rayon $R$ (cercle circonscrit) et l'apothème $a$ (cercle inscrit) sur une figure.
  \item $\square$ Je sais appliquer les formules $a = R\cos(180^\circ/n)$ et $a = \frac{c}{2}\cot(180^\circ/n)$ (ou Pythagore) pour trouver l'apothème.
  \item $\square$ Je sais calculer le périmètre $P = n \times c$ et l'aire $\mathcal{A} = \frac{P \times a}{2}$.
  \item $\square$ Je connais les aires exactes du triangle équilatéral ($\frac{c^2\sqrt{3}}{4}$), du carré ($c^2$) et de l'hexagone régulier ($\frac{3c^2\sqrt{3}}{2}$).
  \item $\square$ Je sais décomposer un polygone régulier en $n$ triangles isocèles (ou rectangles) pour retrouver les formules.
  \item $\square$ Je sais résoudre un problème concret (pavage, terrain, objet) en modélisant par un polygone régulier.
  \item $\square$ Je sais justifier chaque étape de mon raisonnement avec le vocabulaire précis (médiatrice, bissectrice, hauteur, triangle rectangle, trigonométrie).
  \item $\square$ Je vérifie la cohérence de mes résultats (ordre de grandeur, unité de mesure, somme des angles).
  \item $\square$ Je sais construire un carré, un hexagone ou un octogone régulier à la règle et au compas.
\end{itemize}
\end{aretenir}

\findecours

\end{document}
